\documentclass[11pt]{article}
\usepackage[T1]{fontenc}
\usepackage{lmodern}
\usepackage[a4paper,margin=27mm]{geometry}
\usepackage{microtype}
\usepackage{amsmath}
\usepackage{amsthm}
\usepackage{amssymb}
\usepackage{longtable}
\usepackage{array}
\usepackage[hidelinks]{hyperref}
\setlength{\emergencystretch}{2em}

% A formal name, in brackets, breakable after each dot and underscore.
\makeatletter
\newcommand{\formal@break}[1]{\@tfor\formal@tok:=#1\do{%
  \formal@tok
  \expandafter\ifx\formal@tok.\allowbreak\else
  \expandafter\ifx\formal@tok\_\allowbreak\fi\fi}}
\DeclareRobustCommand{\formal}[1]{\textnormal{[\texttt{\small\formal@break{#1}}]}}
\makeatother
\newcommand{\R}{\mathbb{R}}
\newcommand{\N}{\mathbb{N}}
\newcommand{\E}{\mathbb{E}}
\newcommand{\Prob}{\mathbb{P}}
\newcommand{\Sym}{\mathcal{S}}
\newcommand{\Spsd}{\mathcal{S}_{+}}
\newcommand{\ip}[2]{\langle #1, #2\rangle}
\newcommand{\norm}[1]{\lVert #1\rVert}
\newcommand{\abs}[1]{\lvert #1\rvert}
\newcommand{\tr}{\operatorname{tr}}
\newcommand{\op}{\mathrm{op}}
\newcommand{\Fro}{F}
\newcommand{\sqd}{\sqrt{d}}
\newcommand{\Id}{I}
\newcommand{\cone}{\mathcal{K}}
\newcommand{\gauss}{\gamma}
\newcommand{\Dd}{D}
\newcommand{\one}{\mathbf{1}}
\newcommand{\Fit}{\mathrm{Fit}}
\newcommand{\Cond}{\mathrm{Cond}}
\newcommand{\Bal}{\mathrm{Bal}}
\newcommand{\sign}{\operatorname{sign}}
\newcommand{\diag}{\operatorname{diag}}
\newcommand{\rank}{\operatorname{rank}}
\newcommand{\supp}{\operatorname{supp}}

\theoremstyle{plain}
\newtheorem{theorem}{Theorem}[section]
\newtheorem{proposition}[theorem]{Proposition}
\newtheorem{lemma}[theorem]{Lemma}
\newtheorem{corollary}[theorem]{Corollary}
\theoremstyle{definition}
\newtheorem{definition}[theorem]{Definition}
\newtheorem{remark}[theorem]{Remark}
\newtheorem{convention}[theorem]{Convention}

\title{The ellipsoid fitting transition at $n = d^2/4$:\\ the proof as formalized}
\author{}
\date{}

\begin{document}
\maketitle

\begin{abstract}
This document presents, as a mathematical paper, the proof of the sharp
threshold for the ellipsoid fitting problem that has been formalized in
Rocq (the repository \texttt{EllipsoidFitting} with its library
\texttt{GaussianAnalysis}). It follows the route of the formal proof,
which departs from the informal proof of arXiv:2608.10184v1 in several
places: the spectral edge of the quadratic feature map is proved by
self-contained moment methods; Gordon's inequality is used in expectation
on finite nets; the fit of the satisfiable phase is produced from a kernel
element of a bounded-condition cone; the unsatisfiable phase works with a
constant spectral cutoff, nets of orientations, finite priors of
candidates, and a Lindeberg replacement with a replica expansion in place
of the free-entropy universality of Bandeira and Maillard. Each stated
result carries, in brackets, the name of the formal theorem and, when
there is one, the number of the corresponding result of the paper.
\end{abstract}

\tableofcontents

\section{The three theorems}\label{sec:main}

\subsection{The formal setting}

Throughout, $d \ge 1$ is the dimension and $n$ is the number of
observations; in the limit theorems $n = n(d)$ is an arbitrary sequence of
natural numbers. The only probability space is the following one.

\begin{definition}[Samples and their law]\label{def:sample}
A \emph{sample} of $n$ points of $\R^d$ is a tuple
$x = (x_0, \dots, x_{nd - 1})$ of $nd$ real numbers, carrying the product
Borel $\sigma$-algebra of the coordinates
\formal{PaperMainDefinitions.paper\_main\_sample}. Its $i$-th point
($0 \le i < n$) is the column vector
\[
x_i = (x_{id}, x_{id + 1}, \dots, x_{id + d - 1})^{\top} \in \R^d
\tag*{\formal{PaperMainDefinitions.paper\_main\_point}.}
\]
The law of the sample is the probability measure $\Prob_{n, d}$ on
$nd$-tuples defined by recursion on the number of coordinates: the law of
the empty tuple is the Dirac mass, and the law of $(m+1)$-tuples is the
image of the product of the standard normal law $N(0, 1)$ on $\R$ with the
law of $m$-tuples under $(t, u) \mapsto (t, u)$
\formal{GaussianProbabilityProductDefinitions.gaussian\_product\_law}. It
is the law of $nd$ independent standard normal coordinates, and we write
$\Prob_{n,d}(E) = $ \formal{PaperMainDefinitions.paper\_main\_probability}
for the probability of a measurable event $E$ of samples.
\end{definition}

Everything below is stated for the matrix space $\R^{d \times d}$ with the
following vocabulary, which is that of the formal statement
\formal{PaperMainDefinitions}. For $S \in \R^{d\times d}$ and
$v \in \R^d$, the quadratic form is $v^{\top} S v = \sum_{i, j} v_i S_{ij} v_j$.
$S$ is \emph{symmetric} when $S_{ij} = S_{ji}$ for all $i, j$; $S$ is
\emph{positive semidefinite}, written $S \succeq 0$, when it is symmetric
and $v^{\top} S v \ge 0$ for every $v$; $S$ is \emph{positive definite},
$S \succ 0$, when it is symmetric and $v^{\top} S v > 0$ for every
$v \ne 0$; $A \preceq B$ means $B - A \succeq 0$; $\tr S = \sum_i S_{ii}$.

\begin{definition}[The three events]\label{def:events}
For $n, d \in \N$ and real numbers $\lambda_-, \lambda_+$:
\begin{align*}
\Fit_{n,d} &= \{x : \exists S \succeq 0,\ x_i^{\top} S x_i = d \text{ for all } i < n\}
  \tag*{\formal{PaperMainDefinitions.paper\_main\_fit\_event}}\\
\Cond_{n,d}(\lambda_-, \lambda_+) &= \{x : \exists S \succeq 0,\ x_i^{\top} S x_i = d \ (i < n),\ \tr S = d,\ \lambda_- \Id \preceq S \preceq \lambda_+ \Id\}
  \tag*{\formal{PaperMainDefinitions.paper\_main\_conditioned\_event}}\\
\Bal_{n,d} &= \Bigl\{x : \exists y \in \R^n,\ \sum_{i<n} y_i = 0,\ \sum_{i<n} y_i x_i x_i^{\top} \succ 0\Bigr\}
  \tag*{\formal{PaperMainDefinitions.paper\_main\_balanced\_pd\_event}}
\end{align*}
\end{definition}

The asymptotic hypotheses and conclusions are written out as follows
\formal{PaperMainDefinitions}: for a real sequence $u = (u_d)$ and a real
$c$,
\begin{itemize}
\item $u_d \le c$ \emph{eventually} means: there is $D$ with $u_d \le c$ for all $d \ge D$;
\item $\limsup_d u_d < c$ means: there is $c' < c$ with $u_d \le c'$ eventually;
\item $\limsup_d u_d \le c$ means: for every $\varepsilon > 0$, $u_d \le c + \varepsilon$ eventually;
\item $\liminf_d u_d > c$ means: there is $c' > c$ with $u_d \ge c'$ eventually;
\item for a sequence $p = (p_d)$ of extended reals and a real $\ell$,
  $p_d \to \ell$ means: for every $\varepsilon > 0$ there is $D$ with
  $\abs{p_d - \ell} < \varepsilon$ for all $d \ge D$ (an infinite $p_d$ never
  satisfies this).
\end{itemize}
The \emph{density} of a sequence $n = n(d)$ is $n(d)/d^2$
\formal{PaperMainDefinitions.paper\_main\_density}.

\subsection{Statements}

The formal main theorem \formal{MainTheorem.ellipsoid\_fitting\_main} is a
proof of the conjunction of the following three statements
\formal{MainStatement.main\_statement}; the formal proof of the
conjunction is \formal{PaperMainResults.paper\_main\_results}. Each
statement asserts the measurability of the events it measures as part of
its conclusion.

\begin{theorem}[Sharp threshold \formal{PaperMainResults.paper\_main\_results}; paper Theorem~1.2]\label{thm:main}
For every $n, d$ the event $\Fit_{n,d}$ is measurable, and for every
sequence $n = n(d)$:
\begin{enumerate}
\item[(a)] if $\limsup_d n(d)/d^2 < 1/4$, then $\Prob_{n(d), d}(\Fit_{n(d), d}) \to 1$;
\item[(b)] if $\liminf_d n(d)/d^2 > 1/4$, then $\Prob_{n(d), d}(\Fit_{n(d), d}) \to 0$.
\end{enumerate}
\end{theorem}

\begin{theorem}[Well-conditioned fits \formal{PaperMainConditionedLimit.paper\_main\_conditioned\_limit}; paper Theorem~1.3]\label{thm:cond}
For every real $\alpha^* < 1/4$ there are reals $0 < \lambda_- \le \lambda_+$
such that $\Cond_{n,d}(\lambda_-, \lambda_+)$ is measurable for all $n, d$,
and for every sequence $n = n(d)$ with $\limsup_d n(d)/d^2 \le \alpha^*$,
\[
\Prob_{n(d), d}\bigl(\Cond_{n(d), d}(\lambda_-, \lambda_+)\bigr) \to 1 .
\]
\end{theorem}

\begin{theorem}[Balanced positive definite combinations \formal{PaperMainFinalAssembly.paper\_main\_results\_from\_conditioned}; paper Corollary~1.4]\label{thm:bal}
For every $n, d$ the event $\Bal_{n,d}$ is measurable, and for every
sequence $n = n(d)$: if $\limsup_d n(d)/d^2 < 1/4$ then
$\Prob_{n(d), d}(\Bal_{n(d), d}) \to 0$, and if $\liminf_d n(d)/d^2 > 1/4$
then $\Prob_{n(d), d}(\Bal_{n(d), d}) \to 1$.
\end{theorem}

The constants $\lambda_\pm$ of Theorem~\ref{thm:cond} have the form
$\lambda_- = (2\kappa + 1)^{-1}$, $\lambda_+ = 2\kappa + 1$ for a
$\kappa = \kappa(\alpha^*) \ge 1$ produced by the proof; they are not
computed.

\subsection{Plan}

Section~\ref{sec:tools} presents the general tools of the library:
the standard Gaussian measure on $\R^m$, Gaussian integration by parts,
interpolation and the comparison inequalities of Slepian, Sudakov and
Fernique, and Gordon, escape through a mesh, concentration of Lipschitz
functions, the spectral theorem for real symmetric matrices, nets, and the
moment inequalities. Section~\ref{sec:prelim} develops the objects of the
problem: the quadratic features of the sample and their Gram matrix, the
bounded-condition cones and their Gaussian widths, the elementary estimates
on the sample, the spectral edges of the feature map, the homogeneous
formulation of the problem and the alternative behind
Theorem~\ref{thm:bal}, the measurability of the events, and the reduction
to fixed densities. Section~\ref{sec:sat} proves the satisfiable phase,
Section~\ref{sec:unsat} the unsatisfiable one, and Section~\ref{sec:assembly}
assembles the three theorems.

Where the formal proof follows another route than the paper, this is said
in one sentence at the place in question, and the formal argument is given.

Appendix~\ref{sec:correspondence} lists the 39 numbered results of the
paper with their counterparts in this document and their relation to the
formal statements, following the reconciled comparison
\texttt{doc/paper\_correspondence.md} of the repository. By statement, 18 of them are the same as in the paper, 15 are weaker,
6 are stronger and none absent; by argument, 26 follow the paper's route and
13 are proved by another argument. The appendix also lists the results
the paper cites without proof and what the formal proof does instead,
and the changes of route.

\section{Notation and conventions}\label{sec:notation}

\begin{convention}[Vectors, matrices, norms]
Vectors are columns. $\ip{u}{v} = u^{\top} v$ and $\norm{u}_2 = \ip{u}{u}^{1/2}$ on $\R^m$.
$\Sym^d$ is the space of real symmetric $d\times d$ matrices, of dimension
$\Dd_d = d(d+1)/2$; $\Spsd^d \subseteq \Sym^d$ is the cone of positive
semidefinite matrices. On $\R^{d \times d}$, $\ip{A}{B} = \sum_{i,j} A_{ij} B_{ij} = \tr(A^{\top} B)$
is the Frobenius pairing \formal{Geometry.Psd.paper\_main\_pairing} and
$\norm{A}_\Fro = \ip{A}{A}^{1/2}$ the Frobenius norm. For a symmetric
$A$, $\norm{A}_\op = \sup_{\norm{v}_2 = 1} \abs{v^{\top} A v}$; the formal
proof mostly uses the bilinear form of this bound,
$\abs{u^{\top} A v} \le b \norm{u}_2 \norm{v}_2$ for all $u, v$, written
$\norm{A}_\op \le b$ \formal{Cone.ConditionOperator.paper\_main\_cone\_opnorm\_le}.
The unit Frobenius sphere of $\Sym^d$ is $\mathbb{S}_\Fro$, and
$\mathbb{S}^{m-1}$ is the Euclidean unit sphere of $\R^m$.
\end{convention}

\begin{convention}[Gaussian objects]
$\gauss_m$ is the standard Gaussian measure on $\R^m$ (Section~\ref{sec:gauss}).
$g, h, \ldots$ denote standard Gaussian vectors, and $Z$ a standard
Gaussian element of $\Sym^d$, that is $Z = \sum_{k < \Dd_d} z_k E_k$ for an
orthonormal basis $(E_k)$ of $\Sym^d$ and $z \sim \gauss_{\Dd_d}$
\formal{Gaussian.Symmetric.PublicLaw.paper\_main\_public\_symmetric\_gaussian}.
The GOE rows of the paper are $G_i = \sqrt{2/d}\, Z_i$ with independent
copies $Z_i$ of $Z$; in the formal proof they are the \emph{Gaussian
feature rows} \formal{Gaussian.Symmetric.PublicFeature.paper\_main\_public\_gaussian\_row},
whose pairing with a symmetric $T$ is the linear form
$\sqrt{2/d}\, \sum_k z_k \ip{E_k}{T}$ of the coordinates.
\end{convention}

\begin{convention}[Expectations]
Two expectations occur. For a measurable $F$ on tuples,
$\E_{n,d} F = \int F\, d\Prob_{n,d}$ is the Lebesgue integral of
MathComp-Analysis. For a function $F$ of a list of $m$ real numbers, the
\emph{iterated expectation} $\E^{(m)} F$ integrates the coordinates one at
a time against the Gaussian kernel $e^{-t^2/2}/\sqrt{2\pi}$
\formal{Gaussian.Iterated.Definitions.gaussian\_iterated\_expectation},
each level being a Lebesgue integral on the line. For $F$ continuous with
polynomial growth the two agree, and both equal the integral of
$x \mapsto F(\text{coordinates of } x)$ against $\gauss_m$
(Proposition~\ref{prop:bridge}). We write $\E$ for either when no
confusion is possible.
\end{convention}

\begin{convention}[Sample sizes]
A statement "for all large $d$" means: there is $D$ such that the statement
holds for every $d \ge D$. The thresholds $D$ of the formal proof are
existential and, in a few places, explicit but enormous
(for instance $2^{9600}$ in the walk-count estimates of
Section~\ref{sec:edge}); they play no role in the limit theorems.
\end{convention}

\section{General tools}\label{sec:tools}

The results of this section are stated for an arbitrary real closed field
$\R$ of MathComp-Analysis and are the content of the library
\texttt{GaussianAnalysis}. The main proof instantiates them at the reals of
the standard library through a bridge (Proposition~\ref{prop:bridge}).

\subsection{The standard Gaussian measure}\label{sec:gauss}

\begin{definition}[\formal{GaussianAnalysis.gaussian.gauss.gauss}]
For $m \in \N$, $\gauss_m$ is the probability measure on the row vectors
$\R^{1\times m}$, with the Borel $\sigma$-algebra of the product topology,
defined by recursion on $m$: $\gauss_0$ is the Dirac mass at the empty
vector, and $\gauss_{m+1}$ is the image of the product $N(0,1) \otimes \gauss_m$
under $(t, x) \mapsto (t, x)$.
\end{definition}

\begin{proposition}[Boxes and uniqueness]\label{prop:gaussE}
For measurable $A_1, \dots, A_m \subseteq \R$,
$\gauss_m(\prod_i A_i) = \prod_i N(0,1)(A_i)$
\formal{GaussianAnalysis.gaussian.gauss.gauss\_box}. A probability
measure on $\R^{1\times m}$ that gives every box this value is $\gauss_m$
\formal{GaussianAnalysis.gaussian.gauss.gauss\_unique}. Consequently $\gauss_m$
is invariant under permutations of the coordinates
\formal{GaussianAnalysis.gaussian.gauss.integral\_gauss\_col\_perm} and
under $x \mapsto -x$ \formal{GaussianAnalysis.comparison.escape.integral\_gauss\_opp},
the law of a block of $k$ coordinates is $\gauss_k$, and for
$f$ integrable, $\int f\, d\gauss_{m+1} = \int \int f(t, x)\, d\gauss_m(x)\, dN(0,1)(t)$
(Fubini for the first coordinate).
\end{proposition}

\begin{proof}
The box formula is the recursion and the product measure formula. Two
probability measures that agree on the boxes agree on the $\pi$-system of
boxes, which generates the Borel sets of the product, hence on all Borel
sets. The invariances follow since a permutation, or the sign change, maps
boxes to boxes of the same measure.
\end{proof}

\begin{proposition}[Linear images \formal{GaussianAnalysis.comparison.gauss\_law.gauss\_mulmx\_law}]\label{prop:mulmx}
Let $A \in \R^{n\times k}$ and $B \in \R^{m \times k}$ with
$A^{\top} A = B^{\top} B$. Then for every measurable $X \subseteq \R^{1\times k}$,
\[
\gauss_n\{x : xA \in X\} = \gauss_m\{y : yB \in X\}.
\]
In particular, if the columns of $A$ are orthonormal, $x A \sim \gauss_k$
when $x \sim \gauss_n$ \formal{GaussianAnalysis.comparison.gauss\_law.gauss\_orthogonal\_law};
and $\gauss_n$ is invariant under every orthogonal matrix.
\end{proposition}

\begin{proof}
For a test $F$ of class $C^2$ whose derivatives grow at most polynomially,
Kahane's comparison (Theorem~\ref{thm:kahane}) in both directions gives
$\int F(xA)\, d\gauss_n = \int F(yB)\, d\gauss_m$. A bounded Lipschitz $G$
is the uniform limit of its double Gaussian smoothings
$G_{t} = \gauss\text{-}\mathrm{smooth}_t(\gauss\text{-}\mathrm{smooth}_t G)$,
which are of class $C^2$ with bounded derivatives, because integration by
parts moves one derivative inside the second smoothing
\formal{GaussianAnalysis.comparison.gauss\_law.gauss\_smooth\_partial\_stein};
hence the equality holds for $G$. Finally a finite measure on
$\R^{1\times k}$ is determined by the integrals of the Lipschitz functions
with values in $[0,1]$ \formal{GaussianAnalysis.comparison.gauss\_law.measure\_lipschitz\_unique}:
the sets $\{\Phi > 0\}$, for $\Phi \ge 0$ Lipschitz, are closed under
intersection and generate the Borel sets, and their indicators are
increasing limits of $\min(1, i\Phi)$.
\end{proof}

\begin{proposition}[The bridge to the sample law \formal{Gaussian.Law.LibraryBridge.gauss\_gaussian\_product\_law}]\label{prop:bridge}
At the reals of the standard library, let $\rho$ send an $m$-tuple to the
row vector of its coordinates. Then $\gauss_m$ is the image of
$\Prob$-law of $m$-tuples under $\rho$: for every measurable $X \subseteq \R^{1\times m}$,
$\gauss_m(X) = \Prob_{1, m}(\rho^{-1} X)$. Moreover, for $F$ a continuous
function of lists that is integrable at every level of the iterated
expectation, in particular when $F$ has polynomial growth,
\[
\E^{(m)} F = \int F(\rho(x))\, d\gauss_m(x)
\tag*{\formal{Gaussian.Law.LibraryBridge.gaussian\_iterated\_expectation\_gauss}.}
\]
\end{proposition}

\begin{proof}
Both measures give a box the product of the standard normal measures of
its sides (Proposition~\ref{prop:gaussE} and the recursion defining
$\Prob_{1,m}$), so they agree by uniqueness. The iterated expectation is a
Lebesgue integral at each level, and Fubini for the first coordinate,
applied $m$ times, identifies it with the integral against $\gauss_m$.
\end{proof}

All the probabilistic statements below about the sample are proved either
directly on the tuples, or on row vectors and transported through
Proposition~\ref{prop:bridge}. The bridge also carries over continuity,
polynomial growth and partial derivatives of a function of lists to the
corresponding function of row vectors \formal{Gaussian.Law.LibraryBridge.list\_of\_row\_derive}.

\subsection{Integration by parts}

\begin{theorem}[One-dimensional Stein identity \formal{GaussianAnalysis.gaussian.gauss\_ibp.normal\_ibp}]\label{thm:normal_ibp}
Let $f : \R \to \R$ be differentiable with continuous derivative $f'$, and
suppose $\abs{f(t)}, \abs{f'(t)} \le C(1 + \abs{t}^k)$. Then
$\int t f(t)\, dN(0,1)(t) = \int f'(t)\, dN(0,1)(t)$.
\end{theorem}

\begin{proof}
$t e^{-t^2/2}$ is the derivative of $-e^{-t^2/2}$; integrate by parts on
$[-M, M]$ and let $M \to \infty$: the boundary terms $f(\pm M) e^{-M^2/2}$
vanish by the growth bound, and both integrands are integrable on the line
by the growth bound and the Gaussian tail.
\end{proof}

\begin{theorem}[Gaussian integration by parts \formal{GaussianAnalysis.gaussian.gauss\_ibp.gauss\_ibp}, \formal{GaussianAnalysis.gaussian.gauss\_ibp.gauss\_ibp2}; paper Lemma~3.10, second-order identity]\label{thm:ibp}
Let $F : \R^{1\times m} \to \R$ be continuous with partial derivative
$\partial_i F$ along the $i$-th coordinate, continuous, both of
polynomial growth $C(1 + \norm{x}^k)$. Then
\[
\int x_i F(x)\, d\gauss_m(x) = \int \partial_i F(x)\, d\gauss_m(x) .
\]
If moreover $\partial_i F$ has a partial derivative $\partial_i^2 F$
along the same coordinate, continuous and of polynomial growth, then
\[
\int (x_i^2 - 1) F(x)\, d\gauss_m(x) = \int \partial_i^2 F(x)\, d\gauss_m(x) .
\]
\end{theorem}

\begin{proof}
Move the $i$-th coordinate to the first place (Proposition~\ref{prop:gaussE})
and apply Fubini: for fixed values of the other coordinates, the inner
integral is the one-dimensional identity of Theorem~\ref{thm:normal_ibp}.
The second identity is the first applied to $x_i F$, whose derivative is
$F + x_i \partial_i F$, followed by the first applied to $\partial_i F$.
Measurability of a partial derivative of a continuous function comes from
its expression as a limit of difference quotients
\formal{GaussianAnalysis.gaussian.gauss\_ibp.measurable\_derive\_rV}.
\end{proof}

At the reals of the standard library and in the language of iterated
expectations, these are the \emph{coordinate Stein identities}
\formal{Gaussian.Law.LibraryBridge.gaussian\_iterated\_coordinate\_stein\_gauss},
\formal{Gaussian.Law.LibraryBridge.gaussian\_iterated\_coordinate\_second\_stein\_gauss},
which the chaos computations of Sections~\ref{sec:sat} and~\ref{sec:unsat}
use: for $F$ continuous of polynomial growth with continuous derivatives of
polynomial growth,
\[
\E[(x_a x_b - \delta_{ab}) F(x)] = \E[\partial_a \partial_b F(x)] .
\]

\subsection{Interpolation and comparison}\label{sec:comparison}

\begin{theorem}[Gaussian interpolation \formal{GaussianAnalysis.gaussian.gauss\_interpolation.gauss\_interpolation}]\label{thm:interp}
Let $F : \R^{1\times k} \to \R$ be of class $C^2$ with partial derivatives
$\partial_j F$ and $\partial_j \partial_l F$ continuous and of polynomial
growth, and let $M : (u, v) \to \R^{n\times k}$ be differentiable with
derivative $M'$, both bounded on $(u, v)$. Then $t \mapsto \int F(xM(t))\, d\gauss_n(x)$
is differentiable on $(u, v)$ with derivative
\[
\sum_{j, l < k} \bigl(M'(t)^{\top} M(t)\bigr)_{jl} \int \partial_j \partial_l F(xM(t))\, d\gauss_n(x) .
\]
\end{theorem}

\begin{proof}
Differentiate under the integral (dominated convergence with the
polynomial growth and the bounds on $M, M'$): the derivative of
$F(xM(t))$ is $\sum_j \partial_j F(xM(t)) (xM'(t))_j = \sum_j \partial_j F(xM(t)) \sum_i x_i M'(t)_{ij}$.
Integration by parts (Theorem~\ref{thm:ibp}) in $x_i$ turns
$x_i \partial_j F(xM(t))$ into $\sum_l \partial_l \partial_j F(xM(t)) M(t)_{il}$,
and $\sum_i M'(t)_{ij} M(t)_{il} = (M'(t)^{\top} M(t))_{jl}$.
\end{proof}

\begin{theorem}[Kahane's comparison \formal{GaussianAnalysis.gaussian.gauss\_interpolation.gauss\_comparison}]\label{thm:kahane}
Let $F$ be as in Theorem~\ref{thm:interp} and let $A \in \R^{n \times k}$,
$B \in \R^{m \times k}$ satisfy
$\sum_{j,l} (A^{\top} A - B^{\top} B)_{jl}\, \partial_j \partial_l F(z) \le 0$
for every $z$. Then $\int F(xA)\, d\gauss_n \le \int F(yB)\, d\gauss_m$.
\end{theorem}

\begin{proof}
Put the two vectors side by side: for $w = (x, y) \sim \gauss_{n+m}$ and
$M(t) = \bigl(\sqrt{t}\, A;\ \sqrt{1-t}\, B\bigr)$ (stacked), $wM(t) = \sqrt t\, xA + \sqrt{1-t}\, yB$
has covariance $tA^{\top}A + (1-t)B^{\top}B$. Theorem~\ref{thm:interp}
gives $\frac{d}{dt} \int F(wM(t)) = \frac12 \sum_{j,l} (A^{\top}A - B^{\top}B)_{jl} \int \partial_j\partial_l F(wM(t)) \le 0$
on $(0,1)$, and the endpoints are the two sides, by continuity and the
polynomial growth.
\end{proof}

The next three inequalities concern finite centered Gaussian families
written as $x \mapsto xA$ for $x \sim \gauss_n$ and $A \in \R^{n \times p}$,
whose covariance is $A^{\top} A$. The \emph{increment variance} of two
indices is $\sigma_A(j, l) = \E (xA)_j - (xA)_l)^2 = (A^{\top}A)_{jj} + (A^{\top}A)_{ll} - 2(A^{\top}A)_{jl}$.

\begin{theorem}[Sudakov--Fernique \formal{GaussianAnalysis.comparison.slepian.sudakov\_fernique}; Slepian \formal{GaussianAnalysis.comparison.slepian.slepian}]\label{thm:sf}
If $\sigma_A(j, l) \le \sigma_B(j, l)$ for all $j, l < p$, then
$\E \max_j (xA)_j \le \E \max_j (yB)_j$ for $x \sim \gauss_n$, $y \sim \gauss_m$.
In particular, if the variances agree and $(B^{\top} B)_{jl} \le (A^{\top} A)_{jl}$
for all $j, l$, the same conclusion holds.
\end{theorem}

\begin{proof}
For $\beta > 0$ let $\ell_\beta(z) = \beta^{-1} \ln \sum_j e^{\beta z_j}$, the
soft maximum, which satisfies $\max_j z_j \le \ell_\beta(z) \le \max_j z_j + \beta^{-1}\ln p$.
Its Hessian is $\beta(\diag(\pi) - \pi \pi^{\top})$ where
$\pi_j = e^{\beta z_j}/\sum_l e^{\beta z_l}$, so for every symmetric
$\Sigma$, $\sum_{jl} \Sigma_{jl} \partial_j\partial_l \ell_\beta = \beta \sum_{jl} \pi_j \pi_l (\Sigma_{jj} - \Sigma_{jl}) = \frac{\beta}{2} \sum_{jl} \pi_j\pi_l (\Sigma_{jj} + \Sigma_{ll} - 2\Sigma_{jl})$.
With $\Sigma = A^{\top}A - B^{\top}B$ the last factor is
$\sigma_A(j,l) - \sigma_B(j,l) \le 0$, and Theorem~\ref{thm:kahane}
(with the growth of $\ell_\beta$ and its derivatives) gives
$\E \ell_\beta(xA) \le \E \ell_\beta(yB)$. Let $\beta \to \infty$.
\end{proof}

\begin{theorem}[Gordon's min--max comparison \formal{GaussianAnalysis.comparison.gordon.gordon}; paper Lemma~2.3, expectation form]\label{thm:gordon}
Let $p, q \ge 1$ and let $A \in \R^{n \times pq}$, $B \in \R^{m\times pq}$
index the entries by pairs $(u, t)$, $u < p$, $t < q$. Suppose
\[
\sigma_A((u,t),(u,s)) \le \sigma_B((u,t),(u,s)) \quad\text{and}\quad
\sigma_B((u,t),(v,s)) \le \sigma_A((u,t),(v,s)) \text{ for } u \ne v .
\]
Then $\E \min_u \max_t (xA)_{(u,t)} \le \E \min_u \max_t (yB)_{(u,t)}$
for $x \sim \gauss_n$, $y \sim \gauss_m$. With equal variances this reads:
larger covariances inside the rows and smaller covariances across rows
give a smaller expected min--max \formal{GaussianAnalysis.comparison.gordon.gordon\_cov}.
\end{theorem}

\begin{proof}
Replace $\min_u \max_t$ by a smooth min--max $s_\beta$
\formal{GaussianAnalysis.comparison.gordon.smm}: the soft minimum over $u$
of the soft maxima over $t$, which is within $2\beta^{-1}\ln(pq)$ of
$\min_u\max_t z_{(u,t)}$. Its Hessian is a combination of the Gibbs weights
of the two levels, and the sum $\sum_{jl} \Sigma_{jl}\partial_j\partial_l s_\beta$
with $\Sigma = A^{\top}A - B^{\top}B$ splits into a sum over pairs inside a
row, with coefficient $-(\sigma_A - \sigma_B)$ times a nonnegative weight,
and a sum over pairs in different rows, with coefficient $+(\sigma_A - \sigma_B)$
times a nonnegative weight. The hypotheses make both sums nonpositive, so
Theorem~\ref{thm:kahane} gives $\E s_\beta(xA) \le \E s_\beta(yB)$, and
$\beta \to \infty$ concludes. (The formal proof carries out this computation
on the explicit second derivatives of $s_\beta$.)
\end{proof}

The instance that the main proof uses is the comparison of a bilinear
Gaussian process with its additive surrogate.

\begin{corollary}[Bilinear processes \formal{GaussianAnalysis.comparison.escape.auxT\_bilT\_minmax}; paper Lemma~2.3, bilinear case]\label{cor:bil}
Let $z_0, \dots, z_{p-1}$ be unit vectors of $\R^N$ and $v_0, \dots, v_{q-1}$ unit
vectors of $\R^M$, let $\Gamma \in \R^{N \times M}$ have independent
$N(0,1)$ entries and let $g \sim \gauss_N$, $h \sim \gauss_M$ be
independent. Then
\[
\E \min_u \max_t \bigl(\ip{g}{z_u} + \ip{h}{v_t}\bigr) \le \E \min_u \max_t z_u^{\top} \Gamma v_t .
\]
\end{corollary}

\begin{proof}
Both processes have unit variances, and
$\E(z_u^{\top}\Gamma v_t)(z_{u'}^{\top}\Gamma v_s) = \ip{z_u}{z_{u'}} \ip{v_t}{v_s}$
while the surrogate has covariance $\ip{z_u}{z_{u'}} + \ip{v_t}{v_s}$. The
difference of increment variances is $2(1 - \ip{z_u}{z_{u'}})(1 - \ip{v_t}{v_s}) \ge 0$,
which vanishes for $u = u'$: Theorem~\ref{thm:gordon} applies with the
surrogate as $A$ and the bilinear process as $B$.
\end{proof}

\begin{theorem}[Escape through a mesh \formal{GaussianAnalysis.comparison.escape.gordon\_escape}, \formal{GaussianAnalysis.comparison.escape.gordon\_escape\_tail}; paper Lemma~2.4]\label{thm:escape}
Let $T$ be a nonempty set of unit vectors of $\R^M$, let
$w(T) = \E \sup_{v \in T} \ip{h}{v}$ ($h \sim \gauss_M$) be its Gaussian
width, and let $\Gamma \in \R^{N\times M}$, $N \ge 1$, have independent
$N(0,1)$ entries. Then
\[
\E \inf_{v \in T} \norm{\Gamma v}_2 \ge \E\norm{g}_2 - w(T), \qquad g \sim \gauss_N,
\]
and for $u \ge 0$,
$\Prob\{\inf_{v\in T} \norm{\Gamma v}_2 \le \E\norm{g}_2 - w(T) - u\} \le e^{-u^2/2}$.
\end{theorem}

\begin{proof}
For finite families $(v_t) \subseteq T$ and $(z_u)$ a net of the unit
sphere of $\R^N$, Corollary~\ref{cor:bil} with the roles of the two factors
exchanged gives
$\E \min_t \max_u z_u^{\top}\Gamma v_t \ge \E \max_u \ip{g}{z_u} - \E \max_t \ip{h}{v_t}$.
As the net of the sphere becomes finer, $\max_u z_u^{\top} \Gamma v_t \to \norm{\Gamma v_t}_2$
and $\max_u \ip{g}{z_u} \to \norm{g}_2$ (a $\delta$-net of the sphere
captures $(1-\delta)$ of the norm); as the family $(v_t)$ becomes a finer
net of $T$, the minimum tends to $\inf_T$ and $\E\max_t \ip{h}{v_t} \le w(T)$,
since $\norm{\Gamma v - \Gamma v'}_2 \le \norm{\Gamma}_\Fro \norm{v - v'}_2$.
The tail follows from Theorem~\ref{thm:conc}, the function
$\Gamma \mapsto \inf_T \norm{\Gamma v}_2$ being $1$-Lipschitz for the
Frobenius norm.
\end{proof}

\begin{theorem}[Norms of Gaussian matrices \formal{GaussianAnalysis.comparison.escape.gauss\_opnorm\_mean}, \formal{GaussianAnalysis.comparison.escape.gauss\_opnorm\_tail}; paper Lemma~2.5]\label{thm:opnorm}
For $\Gamma$ as above, $\E\norm{\Gamma}_\op \le \E\norm{g}_2 + \E\norm{h}_2 \le \sqrt N + \sqrt M$,
and $\Prob\{\norm{\Gamma}_\op \ge \E\norm{g}_2 + \E\norm{h}_2 + u\} \le e^{-u^2/2}$.
\end{theorem}

\begin{proof}
Theorem~\ref{thm:sf} compares $\max_{u,t} z_u^{\top}\Gamma v_t$ with
$\max_{u,t}(\ip{g}{z_u} + \ip{h}{v_t}) = \max_u\ip{g}{z_u} + \max_t \ip{h}{v_t}$
on nets of the two spheres, and the net argument of the previous proof
concludes; the tail is again Theorem~\ref{thm:conc}.
\end{proof}

The main proof does not use Theorem~\ref{thm:opnorm}: it bounds operator
norms of Gaussian matrices through nets and a union bound, with worse
constants (Lemma~\ref{lem:opnet}).

\subsection{Concentration}

\begin{theorem}[Gaussian concentration \formal{GaussianAnalysis.comparison.gauss\_concentration.gauss\_concentration}]\label{thm:conc}
Let $F : \R^{1\times m} \to \R$ satisfy $\abs{F(x) - F(y)} \le L \norm{x - y}_2$.
Then for every real $t$, $\E e^{t(F - \E F)} \le e^{t^2 L^2/2}$ under
$\gauss_m$. Consequently, for $u \ge 0$, each of $\Prob\{F \ge \E F + u\}$
and $\Prob\{F \le \E F - u\}$ is at most $e^{-u^2/(2L^2)}$, and
$\operatorname{Var} F \le L^2$ \formal{GaussianAnalysis.comparison.gauss\_concentration.gauss\_poincare\_lipschitz}.
\end{theorem}

\begin{proof}
First let $F$ be bounded, of class $C^1$, with $\norm{\nabla F}_2 \le L$,
and $t \ge 0$. Let $y$ be an independent copy of $x$ and
$\varphi(\theta) = \E F(x) e^{tF(\cos\theta\, x + \sin\theta\, y)}$.
After rotating $(x, y)$ by $\theta$, integration by parts in $y$
(Theorem~\ref{thm:ibp} on the pair) shows $\varphi'(\theta) \ge -\sin\theta\, tL^2\, \E e^{tF}$;
comparing $\varphi(0)$ with $\varphi(\pi/2) = \E F \cdot \E e^{tF}$ gives
Herbst's covariance inequality $\E F e^{tF} \le (\E F + tL^2)\E e^{tF}$.
With $\psi(t) = \ln \E e^{t(F - \E F)}$ this reads
$\psi'(t) \le tL^2$, and $\psi(0) = 0$ gives $\psi(t) \le t^2L^2/2$.
A general Lipschitz $F$ is the pointwise limit of the Gaussian smoothings
of $F$ clipped to $[-k-1, k+1]$, which are smooth with gradient at most
$L$; Fatou's lemma passes the bound to $F$. The tails are Chernoff's
bound, and the variance bound follows from the two exponential bounds at
$\pm t$, $\cosh(y) \ge 1 + y^2/2$, and $t \to 0$.
\end{proof}

\begin{theorem}[Chi-square tails \formal{GaussianAnalysis.comparison.chi\_square.gauss\_centered\_chi2\_abs\_tail}]\label{thm:chi2}
Let $a \in \R^m$ with $\abs{a_i} \le L$ and let $Q = \sum_i a_i (x_i^2 - 1)$,
$x \sim \gauss_m$. For $u > 0$ with $\norm{a}_2^2 + Lu > 0$,
\[
\Prob\{\abs{Q} \ge u\} \le 2\exp\Bigl(-\frac{u^2}{4(\norm{a}_2^2 + Lu)}\Bigr).
\]
\end{theorem}

\begin{proof}
For $2\abs{t}L < 1$, $\E e^{tQ} = \prod_i e^{-ta_i}(1 - 2ta_i)^{-1/2} \le \exp\bigl(t^2\norm{a}_2^2/(1 - 2\abs{t}L)\bigr)$,
from $\E e^{sx^2} = (1-2s)^{-1/2}$ for $s < 1/2$
\formal{GaussianAnalysis.comparison.chi\_square.integral\_normal\_expR\_sqr}
and $-\ln(1-u) - u \le u^2/(2(1-\abs{u}))$. Chernoff's bound with
$t = u/(2(\norm{a}_2^2 + Lu))$ gives each one-sided tail.
\end{proof}

\begin{proposition}[Linear forms \formal{GaussianAnalysis.comparison.gauss\_linear.gauss\_mxdot\_abs\_tail}]\label{prop:linear}
For a unit vector $a \in \R^m$ and $u > 0$, $\gauss_m\{\abs{\ip{x}{a}} \ge u\} \le 2e^{-u^2/2}$.
\end{proposition}

\begin{proof}
$\ip{x}{a} \sim N(0,1)$ by Proposition~\ref{prop:mulmx}, and the normal
tail is at most $e^{-u^2/2}$ on each side.
\end{proof}

\subsection{The spectral theorem and operator norms}\label{sec:spectral}

MathComp's spectral theorem is stated over an algebraically closed field;
the real-closed package is not available, so the library proves the real
spectral theorem directly.

\begin{theorem}[Real spectral theorem \formal{GaussianAnalysis.matrix.real\_spectral.symmetric\_spectral}]\label{thm:spectral}
Every symmetric $A \in \R^{m\times m}$ can be written $A = P^{\top} \diag(a) P$
with $P$ orthogonal ($P^{\top}P = \Id$) and $a \in \R^m$.
\end{theorem}

\begin{proof}
By induction on $m$. For $m \ge 1$, the quadratic form $v \mapsto v^{\top} A v$
attains its maximum on the compact unit sphere at some $v$
\formal{GaussianAnalysis.matrix.real\_spectral.rayleigh\_max}; a first-order
condition shows $A v = \lambda v$ with $\lambda = v^{\top} A v$
\formal{GaussianAnalysis.matrix.real\_spectral.symmetric\_eigenvector}:
indeed, for every $w$ and small $\varepsilon$, $(v + \varepsilon w)^{\top}A(v + \varepsilon w) \le \lambda \norm{v + \varepsilon w}_2^2$,
whose first-order term gives $\ip{w}{Av - \lambda v} = 0$. The Householder
reflection $H = \Id - 2ww^{\top}/\norm{w}_2^2$ with $w = v - e_1$
(or $H = \Id$ if $w = 0$) is orthogonal and symmetric and sends $v$ to $e_1$
\formal{GaussianAnalysis.matrix.real\_spectral.householder\_swap}; then
$HAH$ has $e_1$ as eigenvector, so it is block diagonal with blocks
$\lambda$ and a symmetric $(m-1)\times(m-1)$ matrix, to which the induction
hypothesis applies.
\end{proof}

\begin{corollary}[\formal{GaussianAnalysis.matrix.real\_spectral.spectral\_form}, \formal{GaussianAnalysis.matrix.real\_spectral.spectral\_trace}, \formal{GaussianAnalysis.matrix.real\_spectral.spectral\_mxdot}, \formal{GaussianAnalysis.matrix.real\_spectral.spectral\_form\_abs}]\label{cor:spectral}
With $A = P^{\top}\diag(a)P$: $v^{\top} A v = \sum_i a_i (Pv)_i^2$,
$\tr A = \sum_i a_i$, $\norm{A}_\Fro^2 = \sum_i a_i^2$, and if
$\abs{u^{\top} A w} \le b\norm{u}_2\norm{w}_2$ for all $u, w$ then
$\abs{a_i} \le b$ for all $i$. For $A \succeq 0$ all $a_i \ge 0$.
\end{corollary}

The operator norm of the library, $\norm{A}_\op = \sup_{\norm{v}=1}\norm{Av}_2$
\formal{GaussianAnalysis.matrix.opnorm.opnorm}, satisfies the following
two bounds, which replace spectral arguments on the traces of powers.

\begin{proposition}[\formal{GaussianAnalysis.matrix.opnorm.opnorm\_le\_trace\_exp}, \formal{GaussianAnalysis.matrix.opnorm.opnorm\_le\_form\_exp}]\label{prop:opnorm_trace}
Let $A$ be symmetric and $p \ge 1$, $t \ge 0$. If $\tr(A^{2p}) \le t^{2p}$ then
$\norm{A}_\op \le t$. If $\abs{v^{\top} A^{2p} v} \le t^{2p}\norm{v}_2^2$ for
every $v$ then $\norm{A}_\op \le t$.
\end{proposition}

\begin{proof}
For a unit $v$, $\norm{Av}_2^{2p} = (v^{\top}A^2 v)^p \le v^{\top} A^{2p} v$ by
Jensen's inequality along the spectral decomposition ($A^2 = P^{\top}\diag(a^2)P$
and $u \mapsto u^p$ is convex), and $v^{\top}A^{2p}v \le \tr A^{2p}$ since
$A^{2p} \succeq 0$.
\end{proof}

\begin{proposition}[\formal{GaussianAnalysis.matrix.real\_spectral.mxtrace\_ABAB\_le}]\label{prop:abab}
For symmetric $A, B$, $\tr(ABAB) \le \tr(A^2 B^2)$.
\end{proposition}

\begin{proof}
$\tr(A^2B^2) - \tr(ABAB) = \frac12 \tr\bigl((AB - BA)^{\top}(AB - BA)\bigr) \ge 0$,
using $(AB - BA)^{\top} = BA - AB$ for symmetric $A, B$.
\end{proof}

\subsection{Nets}\label{sec:nets}

\begin{theorem}[\formal{GaussianAnalysis.matrix.nets.net\_exists}, \formal{GaussianAnalysis.matrix.nets.unit\_sphere\_net}]\label{thm:nets}
Let $S \subseteq \R^m$ be contained in the ball of radius $r$ and let
$\varepsilon > 0$. There is a finite subset $N \subseteq S$, of cardinality at
most $(2e(1 + 2r/\varepsilon))^m$, such that every point of $S$ is within
$\varepsilon$ of a point of $N$. In particular the unit sphere has an
$\varepsilon$-net of unit vectors of cardinality at most $(2e(1 + 2/\varepsilon))^m$,
and if $\abs{u^{\top} A v} \le b$ for all $u, v$ in a $1/4$-net of the unit
sphere of unit vectors, then $\norm{A}_\op \le 2b$
\formal{GaussianAnalysis.matrix.nets.opnorm\_net\_sphere}.
\end{theorem}

\begin{proof}
Take a maximal $\varepsilon$-separated subset of $S$; the balls of radius
$\varepsilon/2$ around its points are disjoint and contained in the ball of
radius $r + \varepsilon/2$. The library counts the points with the volume
of a lattice box rather than Lebesgue measure: each ball of radius
$\varepsilon/2$ contains a cube of side $\varepsilon/\sqrt m$, and the cubes
of a grid of mesh $\varepsilon/\sqrt m$ that meet the big ball number at most
$(2e(1 + 2r/\varepsilon))^m$ (a binomial count). For the operator norm, write
a unit $x$ as $u + (x - u)$ with $u$ in the net: $\abs{x^{\top}Ay} \le b + \frac14\norm{A}_\op + \frac14 \norm{A}_\op + \frac1{16}\norm{A}_\op$,
and the supremum over unit $x, y$ gives $\norm{A}_\op \le 2b$ after
rearranging (the library's constant).
\end{proof}

\subsection{Moment inequalities}

\begin{theorem}[Paley--Zygmund \formal{GaussianAnalysis.probability.moments.paley\_zygmund}]\label{thm:pz}
On a probability space, let $X$ be square integrable and $0 \le \theta < \E X$.
Then $\Prob\{X \ge \theta\} \ge (\E X - \theta)^2 / \E X^2$.
\end{theorem}

\begin{proof}
For every $s \ge 0$ and real $z$, $2sz \le 2s\theta + s^2 z^2 + \one\{z \ge \theta\}$
(if $z < \theta$ the first term is below the second; otherwise use
$2sz \le s^2z^2 + 1$). Integrate at $z = X$ and optimize in $s$.
\end{proof}

\begin{theorem}[Lindeberg replacement \formal{GaussianAnalysis.probability.lindeberg.lindeberg\_replacement}]\label{thm:lindeberg}
Let $\mu_i, \nu_i$ ($i < n$) be probability measures on a measurable space
$T$, with product laws $\mu = \bigotimes \mu_i$, $\nu = \bigotimes \nu_i$ on
$n$-tuples, and let $F$ be bounded measurable on $n$-tuples. If replacing
the $i$-th coordinate of any tuple by a $\mu_i$-distributed or a
$\nu_i$-distributed coordinate changes the expectation of $F$ by at most
$B_i$, then $\abs{\int F\, d\mu - \int F\, d\nu} \le \sum_i B_i$.
\end{theorem}

\begin{proof}
Telescope through the hybrid laws $\mu_0 \otimes\cdots\otimes\mu_{i-1}\otimes\nu_i\otimes\cdots$
and apply Fubini at the replaced coordinate.
\end{proof}

The main proof uses its own instance of this principle for Gaussian rows
(Proposition~\ref{prop:replacement}); the library's general form is not
referenced by the main theorem.

\section{Preliminaries on the problem}\label{sec:prelim}

From now on the real field is that of the standard library, the sample
$x$ of $n$ points of $\R^d$ has the law $\Prob_{n,d}$ of
Definition~\ref{def:sample}, and $d = m + 1$ is positive whenever a
statement needs it.

\subsection{Features and Gram matrices}\label{sec:features}

\begin{definition}[Features]
The \emph{raw feature} of the point $x_i$ is $Q_i = x_i x_i^{\top} - \Id \in \Sym^d$
and its \emph{head feature} is
\[
W_i = \frac{x_i x_i^{\top} - \Id}{\sqd}
\tag*{\formal{Geometry.HeadFeature.paper\_main\_head\_feature}.}
\]
For $S \in \Sym^d$, $\ip{Q_i}{S} = x_i^{\top} S x_i - \tr S$, so the
constraints $x_i^{\top} S x_i = d$ of $\Fit_{n,d}$ read $\ip{W_i}{S} = (d - \tr S)/\sqd$
for all $i$. The \emph{feature Gram matrix} is the $n\times n$ matrix
$Q Q^{*} = (\ip{Q_i}{Q_j})_{ij}$ \formal{Feature.GramBlocks.paper\_main\_feature\_gram}
and, for a finite family $W = (W_j)_{j < m}$ of symmetric matrices, its
Gram matrix is $(\ip{W_j}{W_l})_{jl}$
\formal{Geometry.HeadCorrection.paper\_main\_head\_gram}. A matrix $M$ is
\emph{coercive with constant $\mu$} when $u^{\top} M u \ge \mu \norm{u}_2^2$
for every $u$ \formal{Geometry.HeadCoefficients.paper\_main\_head\_coercive}.
\end{definition}

\begin{lemma}[Blocks \formal{Feature.GramBlocks.paper\_main\_feature\_gram\_blocks}; paper (2.14)]\label{lem:blocks}
$Q Q^{*} = D D^{\top} + V V^{\top}$, where $D \in \R^{n\times d}$ has entries
$D_{ia} = x_{ia}^2 - 1$ and $V \in \R^{n \times d^2}$ has entries
$\sqrt 2\, x_{ia} x_{ib}$ at the positions $(a, b)$ with $a < b$ and $0$ at the
others. Moreover $(V V^{\top})_{ii} = \norm{x_i}_2^4 - \sum_a x_{ia}^4$
\formal{Feature.GramBlocks.paper\_main\_squarefree\_gram\_diag}, and
$V V^{\top} \preceq Q Q^{*}$ \formal{Feature.GramBlocks.paper\_main\_squarefree\_gram\_le}.
\end{lemma}

\begin{proof}
$\ip{Q_i}{Q_j} = \sum_{a} (x_{ia}^2 - 1)(x_{ja}^2 - 1) + 2\sum_{a < b} x_{ia}x_{ib}x_{ja}x_{jb}$,
by expanding $\ip{x_ix_i^{\top} - \Id}{x_jx_j^{\top} - \Id} = (x_i^{\top}x_j)^2 - \norm{x_i}^2 - \norm{x_j}^2 + d$
and $(x_i^{\top}x_j)^2 = \sum_{a,b} x_{ia}x_{ib}x_{ja}x_{jb}$.
\end{proof}

\begin{lemma}[From raw to head features \formal{Feature.SpectralBridge.paper\_main\_raw\_feature\_coercive}, \formal{Feature.SpectralBridge.paper\_main\_raw\_feature\_upper}; paper Corollary~2.9]\label{lem:bridge_features}
If $Q Q^{*}$ is coercive with constant $c d^2$ then the Gram matrix of
$(W_i)_{i<n}$ is coercive with constant $c d$; if $Q Q^{*} \preceq C d^2 \Id$
then that Gram matrix is at most $C d\, \Id$. Consequently, under both
bounds, $\norm{\sum_i c_i W_i}_\Fro \le \sqrt{Cd}\,\norm{c}_2$ for every
$c \in \R^n$ \formal{Feature.SpectralBridge.paper\_main\_raw\_feature\_operator}.
\end{lemma}

\begin{proof}
$\ip{W_i}{W_j} = \ip{Q_i}{Q_j}/d$, and $\norm{\sum_i c_i W_i}_\Fro^2 = c^{\top}(W W^{*})c$.
\end{proof}

\subsection{Bounded-condition cones, support functions and widths}\label{sec:cones}

\begin{definition}[Cones and sections]
For $\kappa \ge 1$, a matrix $S$ satisfies the \emph{condition $\kappa$}
when $\lambda \Id \preceq S \preceq \kappa\lambda \Id$ for some $\lambda > 0$
\formal{Geometry.HeadCorrection.paper\_main\_head\_condition}; this is the
paper's cone $\cone_\kappa = \{S \succeq 0 : \lambda_{\max}(S) \le \kappa \lambda_{\min}(S)\}$
without its apex. The \emph{cone sphere} is
$\cone_\kappa \cap \mathbb{S}_\Fro$ \formal{Geometry.HeadWidth.paper\_main\_head\_cone\_sphere},
and for a family $W = (W_j)_{j<m}$ the \emph{section} is
$\cone_\kappa \cap \mathbb{S}_\Fro \cap \{S : \ip{W_j}{S} = 0,\ j < m\}$
\formal{Geometry.HeadWidth.paper\_main\_head\_section}. The \emph{support
value} of $Z \in \Sym^d$ on a set $A$ is $h_A(Z) = \sup_{S \in A}\ip{Z}{S}$
\formal{Geometry.HeadWidth.paper\_main\_head\_support}.
\end{definition}

\begin{lemma}[Spectral box \formal{Geometry.HeadSpectral.paper\_main\_head\_condition\_box}; paper (3.13)]\label{lem:box}
If $S$ is on the cone sphere $\cone_\kappa \cap \mathbb{S}_\Fro$, then
$(\kappa\sqd)^{-1}\Id \preceq S \preceq (\kappa/\sqd)\Id$. In particular
$\abs{u^{\top}Sv} \le (\kappa/\sqd)\norm{u}_2\norm{v}_2$
\formal{Spectral.OperatorTransport.paper\_main\_matrix\_operator\_cone\_sphere}.
The cone sphere is nonempty: it contains $\Id/\sqd$
\formal{Geometry.HeadWidthCompact.paper\_main\_head\_cone\_sphere\_nonempty}.
\end{lemma}

\begin{proof}
From $\lambda\Id \preceq S \preceq \kappa\lambda\Id$ and
$1 = \norm{S}_\Fro^2 = \tr(S^2)$: $d\lambda^2 \le \tr(S^2) \le d\kappa^2\lambda^2$
gives $\lambda \le 1/\sqd \le \kappa\lambda$, hence the box. (The formal
proof avoids eigenvalues: $\tr(S^2) \le \kappa\lambda\tr S \le d\kappa^2\lambda^2$
and $\tr(S^2) \ge \lambda \tr S \ge d\lambda^2$ by $\tr(S^2) = \tr(S^{1/2}SS^{1/2})$
written with the Loewner bounds.)
\end{proof}

\begin{definition}[Gaussian width]
Let $Z$ be the standard Gaussian element of $\Sym^d$ (Section~\ref{sec:notation}),
seen as a function of $z \sim \gauss_{\Dd_d}$. For a nonempty $A \subseteq \mathbb{S}_\Fro$,
$w(A) = \E^{(\Dd_d)} h_A(Z)$ is its \emph{Gaussian width}. For the cone
sphere we write $w_\kappa = w(\cone_\kappa\cap\mathbb{S}_\Fro)$
\formal{Gaussian.Width.Condition.paper\_main\_gaussian\_condition\_width},
and for a section $w(\cone_\kappa \cap \mathbb{S}_\Fro \cap \ker W)$ is the
\emph{section width} \formal{Head.SectionWidthLaw.paper\_main\_head\_section\_width},
an extended real which is finite and at most $\E\norm{z}_2$ since
$\abs{h_A(Z)} \le \norm{Z}_\Fro = \norm{z}_2$ for $A \subseteq \mathbb{S}_\Fro$
\formal{Head.SectionWidthLaw.paper\_main\_head\_section\_width\_bound}.
\end{definition}

\begin{lemma}[Support functions are Lipschitz \formal{Geometry.HeadWidthLipschitz.paper\_main\_head\_support\_lipschitz}]\label{lem:support_lip}
For $A \subseteq \mathbb{S}_\Fro$ nonempty, $\abs{h_A(Z) - h_A(Y)} \le \norm{Z - Y}_\Fro$;
in the Gaussian coordinates this is the Euclidean distance of $z$ and $y$
\formal{Gaussian.Symmetric.PublicLaw.paper\_main\_public\_symmetric\_gaussian\_isometry}.
\end{lemma}

\begin{proposition}[Moreau decomposition \formal{Cone.Projection.paper\_main\_psd\_moreau}]\label{prop:moreau}
Every symmetric $Z$ has a unique decomposition $Z = Z_+ - Z_-$ with
$Z_+, Z_- \succeq 0$ and $\ip{Z_+}{Z_-} = 0$; $Z_+$ is the nearest point of
$\Spsd^d$ to $Z$ in Frobenius norm, the map $Z \mapsto Z_+$ is
$1$-Lipschitz, and $(-Z)_+ = Z_-$.
\end{proposition}

\begin{proof}
With $Z = P^{\top}\diag(a)P$ (Theorem~\ref{thm:spectral}) put
$Z_+ = P^{\top}\diag(\max(a,0))P$ and $Z_- = P^{\top}\diag(\max(-a,0))P$.
The nearest-point property is the variational characterization:
$Z_+ \succeq 0$ and $\ip{Z - Z_+}{S - Z_+} \le 0$ for every $S \succeq 0$,
since $\ip{-Z_-}{S} \le 0$ and $\ip{Z_-}{Z_+} = 0$. Uniqueness and the
Lipschitz bound follow from the characterization as for any projection on
a closed convex set.
\end{proof}

\begin{proposition}[Width of the positive semidefinite sphere \formal{Gaussian.Width.Psd.paper\_main\_gaussian\_psd\_width\_bounds}; paper Lemma~2.1]\label{prop:psdwidth}
$\E\norm{Z_+}_\Fro^2 = \Dd_d/2$ and
$\sqrt{\Dd_d/2 - 1} \le w(\Spsd^d \cap \mathbb{S}_\Fro) \le \sqrt{\Dd_d/2}$.
\end{proposition}

\begin{proof}
$\norm{Z}_\Fro^2 = \norm{Z_+}_\Fro^2 + \norm{Z_-}_\Fro^2$ by orthogonality,
$\E\norm{Z}_\Fro^2 = \E\norm{z}_2^2 = \Dd_d$, and $Z \mapsto -Z$ preserves
the law (Proposition~\ref{prop:gaussE}) and exchanges $Z_+$ with $Z_-$;
hence $\E\norm{Z_+}_\Fro^2 = \Dd_d/2$
\formal{Gaussian.Width.Psd.paper\_main\_gaussian\_psd\_projection\_second\_moment}.
Next, $\max(h(Z), 0) = \norm{Z_+}_\Fro$ for $h = h_{\Spsd^d\cap\mathbb{S}_\Fro}$
\formal{Cone.ProjectionSupport.paper\_main\_psd\_sphere\_support\_max}:
for $S \succeq 0$ of unit norm, $\ip{Z}{S} = \ip{Z_+}{S} - \ip{Z_-}{S} \le \norm{Z_+}_\Fro$,
with equality at $S = Z_+/\norm{Z_+}_\Fro$ when $Z_+ \ne 0$, and $h(Z) \le 0$
when $Z_+ = 0$, that is $Z \preceq 0$. So $\E\max(h,0)^2 = \Dd_d/2$. The
function $h$ is $1$-Lipschitz (Lemma~\ref{lem:support_lip}), so is
$\max(h, 0)$, and the Gaussian Poincar\'e inequality (Theorem~\ref{thm:conc})
gives $\operatorname{Var}\max(h,0) \le 1$, whence
$(\E\max(h,0))^2 \ge \Dd_d/2 - 1$; also $\E h \le \E\max(h,0)$ and
$(\E h)^2 \le \E h^2 \le \E\max(h,0)^2$ once $\E h \ge 0$, which holds since
$h(Z) \ge \ip{Z}{\Id}/\sqd = \tr Z/\sqd$ has mean $0$. Finally
$\E\max(h,0) \ge \E h$ needs the reverse: $\max(h,0) - h = \max(-h, 0) \le \norm{Z_-}_\Fro \one\{Z_+ = 0\}$
vanishes unless $Z \preceq 0$, and then equals $-h(Z) \le \norm{Z}_\Fro\one\{Z\preceq 0\}$;
the formal proof instead bounds $w$ from below through
$\E h \ge \sqrt{\E\max(h,0)^2 - 1}$, using $\E\max(h,0)^2 - 1 \le (\E h)^2$
for a $1$-Lipschitz $h$ with $\E h \ge 0$
\formal{Gaussian.Width.SignedSupport.paper\_main\_gaussian\_mean\_positive\_square\_lower}.
\end{proof}

The cited width-dimension inequality of the paper (its line 365) is thus
replaced by the Poincar\'e inequality for the Lipschitz function $\max(h,0)$.

\begin{proposition}[Width of bounded-condition cones \formal{Gaussian.Width.ConditionBound.paper\_main\_gaussian\_condition\_width\_lower}, \formal{Gaussian.Width.ConditionBound.paper\_main\_gaussian\_condition\_width\_large}; paper Lemma~2.2]\label{prop:conewidth}
For every $s > 0$ and $d \ge 1$, with $\kappa = (20 + s)/s$,
\[
w_\kappa \ge \Bigl(\frac12 - s\Bigr) d - 3 - \frac{864}{d} .
\]
Consequently, for every $\varepsilon > 0$ there is $\kappa \ge 1$ with
$w_\kappa \ge (1/2 - \varepsilon) d$ for all large $d$.
\end{proposition}

The paper proves this through the semicircle law and convergence in
probability; the formal proof is a direct computation with a crude
operator-norm event.

\begin{proof}
\emph{The shift.} Let $Z$ be symmetric with $Z_+ \preceq C\sqd\,\Id$, let
$t = s\sqd$ and $Q = Z_+ + t\Id$. Then $t\Id \preceq Q \preceq (C + s)\sqd\,\Id = \kappa t\Id$
with $\kappa = (C+s)/s$, so $P = Q/\norm{Q}_\Fro$ lies on the cone sphere
\formal{Cone.Condition.paper\_main\_cone\_shift\_scaled}. Its support value is
$\ip{Z}{P} = (\norm{Z_+}_\Fro^2 + t\tr Z)/\norm{Q}_\Fro$ (using
$\ip{Z}{Z_+} = \norm{Z_+}_\Fro^2$), and $\norm{Q}_\Fro \le \norm{Z_+}_\Fro + t\sqd$.
Writing $a = \norm{Z_+}_\Fro$ and $b = \abs{\tr Z}$: if $a^2 - tb \ge 0$ then
$h(Z) \ge (a^2 - tb)/(a + t\sqd) \ge a - t\sqd - b/\sqd$, since
$(a - t\sqd - b/\sqd)(a + t\sqd) = a^2 - t^2d - b(a + t\sqd)/\sqd \le a^2 - tb$;
if $a^2 - tb < 0$ then $a < \sqrt{tb}$, so $a - t\sqd - b/\sqd \le -b/\sqd$
(by the inequality $\sqrt{tb} \le \frac12(t\sqd + b/\sqd)$ between the
geometric and the arithmetic means), while
$h(Z) \ge \ip{Z}{\Id/\sqd} = \tr Z/\sqd \ge -b/\sqd$. In both cases
\[
h_{\cone_\kappa\cap\mathbb{S}_\Fro}(Z) \ge \norm{Z_+}_\Fro - s d - \frac{\abs{\tr Z}}{\sqd}
\tag*{\formal{Cone.Condition.paper\_main\_cone\_shift\_cone\_support}.}
\]

\emph{The operator-norm event.} With $C = 20$: outside a measurable set $A$
with $\gauss_{\Dd_d}(A) \le 2e^{-10d}$, the Gaussian $Z$ satisfies
$\abs{u^{\top}Zv} \le 20\sqd\,\norm{u}_2\norm{v}_2$ for all $u, v$
\formal{Gaussian.Symmetric.OperatorTail.paper\_main\_symmetric\_gaussian\_operator\_good\_set},
hence $Z_+ \preceq 20\sqd\,\Id$. (Proof: $u^{\top}Zv = \ip{Z}{(uv^{\top} + vu^{\top})/2}$
is a Gaussian linear form of $z$ with variance at most $\norm{u}^2\norm{v}^2$
\formal{Spectral.SymmetricBilinear.paper\_main\_symmetric\_bilinear\_coefficients\_le};
a $1/4$-net of the unit sphere of $\R^d$ with $\ln\abs{\text{net}} \le 20d\ln 2$
\formal{Spectral.OperatorNet.paper\_main\_symmetric\_operator\_net\_quarter\_stdlib},
Proposition~\ref{prop:linear} at level $10\sqd$ and a union bound give
$2\abs{\text{net}}^2 e^{-50d} \le 2e^{-10d}$, and Theorem~\ref{thm:nets}
doubles the level.)

\emph{The expectation.} Using $\abs{\tr Z}/\sqd \le 1/2 + (\tr Z)^2/(2d)$
and, on $A$, the bound $\abs{h} \le \norm{Z}_\Fro$ together with
$\norm{Z_+}_\Fro \le \norm{Z}_\Fro \le \norm{Z}_\Fro^2/(2a) + a/2$ for $a > 0$,
one gets pointwise, for every $Z$
\formal{Gaussian.Width.Condition.paper\_main\_gaussian\_condition\_point\_bound},
\[
\norm{Z_+}_\Fro \le h(Z) + sd + \tfrac12 + \frac{(\tr Z)^2}{2d} + \frac{\norm{Z}_\Fro^2}{2a} + 2a\one_A(Z).
\]
Take expectations: $\E(\tr Z)^2 = d$ \formal{Gaussian.Symmetric.PublicLaw.paper\_main\_public\_symmetric\_gaussian\_moments},
$\E\norm{Z}_\Fro^2 = \Dd_d \le d^2$, and
$\E\norm{Z_+}_\Fro \ge \sqrt{\Dd_d/2 - 1} \ge d/2 - 1$
\formal{Gaussian.Width.ConditionBound.paper\_main\_gaussian\_psd\_projection\_mean\_lower}.
With $a = d^2$ this gives
$w_\kappa \ge (1/2 - s)d - 2 - \Dd_d/(2d^2) - 4d^2 e^{-10d}$
\formal{Gaussian.Width.ConditionBound.paper\_main\_gaussian\_condition\_width\_finite},
and $\Dd_d/(2d^2) \le 1/2$, $d^2e^{-10d} \le 216/d$ give the claim. For the
second statement take $s = \varepsilon/2$ and $d$ large.
\end{proof}

\subsection{Elementary estimates on the sample}\label{sec:data}

\begin{lemma}[Radius \formal{Gaussian.Sample.Radius.paper\_main\_data\_radius\_probability}, \formal{Gaussian.Sample.Radius.paper\_main\_data\_radius\_good}; paper Lemma~2.10, (2.17)]\label{lem:radius}
The event that some point has $\norm{x_i}_2 \ge 2\sqd$ is measurable with
probability at most $n e^{-d/2}$; outside it, $\norm{x_i}_2^2 \le 4d$ for
every $i < n$.
\end{lemma}

\begin{proof}
$x \mapsto \norm{x}_2$ is $1$-Lipschitz on $\R^d$ with mean at most $\sqd$
(Jensen), so $\gauss_d\{\norm{x}_2 \ge 2\sqd\} \le e^{-d/2}$
(Theorem~\ref{thm:conc}); a union bound over the $n$ points, whose law is
$\gauss_d$ each (Proposition~\ref{prop:gaussE}), concludes.
\end{proof}

\begin{lemma}[Coordinates \formal{Data.SampleTail.paper\_main\_data\_coordinate\_tail}; paper Lemma~2.10, (2.18)]\label{lem:coord}
For $u > 0$, the event that some coordinate satisfies $\abs{x_{ia}} \ge u$
is measurable with probability at most $2nd\, e^{-u^2/2}$.
\end{lemma}

\begin{proof}
Each coordinate is a unit linear form of the sample (Proposition~\ref{prop:linear});
union bound.
\end{proof}

The formal proof uses this lemma at the level $u = d^{1/4}$
\formal{Clipping.Limits.paper\_main\_coordinate\_clipping\_tends\_to\_zero},
where the probability is $2nd\,e^{-\sqd/2} \to 0$ for $n \le Cd^2$; it is
the \emph{coordinate clipping event}, and all the coordinate-wise clippings
below are at this level.

\begin{lemma}[Coherence \formal{Data.SampleTail.paper\_main\_data\_coherence\_tail}, \formal{Data.CoherenceLimit.paper\_main\_data\_coherence\_log\_good}; paper Lemma~2.10, (2.19)]\label{lem:coherence}
$\ip{W_i}{W_j} = \bigl((x_i^{\top}x_j)^2 - \norm{x_i}_2^2 - \norm{x_j}_2^2 + d\bigr)/d$.
For $B > 0$, the event that some pair $i < j$ has $\norm{x_i}_2^2 \le 4d$ and
$\abs{x_i^{\top}x_j} \ge B$ is measurable with probability at most
$2n^2 e^{-B^2/(8d)}$. Outside it and outside the radius event,
$\abs{\ip{W_i}{W_j}} \le B^2/d + 9$ for all $i \ne j$. With
$B^2 = 64d\ln d$ and $d > 1$ the probability is at most $2n^2/d^8$ and the
bound is $73\ln d$.
\end{lemma}

\begin{proof}
Condition on $x_i$ through Fubini for the block of $x_j$: for fixed
$v = x_i$ with $0 < \norm{v}_2^2 \le 4d$, the section is the event
$\abs{\ip{v/\norm{v}_2}{x_j}} \ge B/\norm{v}_2$, of probability at most
$2e^{-B^2/(2\norm{v}^2)} \le 2e^{-B^2/(8d)}$ (Proposition~\ref{prop:linear});
the formal proof splits the sample before row $j$ and bounds the section
probability \formal{Head.LawBounds.paper\_main\_split\_points\_section\_bound}.
Union bound over pairs.
\end{proof}

\begin{lemma}[Sparse column norms \formal{Data.HeadOperatorUniform.paper\_main\_data\_head\_operator\_uniform}, \formal{Head.DataLog.paper\_main\_head\_data\_log\_event}; paper Lemma~2.10, (2.20)]\label{lem:sparse}
For $u > 0$ there is a measurable event of probability at most $2e^{-u}$
outside which, for every $m \le K$ and every injective selection
$J$ of $m$ points, the $d\times m$ matrix $X_J$ of the selected points satisfies
\[
X_J X_J^{\top} \preceq 8\bigl(20\ln 2\,(d + K) + \ln(K+1) + K\ln(n+1) + u\bigr)\,\Id .
\]
Consequently, for $K = K_0 d\ln d$ with $K_0 \ge 0$ and $n \le Cd^2$ there is
$L = L(C, K_0)$ such that, outside an event of probability at most $2/d$,
$X_J X_J^{\top} \preceq L d(\ln d)^2\,\Id$ for every $J$ with $\abs{J} \le K_0 d\ln d$.
\end{lemma}

\begin{proof}
For a fixed selection, $u^{\top}X_J v$ is a unit linear form of the sample
when $u, v$ are unit (the selection is injective), so for $1/4$-nets $z$ of
$\mathbb{S}^{d-1}$ and $w$ of $\mathbb{S}^{m-1}$ the event
$\max_{y,t}\abs{z_y^{\top}X_Jw_t} \ge b$ has probability at most
$2\abs{z}\abs{w}e^{-b^2/2}$ \formal{Data.HeadOperator.pmdo\_test\_tail},
and outside it $\norm{X_J}_\op \le 2b$ (Theorem~\ref{thm:nets}, with nets
of entropy $\ln\abs{z} + \ln\abs{w} \le 20\ln2\,(d + m)$
\formal{Gordon.OperatorNet.paper\_main\_gordon\_operator\_discretization}).
Take $b^2 = 2(20\ln 2\,(d + K) + t)$ with $t = \ln(K+1) + K\ln(n+1) + u$;
the number of selections of at most $K$ points is at most $(K+1)(n+1)^K$,
and the union bound gives $2e^{-u}$. The second statement is the first with
$u = \ln d$.
\end{proof}

\begin{lemma}[Hanson--Wright for Gaussian vectors \formal{Gaussian.Orthogonal.HansonWright.paper\_main\_gaussian\_quadratic\_tail}; paper Lemma~2.11]\label{lem:hw}
Let $P \in \Sym^d$ with $\abs{u^{\top}Pv} \le b\norm{u}_2\norm{v}_2$ for all
$u, v$, let $x \sim \gauss_d$ and $u > 0$ with $\norm{P}_\Fro^2 + bu > 0$. Then
\[
\Prob\bigl\{\abs{x^{\top}Px - \tr P} \ge u\bigr\} \le 2\exp\Bigl(-\frac{u^2}{4\norm{P}_\Fro^2 + 4bu}\Bigr).
\]
For $W = (xx^{\top} - \Id)/\sqd$ and $P$ on the cone sphere $\cone_\kappa\cap\mathbb{S}_\Fro$, this gives
$\Prob\{\abs{\ip{W}{P}} \ge r\} \le 2\exp\bigl(-dr^2/(4 + 4\kappa r)\bigr)$ for $r > 0$
\formal{Gaussian.Sample.ResidualTail.paper\_main\_gaussian\_residual\_tail}.
\end{lemma}

\begin{proof}
Write $P = O^{\top}\diag(a)O$ (Theorem~\ref{thm:spectral}); then
$x^{\top}Px - \tr P = \sum_i a_i((Ox)_i^2 - 1)$ with $\sum a_i^2 = \norm{P}_\Fro^2$
and $\abs{a_i} \le b$ (Corollary~\ref{cor:spectral}), and $Ox \sim \gauss_d$
(Proposition~\ref{prop:mulmx}). Theorem~\ref{thm:chi2} concludes. For the
second claim, $\ip{W}{P} = (x^{\top}Px - \tr P)/\sqd$, $\norm{P}_\Fro = 1$ and
$b = \kappa/\sqd$ (Lemma~\ref{lem:box}).
\end{proof}

\subsection{The spectral edges of the feature map}\label{sec:edge}

This section proves the two-sided bounds on $QQ^{*}$ that the two phases
use. The paper relies on the extreme-eigenvalue theorem of Kogan, Nandy and
Huang and on the matrix Bernstein inequality; the formal proof is
self-contained. It gives the upper edge for every $n \le d^2$ and the lower
edge only for $4n \le d^2$, which is all the satisfiable phase needs.

\subsubsection{The diagonal block}

\begin{proposition}[\formal{Feature.DiagonalLimit.paper\_main\_diag\_trace\_event\_bound}, \formal{Feature.DiagonalLimit.paper\_main\_diag\_trace\_event\_markov}; paper Lemma~2.7]\label{prop:diag}
Let $Z = D^{\top}D - 2n\Id$. The event $\tr Z^4 > n^4$ is measurable, and
outside it $D^{\top}D \preceq 3n\Id$ and $DD^{\top} \preceq 3n\Id$. For $n \ge 1$,
\[
\Prob\{\tr Z^4 > n^4\} \le \frac{n(2^{21}d^4 + 16d) + 3n(n-1)d(4d + 52)^2}{n^4},
\]
which is at most $\bigl((2^{21} + 16)/c^3 + 9408/c^2\bigr)/d$ when $n \ge cd^2$.
\end{proposition}

\begin{proof}
\emph{Deterministic part.} $Z$ is symmetric, so $\tr Z^4 \le n^4$ gives
$-n\Id \preceq Z \preceq n\Id$ (Proposition~\ref{prop:opnorm_trace} with
$p = 2$) and $D^{\top}D \preceq 3n\Id$; then $DD^{\top} \preceq 3n\Id$ since
$\norm{D^{\top}u}_2^2 = u^{\top}DD^{\top}u$ and $\norm{Dv}_2^2 \le 3n\norm{v}_2^2$
give $\norm{DD^{\top}u}_2^2 = \norm{D(D^{\top}u)}_2^2 \le 3n\, u^{\top}DD^{\top}u$,
hence $u^{\top}DD^{\top}u \le 3n\norm{u}_2^2$ by Cauchy--Schwarz
\formal{Feature.DiagonalTrace.paper\_main\_gram\_transfer}.

\emph{Moments.} Let $v(t) = t^2 - 1$; for $t \sim N(0,1)$, $\E v = 0$,
$\E v^2 = 2$, $\E v^4 = 60$ and $\E v^8 \le 2^{21}$ (the Gaussian moments
$\E t^{2k} = (2k-1)!!$). Let $u_i \in \R^d$ have entries $v(x_{ia})$ and
$M_i = u_iu_i^{\top} - 2\Id$, so $Z = \sum_i M_i$
\formal{Feature.DiagonalMoment.paper\_main\_diag\_block\_centered}. One row:
$\E M_i = 0$, $\E M_i^2 = (4d + 52)\Id$ (the diagonal entry $a$ of $M^2$ is
$v_a^2\sum_c v_c^2 - 4v_a^2 + 4$, of mean $60 + 2(d-1)\cdot 2 - 8 + 4 = 4d + 52$;
off-diagonal entries have mean $0$), and
$\E\tr M_i^4 \le 2^{21}d^4 + 16d$ (expand $M^4 = \norm{u}^4 uu^{\top} - \ldots$
through $(uu^{\top} - 2\Id)^4 = (\norm{u}_2^2 - 2)^4 \frac{uu^{\top}}{\norm{u}_2^2} + 16(\Id - \frac{uu^{\top}}{\norm{u}_2^2})$, so
$\tr M^4 = (\norm{u}_2^2 - 2)^4 + 16(d - 1)$ \formal{Feature.DiagonalTrace.paper\_main\_rank\_one\_trace\_fourth},
and $\E(\sum_a v_a^2 - 2)^4 \le d^3\sum_a\E(v_a^2 - 2/d)^4 \le \ldots \le 2^{21}d^4$).
By induction over the rows, with the expansion
$\tr(S + M)^4 = \tr S^4 + 4\tr S^3M + 4\tr S^2M^2 + 2\tr SMSM + 4\tr SM^3 + \tr M^4$
\formal{Feature.DiagonalTrace.paper\_main\_trace\_fourth\_expand}, the
independence of the last row, $\E M = 0$, $\E M^2 = (4d+52)\Id$ and
$\tr(SMSM) \le \tr(S^2M^2)$ (Proposition~\ref{prop:abab}):
$\E\tr S_k^2 = kd(4d + 52)$ and
$\E\tr S_{k+1}^4 \le \E\tr S_k^4 + 6(4d+52)\E\tr S_k^2 + 4\E\tr(S_k\,\E M^3) + \E\tr M^4$,
where $\E\tr(S_kC_3) = \tr(\E S_k\, C_3) = 0$. Summing gives the stated bound
\formal{Feature.DiagonalMoment.paper\_main\_diag\_trace\_moment}.

\emph{Probability.} The function $\tr Z^4$ of the sample is a polynomial,
hence measurable, and Markov's inequality with the moment bound gives the
claim. The rate for $n \ge cd^2$ is arithmetic.
\end{proof}

\subsubsection{The upper edge of the square-free block}

Let $G = VV^{\top}$, with entries $G_{ij} = \sum_{a\ne b}x_{ia}x_{ib}x_{ja}x_{jb}$
(the sum over ordered pairs absorbs the factor $2$ of $\sqrt2^2$)
\formal{SquareFree.WalkExpansion.paper\_main\_squarefree\_gram\_list}, and
let $\mathcal{W} = G - \diag(G)$ be its off-diagonal part.

\begin{proposition}[Walk expansion \formal{SquareFree.WalkExpansion.paper\_main\_squarefree\_offdiag\_walks}]\label{prop:walks}
For $L \ge 1$, $\E\tr\mathcal{W}^L$ is the sum, over closed walks
$i_0, i_1, \dots, i_{L-1}, i_0$ of rows with consecutive rows distinct and
over sequences of ordered label pairs $(a_s, b_s)$, $a_s \ne b_s$, $s < L$,
of the Gaussian moment of the monomial
$\prod_{s<L} x_{i_sa_s}x_{i_sb_s}x_{i_{s+1}a_s}x_{i_{s+1}b_s}$. The moment of a
monomial $\prod_c x_c^{n_c}$ in the coordinates is $\prod_c (n_c - 1)!!$ if
every $n_c$ is even and $0$ otherwise.
\end{proposition}

\begin{proof}
Expand the entries of $\mathcal{W}^L$ along paths and the entries of $G$
along label pairs; the moment of a monomial is the product of the
one-dimensional moments by Fubini (Proposition~\ref{prop:gaussE}).
\end{proof}

\begin{proposition}[Trace moment of the off-diagonal part \formal{SquareFree.EdgeUpper.paper\_main\_squarefree\_offdiag\_trace\_le}]\label{prop:tracemoment}
If $0 < n \le d^2$, $L \ge 32$ and $16L^2(2L)^{45} \le d$, then
$\E\tr\mathcal{W}^L \le n(52d^2)^L$.
\end{proposition}

\begin{proof}
A walk with a nonzero moment is \emph{even}: for every row $r$ and color
(coordinate index) $c$, the degree
$\sum_s(\one\{i_s = r\} + \one\{i_{s+1} = r\})(\one\{a_s = c\} + \one\{b_s = c\})$
is even \formal{SquareFree.WalkWeight.pmsw\_word\_moment\_even}. For such a
walk with $R$ distinct rows and $C$ distinct colors, the moment is at most
$(4L)^{2e}$ where $e = 2L + 2 - 2R - C \ge 0$ is the \emph{deficiency}
\formal{SquareFree.WalkBridge.pmsw\_bridge\_weight}: the moment is at most
$(4L)^X$ with $X = \sum_{r,c}(\deg(r,c)/2 - \one\{\deg(r,c) > 0\})$ the
multiplicity excess, since $(m-1)!! \le m^{m/2 - 1}$ for even $m \ge 2$ and
the monomial has $4L$ factors \formal{SquareFree.WalkWeight.pmsw\_moment\_le},
and $X \le 2e$ \formal{SquareFree.WalkExcess.pmsc\_X\_le}, by an induction
that removes a row visited once (contraction or excursion). The same
induction gives $2R + C \le 2L + 2$ \formal{SquareFree.WalkCount.pmsc\_rows\_colors\_le},
the degree-two case of the Kogan--Nandy--Huang count, so $e \ge 0$.

The walks are then counted by a \emph{twin code}
\formal{SquareFree.WalkShapeCode}: each step is encoded either as a closer
(it repeats the label pair of an earlier open step at the current row,
given by its rank among at most three such steps, or by its index) or as
an explicit step (its target row, fresh or by index, and its two colors,
fresh or by an earlier occurrence). The code determines the walk from its
root \formal{SquareFree.WalkShapeCode.pmsk\_inj}. A walk in which every
non-root row is visited at least twice (a \emph{terminal} walk) has at most
$45e$ index slots in its code \formal{SquareFree.WalkShapeBound.pmsb\_point},
since lonely steps number at most $27e$ \formal{SquareFree.WalkShapeTwin.pmst\_lonely\_le};
with $n$ choices for a fresh row, $d$ for a fresh color and $y = 2(L+1)$
for an index slot, the weighted sum of the codes of terminal walks of
length $L$ is at most $n(13d^2)^L$ when $n \le d^2$ and
$(4L)^2y^{45} \le d$ \formal{SquareFree.WalkShapeBound.pmsb\_T\_le}.
Non-terminal walks reduce to terminal ones by removing rows visited once
\formal{SquareFree.WalkShapeOnce}; the position recursion
\formal{SquareFree.WalkShapeRecursion.pmsr\_bound} costs a factor $2^L$
and the weight $B^{2e}$ is accounted for by the reduction, giving the bound
$2n(4L)^4(26d^2)^L$ on the sum over all walks \formal{SquareFree.WalkShapeBound.pmsb\_S\_le}.
Finally $2(4L)^4 \le 2^L$ for $L \ge 32$ \formal{SquareFree.EdgeUpper.pmeu\_absorb}
turns $2n(4L)^4(26d^2)^L$ into $n(52d^2)^L$.
\end{proof}

\begin{proposition}[Upper edge \formal{SquareFree.EdgeUpperLimit.paper\_main\_squarefree\_upper\_bound}, \formal{SquareFree.EdgeUpperLimit.paper\_main\_squarefree\_upper\_event\_tends\_to\_zero}; paper Lemma~2.6 and Proposition~2.8, upper part]\label{prop:upperedge}
Let $L(d) = 2^{k+3}$ with $k = \lfloor\log_2\lfloor\log_2 d\rfloor\rfloor$ and
$t = 104d^2$. The event
$E_d^{\mathrm{up}} = \{\tr\mathcal{W}^{L(d)} > t^{L(d)}\} \cup \{\exists i : \abs{G_{ii}/d^2 - 1} > 1\}$
is measurable; if $0 < n(d) \le d^2$ eventually then $\Prob(E_d^{\mathrm{up}}) \to 0$;
and outside $E_d^{\mathrm{up}}$, for $d \ge 1$, $VV^{\top} \preceq 106d^2\,\Id$.
Consequently, with Proposition~\ref{prop:diag}, outside an event of
probability tending to zero, $QQ^{*} \preceq (3n + 106d^2)\Id \preceq 109d^2\,\Id$
\formal{Feature.UpperProbability.paper\_main\_feature\_upper\_edge}.
\end{proposition}

\begin{proof}
For $d \ge 2^{9600}$, $L(d) \ge 32$, $16L^2(2L)^{45} \le d$ and $d^3 \le 2^{L}$
\formal{SquareFree.EdgeUpperLimit.pmel\_L\_facts}. Markov's inequality and
Proposition~\ref{prop:tracemoment} give
$\Prob\{\tr\mathcal{W}^L > t^L\} \le n(52d^2)^L/(104d^2)^L = n2^{-L} \le n/d^3 \le 1/d$.
Outside the event, since $L$ is even (a power of two),
$-t\Id \preceq \mathcal{W} \preceq t\Id$ (Proposition~\ref{prop:opnorm_trace}).
The diagonal: $G_{ii} = \norm{x_i}_2^4 - \sum_ax_{ia}^4$ (Lemma~\ref{lem:blocks});
outside the row event $\{\abs{\norm{x_i}_2^2 - d} \ge \eta d\}$, of
probability at most $2\exp(-(\eta d)^2/(4d + 4\eta d))$ by Theorem~\ref{thm:chi2}
with the weights $\one_{\text{row } i}$, and outside the coordinate event at
level $U = d^{1/4}$ (so that $\sum_a x_{ia}^4 \le U^2\norm{x_i}_2^2$),
$\abs{G_{ii}/d^2 - 1} < \eta(2 + \eta) + U^2(1 + \eta)/d \le 1$ for
$\eta = 1/6$ and large $d$ \formal{SquareFree.Diagonal.paper\_main\_row\_diag\_bound};
both events have probability tending to zero for $n \le d^2$
\formal{SquareFree.DiagonalLimit.paper\_main\_squarefree\_diag\_tends\_to\_zero}.
So $G_{ii} \le 2d^2$ and $v^{\top}Gv = v^{\top}\mathcal{W}v + \sum_i v_i^2G_{ii} \le (104 + 2)d^2\norm{v}_2^2$.
\end{proof}

\subsubsection{The lower edge of the square-free block}

Here $X \in \R^{n\times d^2}$ has entries $X_{i,(a,b)} = x_{ia}x_{ib}$ for
$a < b$ and $0$ otherwise, so that $VV^{\top} = 2XX^{\top} = 2\sum_c R_c$ with
$R_c = X_cX_c^{\top}$ for the columns $c = (a, b)$
\formal{Column.ConcentrationFrame.pmccf\_rankone\_sum}. Let
$M = \binom d2$, $\kappa = (nM)^{-1/2}$ and $s = \kappa n = \sqrt{n/M}$, so
that $\kappa^2nM = 1$ and $1/\kappa = sM$.

\begin{definition}[Non-backtracking matrices \formal{SquareFree.EdgeLower.pmlo\_T}]
Let $D_c = \diag(R_c)$ and $P_c = \kappa(R_c - D_c)$. Define $T_0 = \Id$,
$T_0^c = 0$ and, for $\ell \ge 0$, $T^c_{\ell+1} = P_c(T_\ell - T^c_\ell)$,
$T_{\ell+1} = \sum_cT^c_{\ell+1}$. Thus $T_1 = \sum_cP_c = \kappa(\sum_cR_c - \sum_cD_c)$,
and $(T_\ell)_{ij}$ is $\kappa^\ell$ times the sum over walks of $\ell$ steps
from $i$ to $j$ with consecutive rows distinct and consecutive columns
distinct of the products $\prod_sX_{r_sc_s}X_{r_{s+1}c_s}$
\formal{SquareFree.WalkNB.pmnb\_T\_walks}; $T_\ell$ is symmetric.
\end{definition}

\begin{lemma}[The recursion of Kogan, Nandy and Huang \formal{SquareFree.EdgeLower.pmlo\_knh}]\label{lem:knh}
Let $S_j = \sum_cP_c^j$, $E_k = S_{k+2} - s^{k+1}T_1 - s^k\Id$ and
$\Psi_\ell = \sum_{k<\ell}(-1)^kE_kT_{\ell-1-k}$. Then for every $\ell \ge 0$,
\[
(T_1 - s\Id)\,T_{\ell+1} = T_{\ell+2} + T_\ell + \Psi_{\ell+1} + s\Psi_\ell .
\]
\end{lemma}

\begin{proof}
Unrolling $T^c_{\ell} = \sum_{k<\ell}(-1)^kP_c^{k+1}T_{\ell-1-k}$ gives
$T_{\ell+1} = T_1T_\ell - \sum_{k<\ell}(-1)^kS_{k+2}T_{\ell-1-k}$
\formal{SquareFree.EdgeLower.pmlo\_rec\_S}; substitute
$S_{k+2} = E_k + s^{k+1}T_1 + s^k\Id$ and collect the terms.
\end{proof}

\begin{lemma}[Chebyshev step \formal{SquareFree.EdgeLower.pmlo\_cheb\_box}]\label{lem:cheb}
Let $B$ be symmetric, $0 \le s \le 1$, and let $(T_m)$, $(R_m)$ be matrices
with $T_0 = \Id$, $BT_0 = T_1 - sT_0$ and $BT_{m+1} = T_{m+2} + T_m + R_{m+1}$.
Let $K = 2^{k+1}$. If $\abs{v^{\top}T_mv} \le \tau\norm{v}_2^2$ for $m \le K$,
$\norm{B}_\op \le \beta$, $\norm{R_m}_\op \le \rho$ for $0 < m < K$, and
$2^K\tau + \rho_K \le \theta^K$ where $\rho_0 = 0$,
$\rho_{m+1} = \beta\rho_m + 2^m\rho$, then $-\theta\Id \preceq B \preceq \theta\Id$.
\end{lemma}

\begin{proof}
Define coefficients $\gamma_{K,m}$ by the recursion of the Chebyshev-type
polynomials attached to the three-term relation; then
$B^K = \sum_{m\le K}\gamma_{K,m}T_m + \mathrm{rem}_K$ with
$\mathrm{rem}_{K+1} = B\,\mathrm{rem}_K + \sum_{m<K}\gamma_{K,m+1}R_{m+1}$
\formal{SquareFree.EdgeLower.pmlo\_cheb}, and $\sum_m\abs{\gamma_{K,m}} \le 2^K$
for $0 \le s \le 1$ \formal{SquareFree.EdgeLower.pmlo\_gam\_abs}. Hence
$\norm{\mathrm{rem}_K}_\op \le \rho_K$ and
$\abs{v^{\top}B^Kv} \le (2^K\tau + \rho_K)\norm{v}_2^2 \le \theta^K\norm{v}_2^2$;
Proposition~\ref{prop:opnorm_trace} (form version, $K$ even) concludes.
\end{proof}

\begin{proposition}[Box for the first non-backtracking matrix \formal{SquareFree.Box.paper\_main\_squarefree\_box\_from\_remainders}]\label{prop:nbbox}
Let $0 \le s \le 1$, $t \ge 1$, $\eta, \theta \ge 0$, $K = 2^{k+1}$ and
$p = 2^{j+2}$. Suppose $\tr T_m^p \le t^p$ for $1 \le m \le K$, that every
$E_m$ with $m < K$ has $\norm{E_m}_\op \le \eta$, and that
$2^Kt + \rho(t+1, 3K\eta t, K) \le \theta^K$ with $\rho$ as in Lemma~\ref{lem:cheb}.
Then $-\theta\Id \preceq T_1 - s\Id \preceq \theta\Id$.
\end{proposition}

\begin{proof}
The trace bounds give $\norm{T_m}_\op \le t$ and $\abs{v^{\top}T_mv} \le t\norm{v}^2$
for $m \le K$ (Proposition~\ref{prop:opnorm_trace}; $T_0 = \Id$ and $t \ge 1$),
so $B = T_1 - s\Id$ has $\norm{B}_\op \le t + 1$. The remainders of
Lemma~\ref{lem:knh}, $R_{m+1} = \Psi_{m+1} + s\Psi_m$, have operator norm at
most $3K\eta t$ for $m < K$ \formal{Column.RemainderRecurrence.pmcr\_recurrence\_opbound}.
Lemma~\ref{lem:cheb} applies.
\end{proof}

\begin{proposition}[Coercivity from the box \formal{SquareFree.Coercivity.paper\_main\_squarefree\_coercivity}, \formal{SquareFree.Coercivity.paper\_main\_squarefree\_coercivity\_scale}]\label{prop:coercive}
Let $n > 0$, $d \ge 9$ and $4n \le d^2$, so that $0 \le s \le 3/4$. If every
row satisfies $\abs{D_i/M - 1} \le 1/64$, where $D_i = \sum_cX_{ic}^2$, and
$-(49/24)\Id \preceq T_1 - s\Id$, then $QQ^{*}$ is coercive with constant
$d^2/128$.
\end{proposition}

\begin{proof}
From $\kappa(\sum_cR_c - \sum_cD_c) \succeq (s - 49/24)\Id$, $\sum_cD_c = \diag(D_i) \succeq (63/64)M\Id$
and $1/\kappa = sM$: $\sum_cR_c \succeq M(63/64 + s^2 - (49/24)s)\Id \succeq (M/64)\Id$,
since $1 + s^2 - (49/24)s - 1/32 = (s - 3/4)(s - 31/24) \ge 0$ on $[0, 3/4]$.
Then $VV^{\top} = 2\sum_cR_c \succeq (M/32)\Id$, $2M = d(d-1) \ge d^2/2$, and
$QQ^{*} \succeq VV^{\top}$ (Lemma~\ref{lem:blocks}). Finally
$s^2 = n/M \le (d^2/4)/(d(d-1)/2) = d/(2(d-1)) \le 9/16$ for $d \ge 9$.
\end{proof}

The three hypotheses of Proposition~\ref{prop:nbbox} hold with probability
tending to one at the scale $4n \le d^2 \le c_1n$.

\begin{proposition}[Column remainders \formal{Column.RemainderProbability.paper\_main\_column\_remainder\_probability}]\label{prop:remainders}
Let $n(d)$ have density at least $c > 0$ eventually and $0 < n(d) \le d^2$
eventually. For every $K$ and $\eta > 0$ there are measurable events of
probability tending to zero outside which, for all large $d$, every $E_m$
with $m < K$ has $\norm{E_m}_\op \le \eta$.
\end{proposition}

\begin{proof}
The remainder splits as
\[
E_k = \kappa^{k+2}\sum_c\bigl((R_c - D_c)^{k+2} - R_c^{k+2}\bigr) + \kappa^{k+2}\sum_c(Q_c^{k+1} - n^{k+1})R_c + s^k\Bigl(\kappa^2n\sum_cD_c - \Id\Bigr),
\]
where $Q_c = \norm{X_c}_2^2$ and $R_c^{k+1} = Q_c^kR_c$
\formal{Column.Remainder.paper\_main\_column\_remainder\_split}. Each term is
controlled by the \emph{frame} conditions
\formal{Column.ConcentrationFrame.paper\_main\_column\_frame\_good}:
(i) every $Q_c \le C_0n$ and $\abs{Q_c/n - 1} \le \delta$ for $a < b$;
(ii) every $X_{ic}^2 \le d$; (iii) $\sum_cR_c \preceq C_0n\Id$;
(iv) every row sum $D_i \le C_0n$ and $\abs{D_i/M - 1} \le \delta$;
with $C_0 = 2 + 109/c$. The first term is a sum of products of matrices
with $D_c$ inserted, bounded by a Cauchy--Schwarz over the family
\formal{Column.RemainderPowers.pmcr\_power\_sum\_error} using the energy
bounds $\sum_c\norm{R_cu}^2 \le (C_0n)^2\norm{u}^2$ and
$\sum_c\norm{D_cu}^2 \le d\,C_0n\norm{u}^2$; the second by
$\abs{Q_c^{k+1} - n^{k+1}} \le n^{k+1}(2^{k+1} - 1)\delta$ and (iii)
\formal{Column.ConcentrationCoefficient.paper\_main\_column\_coefficient\_opbound};
the third is diagonal with entries $\abs{\kappa^2nD_i - 1} = \abs{D_i/M - 1} \le \delta$.
With $\kappa n \le 2$ and $d \le C_0n\delta^2$ the total is at most
$\bigl((k+2)(4C_0)^{k+2} + 4^{k+1}2C_0 + 2^k\bigr)\delta$
\formal{Column.RemainderFrame.paper\_main\_column\_remainder\_frame\_scaled},
which is at most $\eta$ for $k < K$ when $\delta = \delta(K, C_0, \eta)$ is
small \formal{Column.RemainderTolerance.paper\_main\_remainder\_tolerance\_spec}.
The frame conditions hold outside the union of: the column event
$\{\exists a < b : \abs{Q_{ab} - n} > \min(\delta,1)n\}$, of probability at
most $d^2(10176n + 192n(n-1))/(\min(\delta,1)n)^4$ by Markov with the fourth
moment of $Q_{ab} - n$ \formal{Column.Concentration.paper\_main\_column\_event\_bound}
(the entries $(x_{ia}x_{ib})^2 - 1$ have mean $0$, variance $8$ and fourth
moment $10176$ \formal{Column.ConcentrationMoments}); the upper-edge event
of Proposition~\ref{prop:upperedge}, outside which $\sum_cR_c = \frac12VV^{\top} \preceq \frac12\cdot109d^2\Id \preceq (109/c)n\Id$;
the coordinate clipping event, outside which $(x_{ia}x_{ib})^2 \le d$; and
the diagonal event of Proposition~\ref{prop:upperedge} at level
$\delta/4$, which gives $\abs{D_i/M - 1} \le \delta$ since $D_i = \frac12G_{ii}$
and $2M = d(d-1)$. All have probability tending to zero
\formal{Column.ConcentrationFrameEvents.paper\_main\_column\_frame\_event\_limit}.
\end{proof}

\begin{proposition}[Traces of the non-backtracking matrices \formal{SquareFree.WalkNBMoment.pmnct\_moment\_le}, \formal{SquareFree.WalkNBRate.pmnbr\_rate}, \formal{SquareFree.LowerProbabilityTrace.paper\_main\_squarefree\_trace\_tends\_to\_zero}]\label{prop:nbtrace}
Fix $k$ and let $K = 2^{k+1}$, $t = 12(K+1)$ and $p = 2^{\lfloor\log_2\lfloor\log_2 d\rfloor\rfloor + 4} = 2L(d)$.
If eventually $0 < n(d)$, $4n(d) \le d^2$ and $d^2 \le c_1n(d)$, then the
union over $1 \le m \le K$ of the events $\{\tr T_m^p > t^p\}$ is measurable
and its probability tends to zero; for all large $d$ it is at most $1/d$.
\end{proposition}

\begin{proof}
$\E\tr T_{m}^{q}$ is $\kappa^{mq}$ times a sum over closed walks of $q$
segments of $m$ non-backtracking steps, the junction rows being summed, of
the walk moments \formal{SquareFree.WalkNB.pmnb\_trace\_moment}. Marking the
segment starts turns these into walks that are non-backtracking except at
marks \formal{SquareFree.WalkNBMarks}, and a sharp count of the marked
walks by a twin code with a terminal core
\formal{SquareFree.WalkNBCore.pmnc\_XFZ\_le}, \formal{SquareFree.WalkNBCoreSum.pmns\_T\_le}
gives, with $L = mq$ and $B = 4L$, $\E\tr T_m^q \le F(L, q)(L+1)^2(1 + \delta)^L\Theta$,
where $F(L, q) = 2^q\binom{L+q}{q}$, $\delta = 2Ld(4L)^4n\kappa^2$ and
$\Theta = n\mu^L + n(1 + 13c_0^2)^2 + (4L)^4$, $\mu = 1 + \vartheta(1 + K_L)$,
under the deterministic hypotheses $n \le d^2$, $nM \le \lambda^4$,
$n \le \lambda^2$, $d \le c_0\lambda$, $B^2 \le \lambda\vartheta^{294}$ and
$B^24^{45} \le d$, with $\kappa = \lambda^{-2}$. The paper's constants are
not needed: with $\lambda = (nM)^{1/4}$, $c_0 = 4c_1$,
$\vartheta = 1/(CL_{\max}^4)$ and $L_{\max} = pK$, one has
$F2^p \le t^p$, $L\delta \le 1$ and $L\vartheta(1 + K_L) \le 1$, so Markov
gives for each $m \le K$ a bound
$3(L_{\max}+1)^2\bigl(3n + n(1 + 13c_0^2)^2 + (4L_{\max})^4\bigr)/2^p$ on
$\Prob\{\tr T_m^p > t^p\}$ \formal{SquareFree.WalkNBRate.pmnbr\_term}; summed
over $m \le K$ this is at most $1/d$ once $d \ge 2^{9600}$ and
$A(c_1, k)L(d)^{1178} \le d$ \formal{SquareFree.WalkNBRate.pmnbr\_rate},
which holds for all large $d$ since $L(d) \le 8\log_2d$
\formal{SquareFree.LogarithmicGrowth.pmlg\_eventual}.
\end{proof}

\begin{theorem}[Lower edge \formal{SquareFree.LowerProbability.paper\_main\_squarefree\_lower\_probability}, \formal{Fixed.LowerProbability.paper\_main\_fixed\_lower\_probability}; paper Lemma~2.6 and Proposition~2.8, lower part, for $n/d^2 \le 1/4$]\label{thm:loweredge}
Let $n(d)$ have density at least $c > 0$ eventually and satisfy
$0 < n(d)$, $4n(d) \le d^2$ and $d^2 \le c_1n(d)$ eventually. There are
measurable events $E_d$ with $\Prob(E_d) \to 0$ such that for all large $d$
and every sample outside $E_d$, $QQ^{*}$ is coercive with constant
$d^2/128$. In particular this holds for $n(d) = \lfloor\beta d^2\rfloor$,
$0 < \beta < 1/4$.
\end{theorem}

\begin{proof}
Choose the scalar parameters before the data: with $\theta = 49/24 > 2$ and
$A = 12$, there is $k$ such that $2^Kt + g \le \theta^K$ for some $g > 0$,
$K = 2^{k+1}$, $t = 12(K+1)$ \formal{SquareFree.ChebyshevParameters.paper\_main\_chebyshev\_parameters}
(Bernoulli's inequality: $(\theta/2)^K \ge (1 + K u/2)^2$ with
$u = \theta/2 - 1$ beats $12(K+1)$ for large $K$), and then $\eta > 0$ with
$\rho(t+1, 3K\eta t, K) \le g$ \formal{SquareFree.ChebyshevParametersEdge.paper\_main\_chebyshev\_edge\_parameters\_exact}
(the remainder bound is linear in $\eta$). Let $E_d$ be the union of the
remainder event of Proposition~\ref{prop:remainders} at $(K, \eta)$, the row
event $\{\exists i : \abs{D_i/M - 1} > 1/64\}$ (the frame at tolerance
$1/64$, Proposition~\ref{prop:remainders}) and the trace event of
Proposition~\ref{prop:nbtrace}. Outside $E_d$, Propositions~\ref{prop:nbbox}
and~\ref{prop:coercive} give the coercivity. For $n(d) = \lfloor\beta d^2\rfloor$,
the count conditions hold eventually with $c = \beta/2$ and some $c_1$
\formal{Fixed.SubcriticalCounts.paper\_main\_fixed\_subcritical\_counts}.
\end{proof}

\subsection{The homogeneous formulation and the alternative}\label{sec:alternative}

\begin{definition}
The sample \emph{spans} $\R^d$ when $x_i^{\top}v = 0$ for all $i$ implies
$v = 0$ \formal{Geometry.Homogeneous.paper\_main\_spanning}. A matrix $S$
is in the \emph{kernel} of the sample when all the values $x_i^{\top}Sx_i$
are equal \formal{Geometry.Homogeneous.paper\_main\_kernel}.
\end{definition}

\begin{lemma}[Homogeneous formulation \formal{Geometry.Homogeneous.paper\_main\_homogeneous\_kernel}, \formal{Geometry.Homogeneous.paper\_main\_homogeneous\_witness}; paper Lemma~2.12(i), Lemma~4.2]\label{lem:homog}
For $d > 0$ and a spanning sample, $x \in \Fit_{n,d}$ if and only if some
nonzero $S \succeq 0$ is in the kernel, if and only if there are $S \succeq 0$
with $\norm{S}_\Fro = 1$ and $b \in \R$ with $x_i^{\top}Sx_i - \tr S = b$ for
all $i$. Without any hypothesis, a fit yields such a pair $(S, b)$ with
$S \succeq 0$, $\norm{S}_\Fro = 1$ \formal{Projection.PinchingLoss.paper\_main\_fit\_unit\_witness}.
\end{lemma}

\begin{proof}
A fit $S$ is nonzero (as $d > 0$) and in the kernel. Conversely let
$0 \ne S \succeq 0$ have all values equal to $c$. If $c = 0$ then
$x_i^{\top}Sx_i = 0$ forces $Sx_i = 0$ for all $i$ ($S \succeq 0$), so $S = 0$
by spanning \formal{Geometry.Homogeneous.paper\_main\_psd\_spanning\_zero},
a contradiction; so $c > 0$ and $dS/c$ is a fit. The witness form is the
normalization $S/\norm{S}_\Fro$, $b = (c - \tr S)/\norm{S}_\Fro$, and
conversely $c = \tr S + b$.
\end{proof}

\begin{theorem}[The alternative \formal{Geometry.Separation.paper\_main\_homogeneous\_alternative}, \formal{Geometry.Exclusion.paper\_main\_fit\_balanced\_exclusive}, \formal{Geometry.Separation.paper\_main\_fit\_alternative}; paper Lemma~2.12(ii)]\label{thm:alternative}
For every sample, either some nonzero $S \succeq 0$ is in the kernel, or
$x \in \Bal_{n,d}$; for $d > 0$ not both. Hence for $d > 0$ and a spanning
sample, $x \in \Fit_{n,d}$ if and only if $x \notin \Bal_{n,d}$.
\end{theorem}

\begin{proof}
\emph{Exclusion.} If $S \succeq 0$ is in the kernel with common value $c$
and $y$ has $\sum_iy_i = 0$ with $M_y = \sum_iy_ix_ix_i^{\top} \succ 0$, then
$0 < \ip{M_y}{S} = \sum_iy_ix_i^{\top}Sx_i = c\sum_iy_i = 0$ when $S \ne 0$
(a positive definite $M_y$ pairs positively with a nonzero $S \succeq 0$),
a contradiction.

\emph{Existence.} Suppose no nonzero $S \succeq 0$ is in the kernel, and
$d, n > 0$ (if $d = 0$ the balanced event holds trivially; if $n = 0$ the
hypothesis fails since $\Id$ is in the empty kernel). On the compact set
$\{S \succeq 0 : \tr S = 1\}$ \formal{Geometry.Compact}, the continuous
function $F(S) = \sum_i(a_i(S) - \bar a(S))^2$, $a_i(S) = x_i^{\top}Sx_i$,
$\bar a = \frac1n\sum_ia_i$, attains its minimum at some $S_0$; $F(S_0) > 0$,
else $S_0$ would be a nonzero kernel element. Put $y_i = a_i(S_0) - \bar a(S_0)$,
so $\sum_iy_i = 0$. For $S$ in the set and $\tau \in (0, 1]$,
$S_0 + \tau(S - S_0)$ is in the set and
$F(S_0) \le F(S_0 + \tau(S - S_0)) = F(S_0) + \tau\,c + \tau^2a$ with
$c = 2\sum_iy_i(z_i - y_i)$, $z_i = a_i(S) - \bar a(S)$, and $a \ge 0$; letting
$\tau \to 0$ gives $c \ge 0$, that is $\sum_iy_ia_i(S) = \sum_iy_iz_i \ge \sum_iy_i^2 = F(S_0) > 0$
\formal{Geometry.Separation.pms\_separation}. Testing with $S = ww^{\top}/\norm{w}_2^2$
gives $w^{\top}M_yw > 0$ for every $w \ne 0$.

The last claim combines both parts with Lemma~\ref{lem:homog}.
\end{proof}

\begin{lemma}[Almost sure spanning \formal{Gaussian.Spanning.AlmostSure.paper\_main\_gaussian\_spanning\_ae}; paper line 678]\label{lem:spanning}
If $d \le n$, the sample spans $\R^d$ almost surely.
\end{lemma}

\begin{proof}
The determinant of the $d\times d$ matrix of the first $d$ points is a
polynomial in the coordinates that does not vanish identically (it equals
$1$ at the sample whose first $d$ points are the standard basis), and the
zero set of a nonzero polynomial is $\Prob_{n,d}$-null, by induction on the
number of variables through Fubini (a nonzero one-variable polynomial has
finitely many roots) \formal{Gaussian.Spanning.MultivariateZeros.paper\_main\_gaussian\_multivariate\_null}.
A nonzero determinant means the points span.
\end{proof}

\subsection{Measurability of the events}\label{sec:measurability}

The paper does not discuss the measurability of its events. The formal
proof establishes it by countable rational approximation.

\begin{proposition}[\formal{Measurability.Fit.pm\_fit\_event\_measurable}, \formal{Measurability.Conditioned.pm\_conditioned\_event\_measurable}, \formal{Measurability.Balanced.pm\_balanced\_pd\_event\_measurable}]\label{prop:measurable}
For all $n, d$ and reals $\lambda_\pm$, the events $\Fit_{n,d}$,
$\Cond_{n,d}(\lambda_-, \lambda_+)$ and $\Bal_{n,d}$ are measurable.
\end{proposition}

\begin{proof}
Let $\mathrm{res}(S, x) = \sum_{i<n}(x_i^{\top}Sx_i - d)^2$, a polynomial in
the coordinates for fixed $S$, hence measurable
\formal{Measurability.Coordinates.pm\_fit\_residual\_measurable}. For
$m, k \in \N$ let $\mathcal{I}_{m,k}$ be the countable set of symmetric
rational matrices $q$ with entries bounded by $m$ and
$q + \frac{1}{k+1}\Id \succeq 0$ (a condition on $q$ only), and
\[
F_m = \bigcap_k\bigcup_{q\in\mathcal{I}_{m,k}}\{x : \mathrm{res}(q, x) < \tfrac1{k+1}\}.
\]
Each $F_m$ is measurable, and $\Fit_{n,d} = \bigcup_mF_m$. Indeed, a fit $S$
has entries bounded by some $m - 1$, and for each $k$ a symmetric rational
$q$ within $\eta$ of $S$ entrywise, with $\eta$ small against $d$, $k$ and
the sample, lies in $\mathcal{I}_{m,k}$ with $\mathrm{res}(q,x) < 1/(k+1)$
\formal{Measurability.Fit.pm\_fit\_witness}. Conversely, if
$x \in F_m$, the witnesses $q_k \in \mathcal{I}_{m,k}$ form a bounded
sequence in $\R^{d\times d}$, which has a cluster point $p$ (compactness of
a box \formal{Measurability.Box.pm\_box\_cluster}); by continuity of
$S \mapsto \mathrm{res}(S, x)$, of the entries and of the quadratic form,
$p \succeq 0$ and $\mathrm{res}(p, x) = 0$ \formal{Measurability.Fit.pm\_fit\_cluster},
so $p$ is a fit. The conditioned event adds to the index the conditions
$\abs{\tr q - d} \le \frac1{k+1}$ and
$-\frac{\norm{v}^2}{k+1} \le v^{\top}qv - \lambda_-\norm{v}^2$,
$v^{\top}qv - \lambda_+\norm{v}^2 \le \frac{\norm{v}^2}{k+1}$ for all $v$, which pass
to the cluster point. The balanced event is the union over rational weights
$r$ with $\sum r_i = 0$ and over $k$ of the sets
$\bigcap_{w\in\mathbb{Q}^d}\{x : \frac{\norm{w}^2}{k+1} \le w^{\top}M_r(x)w\}$,
using that a positive definite matrix is uniformly positive,
$v^{\top}Mv \ge \norm{v}^2/(k+1)$ for some $k$ (by compactness of the sphere),
that rational $w$ suffice by continuity, and that rational weights with sum
zero approximate real ones (subtract the mean)
\formal{Measurability.Balanced.pm\_rat\_zero\_sum\_approx}.
\end{proof}

Every exceptional event of the proof is of the form $\{F(x) \le c\}$ or
$\{c < F(x)\}$ for a continuous function $F$ of the coordinates, or a finite
or countable union of such; these are measurable
\formal{Analysis.ContinuousTupleMeasurability}.

\subsection{Reduction to fixed densities}\label{sec:fixed}

\begin{lemma}[Deleting points \formal{Gaussian.Sample.Deletion.paper\_main\_drop\_points\_law}, \formal{Gaussian.Sample.Deletion.paper\_main\_fit\_probability\_antitone\_add}, \formal{Gaussian.Sample.Deletion.paper\_main\_conditioned\_probability\_antitone\_add}]\label{lem:drop}
The map that forgets the first $k$ points of a sample of $k + n$ points is
measurable and sends $\Prob_{k+n,d}$ to $\Prob_{n,d}$. A fit, or a
well-conditioned fit, of the larger sample is one of the smaller sample,
so $\Prob_{k+n,d}(\Fit_{k+n,d}) \le \Prob_{n,d}(\Fit_{n,d})$ and likewise for
$\Cond$.
\end{lemma}

\begin{proof}
Forgetting the first coordinate of a tuple sends the law of $(m+1)$-tuples
to that of $m$-tuples by construction (Definition~\ref{def:sample});
iterate $kd$ times. The inclusions are immediate.
\end{proof}

\begin{proposition}[Fixed densities \formal{Fixed.Density.paper\_main\_conditioned\_from\_fixed}, \formal{Fixed.Density.paper\_main\_unsat\_from\_fixed}]\label{prop:fixed}
Let $n_\beta(d) = \lfloor\beta d^2\rfloor$ \formal{Fixed.Density.paper\_main\_fixed\_count}.
\begin{enumerate}
\item[(a)] Suppose that the events $\Cond_{n,d}(\lambda_-,\lambda_+)$ are
measurable and that for every $0 < \beta < 1/4$ there are
$0 < \lambda_- \le \lambda_+$ with $\Prob_{n_\beta(d),d}(\Cond_{n_\beta(d),d}(\lambda_-,\lambda_+)) \to 1$.
Then Theorem~\ref{thm:cond} holds.
\item[(b)] Suppose that the events $\Fit_{n,d}$ are measurable and that
$\Prob_{n_\beta(d),d}(\Fit_{n_\beta(d),d}) \to 0$ for every $1/4 < \beta < 1/2$.
Then Theorem~\ref{thm:main}(b) holds.
\end{enumerate}
\end{proposition}

\begin{proof}
(a) Given $\alpha^* < 1/4$ pick $\beta$ with $0 < \beta$ and
$\alpha^* < \beta < 1/4$, and the constants $\lambda_\pm$ for $\beta$. If
$\limsup n(d)/d^2 \le \alpha^*$ then $n(d)/d^2 \le \beta$ eventually, hence
$n(d) \le n_\beta(d)$ eventually, and Lemma~\ref{lem:drop} gives
$\Prob_{n(d),d}(\Cond) \ge \Prob_{n_\beta(d),d}(\Cond) \to 1$, while
probabilities are at most $1$. (b) If $\liminf n(d)/d^2 > 1/4$, pick
$1/4 < \beta < 1/2$ below the liminf; then $n_\beta(d) \le n(d)$ eventually,
and $\Prob_{n(d),d}(\Fit) \le \Prob_{n_\beta(d),d}(\Fit) \to 0$.
\end{proof}

\section{The satisfiable phase}\label{sec:sat}

\subsection{Statement and strategy}

Throughout this section $0 \le \alpha < 1/4$, $d = m + 1$, and the sample
has $N \le \alpha d^2$ points (eventually). We write $W_i$ for the head
features and $g = 1/2 - \sqrt\alpha > 0$.

\begin{theorem}[Well-conditioned fits from the feature edges \formal{Conditioned.SpectralReduction.paper\_main\_subcritical\_raw\_spectral\_reduction}, \formal{Head.SubcriticalFit.paper\_main\_head\_subcritical\_fit\_exception}; paper Theorem~3.1]\label{thm:sat}
For every $0 \le \alpha < 1/4$ there is $\kappa \ge 1$ with the following
property. Let $C \ge 0$, $c > 0$, and let $N(d) \le \alpha d^2$ eventually.
Let $B_d$ be measurable events with $\Prob(B_d) \to 0$ such that, for all
large $d$, every sample outside $B_d$ has $QQ^{*}$ coercive with constant
$cd^2$ and $QQ^{*} \preceq Cd^2\,\Id$. Then
\[
\Prob_{N(d),d}\bigl(\Cond_{N(d),d}((2\kappa+1)^{-1}, 2\kappa + 1)\bigr) \to 1 .
\]
\end{theorem}

Combined with the feature edges of Section~\ref{sec:edge} at
$N(d) = \lfloor\beta d^2\rfloor$ this gives Theorem~\ref{thm:cond}
(Section~\ref{sec:assembly}).

\emph{Strategy.} The paper proves $0 \in \mathcal{W}(\mathcal{B})$ for a
Loewner box $\mathcal{B}$ by convex separation (its Lemma~3.2). The formal
proof produces the fit from a \emph{kernel element} of a bounded-condition
cone (Proposition~\ref{prop:kernel}): it suffices to show that, for every
unit $y \in \R^N$, some $P$ of condition $\kappa'$ has
$\sum_iy_i\ip{W_i}{P} > 0$. Each unit $y$ is split into a head $y_J$
supported on a set $J$ of $O(d\ln d)$ coordinates and a tail $y_{J^c}$ that
is either tiny or spread (Lemma~\ref{lem:peel}). For every head $J$ of
that size, the section $\cone_{\kappa}\cap\mathbb{S}_\Fro\cap\ker W_J$ has
Gaussian width at least $(1/2 - \varepsilon)d$ with probability tending to
one, uniformly in $J$ (Theorem~\ref{thm:heads}); on such a section, Gordon's
inequality for the Gaussian surrogate rows, a second-chaos interpolation and
a concentration argument give, for every spread tail direction $t$ capped
by $\theta/\sqd$, a section element $P$ with $\sum_{i\notin J}t_i\ip{W_i}{P} > 0$
(Theorem~\ref{thm:lowinf}), simultaneously over all heads. A tiny tail is
absorbed by a perturbation of a section element (Proposition~\ref{prop:support_all}).

\subsection{Head--tail decomposition}

\begin{lemma}[\formal{Head.TailDecompositionLog.paper\_main\_head\_tail\_regularity}; paper Lemma~3.5]\label{lem:peel}
Let $\theta > 0$, $d \ge 1$ and $y \in \R^N$ with $\norm{y}_2 = 1$. There is a
set $J \subseteq \{0, \dots, N-1\}$ with $\abs{J} \le \frac{16}{\theta^2\ln 2}\,d\ln d$
such that, with $r = \norm{y_{J^c}}_2$, either $r \le d^{-2}$, or $r > 0$ and
$\abs{y_i}/r \le \theta/\sqd$ for all $i \notin J$.
\end{lemma}

\begin{proof}
Peel: start with $J = \varnothing$ and $r = 1$. While $r > d^{-2}$, move to
$J$ every $i \notin J$ with $\abs{y_i} > \theta r/(2\sqd)$; at most
$4d/\theta^2$ coordinates are moved in a round (their squares sum to at
most $r^2$). If the new remainder norm $r'$ satisfies $r' \ge r/2$, stop: every
remaining coordinate is at most $\theta r/(2\sqd) \le \theta r'/\sqd$.
Otherwise $r' < r/2$ and the process continues; after $\ell$ rounds the
remainder is at most $2^{-\ell}$, so it falls below $d^{-2}$ within
$\lceil 2\log_2 d\rceil + 1$ rounds. The size bound follows
\formal{Head.Peeling.peeling\_iteration}.
\end{proof}

\subsection{From positive support values to a kernel element}

\begin{proposition}[Kernel elements \formal{Cone.Kernel.paper\_main\_cone\_kernel\_exists\_support}, \formal{Cone.Kernel.paper\_main\_cone\_kernel\_exists}; paper Lemma~3.2, sufficient direction]\label{prop:kernel}
Let $d > 0$, $\kappa \ge 1$ and $W_0, \dots, W_{m-1} \in \R^{d\times d}$. If for
every unit $y \in \R^m$ some $P$ on the cone sphere $\cone_\kappa\cap\mathbb{S}_\Fro$
has $\sum_iy_i\ip{W_i}{P} > 0$, then some $P$ on the cone sphere has
$\ip{W_i}{P} = 0$ for every $i$.
\end{proposition}

\begin{proof}
The trace-one slice $\Sigma = \{S : \lambda\Id \preceq S \preceq \kappa\lambda\Id \text{ for some } \lambda \ge 0,\ \tr S = 1\}$
(the closure of the cone in the Loewner form, cut at trace one) is compact
\formal{Cone.KernelSlice.paper\_main\_cone\_kernel\_slice\_compact}: it is
bounded since $0 \preceq S$ and $\tr S = 1$ give $\norm{S}_\Fro \le 1$, and
closed since $\lambda$ ranges in $[1/(d\kappa), 1/d]$ on it. The continuous
function $F(S) = \sum_i\ip{W_i}{S}^2$ attains its minimum on $\Sigma$ at
some $S_0$. If $F(S_0) > 0$, put $y_i = \ip{W_i}{S_0}$, not all zero. For $S \in \Sigma$
and $\tau \in (0,1]$ the segment point $S_0 + \tau(S - S_0)$ lies in the
convex set $\Sigma$, and $F(S_0) \le F(S_0 + \tau(S - S_0))$ gives, as
$\tau \to 0$, $\sum_iy_i\ip{W_i}{S - S_0} \ge 0$, that is
$\sum_iy_i\ip{W_i}{S} \ge F(S_0) > 0$ for all $S \in \Sigma$. Every $P$ on the
cone sphere is a positive multiple of a point of $\Sigma$, so
$\sum_i(-y_i)\ip{W_i}{P} < 0$ for all such $P$, contradicting the hypothesis
for the unit vector $-y/\norm{y}_2$. Hence $F(S_0) = 0$; $S_0 \ne 0$
(trace one), and $P = S_0/\norm{S_0}_\Fro$ is a kernel element on the cone
sphere (if $\lambda = 0$ at $S_0$ then $S_0 \preceq 0$, impossible with trace
one, so $S_0$ has condition $\kappa$).
\end{proof}

\begin{lemma}[A kernel element is a well-conditioned fit \formal{Conditioned.SupportFit.paper\_main\_cone\_kernel\_conditioned\_fit}; paper Lemma~3.2, from a kernel element to a fit]\label{lem:kernel_fit}
If $d > 0$, $\kappa \ge 1$ and $P$ on $\cone_\kappa\cap\mathbb{S}_\Fro$ has
$\ip{W_i}{P} = 0$ for all $i < N$, then $x \in \Cond_{N,d}(\kappa^{-1}, \kappa)$.
\end{lemma}

\begin{proof}
$\ip{W_i}{P} = 0$ means $x_i^{\top}Px_i = \tr P$. With $\lambda\Id \preceq P \preceq \kappa\lambda\Id$,
$\lambda > 0$, one has $d\lambda \le \tr P \le d\kappa\lambda$; put
$S = dP/\tr P$. Then $x_i^{\top}Sx_i = d$, $\tr S = d$, and
$\kappa^{-1}\Id \preceq (d\lambda/\tr P)\Id \preceq S \preceq (d\kappa\lambda/\tr P)\Id \preceq \kappa\Id$.
\end{proof}

\begin{proposition}[All dual directions \formal{Head.TailSupportAll.paper\_main\_support\_all}, \formal{Head.TailSupportSplit.paper\_main\_support\_spread}, \formal{Head.TailSupportSplit.paper\_main\_support\_tiny}, \formal{Cone.HeadTailFit.paper\_main\_cone\_head\_tail\_conditioned\_fit}]\label{prop:support_all}
Let $d \ge 2$, $\theta > 0$, $\kappa \ge 1$, $\mu > 0$, $A \ge 0$, and let
$\mathcal{J}$ be the family of sets $J$ of distinct indices with
$\abs{J} \le \frac{16}{\theta^2\ln2}d\ln d$. Suppose:
\begin{enumerate}
\item[(i)] the Gram matrix of $(W_i)_{i<N}$ is coercive with constant $\mu$;
\item[(ii)] for every $J \in \mathcal{J}$ there is $P$ of condition $\kappa$ with $\ip{W_i}{P} = 0$ for $i \in J$;
\item[(iii)] $\norm{\sum_ic_iW_i}_\Fro \le A\norm{c}_2$ for every $c \in \R^N$;
\item[(iv)] for every $J \in \mathcal{J}$ and every unit $t \in \R^N$ supported off $J$ with
$\abs{t_i} \le \theta/\sqd$ for all $i$, there is $P$ of condition $\kappa$ with
$\ip{W_i}{P} = 0$ for $i \in J$ and $\sum_it_i\ip{W_i}{P} > 0$;
\item[(v)] $A\,d^{-2}(\kappa + \tfrac12)\sqd < \tfrac12\cdot\tfrac{\sqrt\mu}{2}$.
\end{enumerate}
Then $x \in \Cond_{N,d}((2\kappa+1)^{-1}, 2\kappa+1)$.
\end{proposition}

\begin{proof}
Scale the matrices of (ii) and (iv) so that their lower Loewner bound is
$\Id$: then $\Id \preceq P \preceq \kappa\Id$. Let $y$ be a unit vector and
decompose it by Lemma~\ref{lem:peel} into $y_J$ and $y_{J^c}$,
$r = \norm{y_{J^c}}_2$. Put $a = 1$, $b = \kappa$, $\delta = 1/2$.

\emph{Spread tail.} If $r > 0$ and $t = y_{J^c}/r$ is capped, (iv) gives $P$
with $\ip{W_i}{P} = 0$ on $J$ and $\sum_it_i\ip{W_i}{P} > 0$; then
$\sum_iy_i\ip{W_i}{P} = r\sum_it_i\ip{W_i}{P} > 0$, and $P$ lies in the box
$[a - \delta, b + \delta]$.

\emph{Tiny tail.} If $r \le d^{-2}$, take $S_0$ from (ii), with $\Id \preceq S_0 \preceq \kappa\Id$,
and let $M_u = \sum_{i\in J}y_iW_i$, $M_v = \sum_{i\notin J}y_iW_i$. Then
$\ip{M_u}{S_0} = 0$, $\norm{M_u}_\Fro^2 \ge \mu\norm{y_J}_2^2 \ge \mu(1 - r^2) \ge \mu/4$
by (i), and $\norm{M_v}_\Fro \le Ar$ by (iii). The perturbation
$S = S_0 + \delta M_u/\norm{M_u}_\Fro$ satisfies
$(a - \delta)\Id \preceq S \preceq (b + \delta)\Id$ (a unit Frobenius matrix
has operator norm at most one) and
\[
\ip{M_u + M_v}{S} \ge \delta\norm{M_u}_\Fro - \norm{M_v}_\Fro\norm{S}_\Fro
\ge \delta\frac{\sqrt\mu}{2} - A\,d^{-2}(b + \delta)\sqd > 0
\]
by (v), using $\norm{S}_\Fro \le (b+\delta)\sqd$
\formal{Head.TailSupport.paper\_main\_support\_perturbation}. Thus in both
cases some $S$ in the box $[1/2, \kappa + 1/2]$, that is of condition
$\kappa' = 2\kappa + 1$, has $\sum_iy_i\ip{W_i}{S} > 0$; normalizing $S$ to
unit Frobenius norm, Proposition~\ref{prop:kernel} gives a kernel element on
$\cone_{\kappa'}\cap\mathbb{S}_\Fro$, and Lemma~\ref{lem:kernel_fit} concludes.
\end{proof}

The rest of the section establishes (ii) and (iv) with probability tending
to one, with $\mu = cd$ and $A = \sqrt{C}\sqd$ (Lemma~\ref{lem:bridge_features}),
for which (v) holds for large $d$.

\subsection{Section widths}\label{sec:widths}

Fix $\kappa_0 \ge 1$ and let $P(z)$ be a measurable choice of an
approximate maximizer of the support value on the cone sphere.

\begin{lemma}[Approximate maximizers \formal{Cone.ApproxSelector.paper\_main\_cone\_selector\_gaussian}; replaces the measurable maximum theorem of the paper, line 964]\label{lem:selector}
For $\kappa_0 \ge 1$ and $\varepsilon > 0$ there is a map $z \mapsto P(z)$ from
$\R^{\Dd_d}$ to $\Sym^d$ with measurable entries such that for every $z$,
$P(z)$ lies on $\cone_{\kappa_0}\cap\mathbb{S}_\Fro$ and
$h_{\kappa_0}(Z) - \varepsilon\norm{Z}_\Fro \le \ip{Z}{P(z)} \le h_{\kappa_0}(Z)$,
where $h_{\kappa_0} = h_{\cone_{\kappa_0}\cap\mathbb{S}_\Fro}$.
\end{lemma}

\begin{proof}
Take a finite $\varepsilon$-net $P_0, \dots, P_{s-1}$ of the cone sphere in
Frobenius norm (Theorem~\ref{thm:nets}, applied to the codes of the
matrices) and let $P(z)$ be the first net point maximizing
$\ip{Z}{P_t}$ \formal{Head.FiniteSupportSelector.paper\_main\_finite\_argmax\_measurable}.
For every $S$ on the cone sphere there is $t$ with $\norm{S - P_t}_\Fro \le \varepsilon$,
so $\ip{Z}{S} \le \ip{Z}{P_t} + \varepsilon\norm{Z}_\Fro$.
\end{proof}

\begin{lemma}[Residuals of the maximizer \formal{Clipping.ResidualConditional.paper\_main\_cone\_selector\_conditional}, \formal{Cone.ApproxSelectorResidual.paper\_main\_cone\_selector\_residual}]\label{lem:residual}
Let $P$ be as in Lemma~\ref{lem:selector} and $r, \tau > 0$. The joint event
$\{(z, x) : \exists i < N,\ \abs{\ip{W_i}{P(z)}} \ge r\}$ is measurable with
$\gauss_{\Dd_d}\otimes\Prob_{N,d}$-probability at most $2N\exp(-dr^2/(4 + 4\kappa_0r))$,
and the event of samples $\{x : \gauss_{\Dd_d}\{z : \exists i,\ \abs{\ip{W_i}{P(z)}} \ge r\} \ge \tau\}$
is measurable with probability at most $\tau^{-1}\cdot 2N\exp(-dr^2/(4 + 4\kappa_0r))$.
\end{lemma}

\begin{proof}
For fixed $z$, $P(z)$ is a fixed element of the cone sphere, so by
Lemma~\ref{lem:hw} and a union bound over the rows (whose law is $\gauss_d$),
the section of the joint event has probability at most
$2N\exp(-dr^2/(4 + 4\kappa_0r))$; Fubini gives the joint bound. The
second event is the set where the section probability is at least $\tau$;
by Markov's inequality its probability is at most $\tau^{-1}$ times the
joint probability \formal{Gaussian.Sample.ConditionalException.paper\_main\_conditional\_exception\_reverse\_markov}.
\end{proof}

We take $r = \sqrt{4 + 4\kappa_0}\,d^{-1/4}$
\formal{Clipping.ResidualRates.paper\_main\_residual\_threshold}, for which
$dr^2/(4 + 4\kappa_0r) \ge \sqd/2$ once $d \ge (4 + 4\kappa_0)^2$, and
$\tau = d^{-4}$; the second event of Lemma~\ref{lem:residual} then has
probability at most $2Nd^4e^{-\sqd/2}$, which tends to zero for $N \le Cd^2$.

\begin{definition}[Regular heads \formal{Head.WidthRegular.paper\_main\_head\_regular}]
Let $K, Q, L \ge 0$ and $c, c_0 > 0$. A selection $s$ of $m$ distinct rows of the sample is
\emph{regular} when $m \le Kd\ln d$, the Gram matrix of all the $W_i$ is
coercive with constant $cd$, $\ip{W_{s(k)}}{W_{s(k)}} \ge c_0d$ for every $k$,
$\abs{\ip{W_{s(k)}}{W_{s(l)}}} \le Q\ln d$ for $k \ne l$, and the selected points
satisfy $X_JX_J^{\top} \preceq Ld(\ln d)^2\,\Id$.
\end{definition}

\begin{proposition}[Deterministic correction \formal{Head.WidthActual.paper\_main\_head\_width\_actual\_selected}, \formal{Geometry.HeadSpectral.paper\_main\_head\_correction\_bounds\_cone}, \formal{Geometry.HeadCorrection.paper\_main\_head\_condition\_transfer}; paper Lemma~3.8]\label{prop:correction}
Let $s$ select $m$ distinct rows, let the Gram matrix $G$ of all the $W_i$
be coercive with constant $\mu > 0$, let $\ip{W_{s(k)}}{W_{s(k)}} \ge q_0 > 0$
and $\abs{\ip{W_{s(k)}}{W_{s(l)}}} \le q_1$ for $k \ne l$, and
$X_JX_J^{\top} \preceq \lambda\Id$. Let $P \in \cone_{\kappa_0}\cap\mathbb{S}_\Fro$
with residuals $r_k = \ip{W_{s(k)}}{P}$, $\abs{r_k} \le r$. Let
$\eta, A_1, A_\infty, \delta \ge 0$ satisfy $mr^2 \le \mu\eta^2$, $\eta < 1$,
$mr \le \mu A_1$, $r + q_1A_1 \le q_0A_\infty$, $(A_\infty\lambda + A_1)/\sqd \le \delta$,
$\delta < (\kappa_0\sqd)^{-1}$ and $\kappa_0/\sqd + \delta \le \kappa((\kappa_0\sqd)^{-1} - \delta)$.
Then there is $\tilde P$ in the section $\cone_\kappa\cap\mathbb{S}_\Fro\cap\ker W_J$
with $\ip{Z}{\tilde P} \ge \ip{Z}{P} - 2\eta\norm{Z}_\Fro$ for every $Z$.
\end{proposition}

\begin{proof}
Let $G_J$ be the Gram matrix of the selected features, coercive with
constant $\mu$ as a principal submatrix of $G$. Solve $G_Ja = r_J$ and put
$E = -\sum_ka_kW_{s(k)}$, so that $\ip{W_{s(k)}}{P + E} = 0$ for all $k$
\formal{Geometry.HeadCorrection.paper\_main\_head\_correction\_kernel}.
Bounds: $\norm{E}_\Fro^2 = a^{\top}G_Ja = a^{\top}r_J \le \norm{a}_2\norm{r_J}_2$
and $\mu\norm{a}_2^2 \le a^{\top}G_Ja$ give $\norm{a}_2 \le \norm{r_J}_2/\mu \le \sqrt m\,r/\mu$
and $\norm{E}_\Fro^2 \le mr^2/\mu \le \eta^2$; $\norm{a}_1 \le \sqrt m\norm{a}_2 \le mr/\mu \le A_1$;
from $(G_J)_{kk}a_k = r_k - \sum_{l\ne k}(G_J)_{kl}a_l$,
$\norm{a}_\infty \le (r + q_1\norm{a}_1)/q_0 \le A_\infty$. Since
$E = -d^{-1/2}\bigl(X_J\diag(a)X_J^{\top} - (\sum_ka_k)\Id\bigr)$,
$\norm{E}_\op \le (A_\infty\lambda + A_1)/\sqd \le \delta$, so
$-\delta\Id \preceq E \preceq \delta\Id$. By Lemma~\ref{lem:box},
$(\kappa_0\sqd)^{-1}\Id \preceq P \preceq (\kappa_0/\sqd)\Id$, hence
$((\kappa_0\sqd)^{-1} - \delta)\Id \preceq P + E \preceq (\kappa_0/\sqd + \delta)\Id$,
which is the condition $\kappa$ by the last hypothesis; in particular
$P + E \ne 0$. Let $\tilde P = (P + E)/\norm{P + E}_\Fro$. From
$\norm{E}_\Fro \le \eta < 1$ and $\norm{P}_\Fro = 1$ one gets
$\norm{\tilde P - P}_\Fro \le 2\eta$ \formal{Geometry.HeadCorrection.paper\_main\_head\_normalize\_close},
whence the support transfer.
\end{proof}

\begin{lemma}[Rates \formal{Head.CorrectionScale.paper\_main\_head\_correction\_eventual}; paper Lemma~3.8, choice of rates]\label{lem:rates}
Let $K, a \ge 0$, $c, c_0 > 0$, $Q, L \ge 0$, $\kappa_0 \ge 1$. There are
$E > 0$, $A, B \ge 0$, $F > 0$ and $D$ such that for every $d \ge D$, with
$t = d^{1/32}$, the choices $\eta = E/t^7$, $A_1 = A/t^7$, $A_\infty = B/(dt^6)$,
$\delta = F/(\sqd\,t^4)$ satisfy the hypotheses of
Proposition~\ref{prop:correction} with $\mu = cd$, $q_0 = c_0d$, $q_1 \le Q\ln d$,
$\lambda \le Ld(\ln d)^2$, $m \le Kd\ln d$, $r = a/t^8$, and $\kappa = 2\kappa_0^2$.
\end{lemma}

\begin{proof}
Each hypothesis is a comparison of powers of $d$ with powers of $\ln d$,
which holds for large $d$: for instance $mr^2 \le Kd\ln d\cdot a^2t^{-16} \le cd\,E^2t^{-14}$
once $K a^2\ln d \le cE^2t^2$; the constants $E, A, B, F$ are chosen in this
order \formal{Head.CorrectionRates.paper\_main\_head\_rate\_constants}. The
residual threshold $r = \sqrt{4 + 4\kappa_0}\,d^{-1/4} = a/t^8$ with
$a = \sqrt{4 + 4\kappa_0}$. The last two hypotheses hold because
$\delta\sqd = Ft^{-4} \to 0$ while $\kappa_0/\sqd + \delta \le 2\kappa_0^2((\kappa_0\sqd)^{-1} - \delta)$
reads $\kappa_0 + \delta\sqd \le 2\kappa_0 - 2\kappa_0^2\delta\sqd$.
\end{proof}

\begin{proposition}[Width transfer \formal{Head.WidthRegular.paper\_main\_head\_width\_regular\_transfer}, \formal{Head.WidthExpectation.paper\_main\_head\_width\_expectation\_transfer}; paper Proposition~3.3, transfer inequality]\label{prop:width_transfer}
Let $\varepsilon > 0$ and $\kappa_0 \ge 1$, and let $P$ be the selector of
Lemma~\ref{lem:selector} at precision $\varepsilon/2$. For all large $d$,
every sample outside the event of Lemma~\ref{lem:residual} (at $r$, $\tau = d^{-4}$)
and every regular selection $s$ of $m$ rows satisfy
\[
w_{\kappa_0} - w\bigl(\cone_{2\kappa_0^2}\cap\mathbb{S}_\Fro\cap\ker W_J\bigr) \le \varepsilon d .
\]
\end{proposition}

\begin{proof}
Let $A_x = \{z : \exists k,\ \abs{\ip{W_{s(k)}}{P(z)}} \ge r\}$; outside the
event of Lemma~\ref{lem:residual}, $\gauss_{\Dd_d}(A_x) < d^{-4}$. For
$z \notin A_x$, Proposition~\ref{prop:correction} with the rates of
Lemma~\ref{lem:rates} gives a section element $\tilde P$ with
$h_{\mathrm{sec}}(Z) \ge \ip{Z}{\tilde P} \ge \ip{Z}{P(z)} - 2\eta\norm{Z}_\Fro \ge h_{\kappa_0}(Z) - (\varepsilon/2 + 2\eta)\norm{Z}_\Fro$.
For every $z$, $\abs{h_{\mathrm{sec}}(Z)} \le \norm{Z}_\Fro \le \norm{Z}_\Fro^2/(2a) + a/2$
for any $a > 0$, and likewise for $h_{\kappa_0}$. Integrating, with
$\E\norm{Z}_\Fro \le \sqrt{\Dd_d} \le d$, $\E\norm{Z}_\Fro^2 = \Dd_d$ and $a = d^2$,
\[
w_{\kappa_0} - w_{\mathrm{sec}} \le (\varepsilon/2 + 2\eta)d + \frac{\Dd_d}{d^2} + d^2\gauss_{\Dd_d}(A_x) \le (\varepsilon/2 + 2\eta)d + 1 + d^{-2},
\]
which is at most $\varepsilon d$ for large $d$ since $\eta \to 0$.
\end{proof}

\begin{theorem}[Uniform section widths \formal{Head.RegularProbability.paper\_main\_head\_width\_coercive\_probability\_uniform}; paper Proposition~3.3]\label{thm:heads}
For every $\varepsilon > 0$ there is $\kappa \ge 1$ such that for all
$C, K \ge 0$ and $c > 0$ and every $N(d) \le Cd^2$ eventually, there are
measurable events $A_d$ with $\Prob(A_d) \to 0$ such that, for all large
$d$, every sample outside $A_d$ whose feature Gram matrix is coercive with
constant $cd$ satisfies: for every selection $s$ of $m \le Kd\ln d$ distinct
rows, the section $\cone_\kappa\cap\mathbb{S}_\Fro\cap\ker W_J$ has width at
least $(1/2 - \varepsilon)d$.
\end{theorem}

\begin{proof}
By Proposition~\ref{prop:conewidth} choose $\kappa_0 \ge 1$ with
$w_{\kappa_0} \ge (1/2 - \varepsilon/2)d$ for all large $d$, and put
$\kappa = 2\kappa_0^2$. Let $A_d$ be the union of the event of
Lemma~\ref{lem:residual}, of the sparse-norm event of Lemma~\ref{lem:sparse}
(with $L = L(C, K)$), of the radius event and of the coherence event
(Lemmas~\ref{lem:radius}, \ref{lem:coherence}); all have probability
tending to zero. Outside $A_d$, a selection of $m \le Kd\ln d$ rows is
regular with $c_0 = c$ (the diagonal entries of a coercive matrix are at
least its constant), $Q = 73$ and the $L$ of Lemma~\ref{lem:sparse}, and
Proposition~\ref{prop:width_transfer} at $\varepsilon/2$ gives
$w_{\mathrm{sec}} \ge w_{\kappa_0} - \varepsilon d/2 \ge (1/2 - \varepsilon)d$.
\end{proof}

The paper's constant is $\kappa_1 = 2\kappa_0$; the formal proof works with
Loewner boxes and obtains $2\kappa_0^2$.

\subsection{The low-influence tail margin}\label{sec:lowinf}

Fix a head $J$ of $k$ rows and write $h$ for the head sample (the rows of
$J$) and $x$ for the tail sample of $q$ rows; the \emph{section} is
$\mathcal{D} = \cone_\kappa\cap\mathbb{S}_\Fro\cap\ker W_J$, assumed
nonempty, with features $V_i = W_i$, $i \notin J$. For a set $Y$ of unit
vectors of $\R^q$ the \emph{section minimax} is
\[
\inf_{y\in Y}\ \sup_{S\in\mathcal{D}}\ \sum_{i<q}y_i\ip{V_i}{S}
\tag*{\formal{Head.SectionMinimax.paper\_main\_head\_section\_minimax}.}
\]
The \emph{capped sphere} is the set of unit $y$ with $\abs{y_i} \le \theta/\sqd$
for all $i$.

\begin{definition}[Soft minimax \formal{Soft.Definitions.soft\_minimax}; paper Lemma~3.12, sandwich bounds]\label{def:softmm}
For finite index lists $Y, S$ and $H : Y \times S \to \R$, with $\beta, \gamma > 0$,
\[
\Phi_{\beta,\gamma}(H) = -\frac1\beta\ln\sum_{y\in Y}\exp\Bigl(-\beta\cdot\frac1\gamma\ln\sum_{s\in S}e^{\gamma H(y,s)}\Bigr).
\]
It satisfies $\min_y\max_sH - \frac{\ln\abs{Y}}{\beta} \le \Phi_{\beta,\gamma}(H) \le \min_y\max_sH + \frac{\ln\abs{S}}{\gamma}$
\formal{Soft.Extrema.soft\_minimax\_bounds}.
\end{definition}

\begin{lemma}[Nets of the section and of the capped sphere \formal{Head.SectionApproxWidth.paper\_main\_head\_section\_approx\_width\_exists}, \formal{Gordon.Entropy.paper\_main\_gordon\_entropy\_net}, \formal{Gordon.Nets.paper\_main\_gordon\_net\_width\_error}]\label{lem:section_nets}
Let $\kappa \ge 1$, $\varepsilon > 0$, $Y$ a nonempty set of unit vectors of
$\R^q$ and $\mathcal{D} \ne \varnothing$. There are unit vectors $y^0, \dots, y^{r-1}$
with $r \le e^{C(\varepsilon)q}$ forming an $\varepsilon$-net of $Y$ and
section elements $B_0, \dots, B_{s-1} \in \mathcal{D}$ with $s \le e^{C(\varepsilon)\Dd_d}$
forming an $\varepsilon$-net of $\mathcal{D}$, where $C(\varepsilon) = (4 + 4/\varepsilon)\ln 2$,
such that:
\begin{enumerate}
\item[(a)] $w(\{B_t\}) \le w(\mathcal{D}) \le w(\{B_t\}) + \varepsilon\sqrt{\Dd_d}$, where
$w(\{B_t\}) = \E\max_t\ip{Z}{B_t}$;
\item[(b)] for every family $V$ with $\norm{\sum_iu_iV_i}_\Fro \le L\norm{u}_2$ and every $\beta, \gamma > 0$,
\[
\Bigl|\ \inf_{y\in Y}\sup_{S\in\mathcal{D}}\sum_iy_i\ip{V_i}{S} - \Phi_{\beta,\gamma}\bigl(\textstyle\sum_iy^a_i\ip{V_i}{B_t}\bigr)\ \Bigr|
\le \frac{C(\varepsilon)\Dd_d}{\gamma} + \frac{C(\varepsilon)q}{\beta} + L\varepsilon .
\]
\end{enumerate}
\end{lemma}

\begin{proof}
The nets come from Theorem~\ref{thm:nets} applied to the codes of the
section elements in the orthonormal basis of $\Sym^d$ (unit vectors of
$\R^{\Dd_d}$) and to $Y$; the library's grid construction
\formal{Gordon.Entropy} gives nets \emph{inside} the sets with the stated
entropy. (a): $\max_t\ip{Z}{B_t} \le h_{\mathcal{D}}(Z)$ since the $B_t$ are
in $\mathcal{D}$, and every $S \in \mathcal{D}$ has a net point with
$\ip{Z}{S - B_t} \le \varepsilon\norm{Z}_\Fro$, whose mean is at most
$\varepsilon\sqrt{\Dd_d}$. (b): replacing $y$ by a net point changes
$\sum_iy_i\ip{V_i}{S}$ by at most $\norm{\sum_i(y - y^a)_iV_i}_\Fro\norm{S}_\Fro \le L\varepsilon$,
and replacing $S$ by a net point changes it by at most
$\norm{\sum_iy_iV_i}_\Fro\varepsilon \le L\varepsilon$; the soft minimax is
within the entropy terms of the finite minimax.
\end{proof}

\begin{proposition}[Gordon's margin for the Gaussian surrogate \formal{Gordon.Width.paper\_main\_gordon\_finite\_width\_margin}, \formal{Gordon.SoftMargin.paper\_main\_gordon\_soft\_width\_margin}, \formal{Soft.FeatureMargin.paper\_main\_public\_goe\_soft\_margin}; replaces paper Lemma~2.3 and Lemma~3.13]\label{prop:gordon_margin}
Let $z^0, \dots, z^{r-1}$ be unit vectors of $\R^N$ and $v^0, \dots, v^{s-1}$ unit
vectors of $\R^M$, and let $\Gamma \in \R^{N\times M}$ have independent
$N(0,1)$ entries. Then, with $w = \E\max_t\ip{h}{v^t}$ for $h \sim \gauss_M$,
\[
\E\min_a\max_t (z^a)^{\top}\Gamma v^t \ge w - \sqrt N,
\qquad
\E\,\Phi_{\beta,\gamma}\bigl((z^a)^{\top}\Gamma v^t\bigr) \ge w - \sqrt N - \frac{\ln r}{\beta}.
\]
Consequently, for the Gaussian feature rows $G_i$, $i < N$, unit vectors
$y^a \in \R^N$ and symmetric $B_t$ with $\norm{B_t}_\Fro = 1$,
\[
\E\,\Phi_{\beta,\gamma}\Bigl(\sum_iy^a_i\ip{G_i}{B_t}\Bigr) \ge \sqrt{\tfrac2d}\,\bigl(w(\{B_t\}) - \sqrt N\bigr) - \frac{\ln r}{\beta}.
\]
\end{proposition}

\begin{proof}
Corollary~\ref{cor:bil} gives
$\E\min_a\max_t(z^a)^{\top}\Gamma v^t \ge \E\min_a\max_t(\ip{g}{z^a} + \ip{h}{v^t}) = \E\max_t\ip{h}{v^t} + \E\min_a\ip{g}{z^a} \ge w - \E\norm{g}_2 \ge w - \sqrt N$
\formal{Gordon.Processes.paper\_main\_gordon\_process\_comparison}, with
$\E\norm{g}_2 \le \sqrt N$ by Jensen. The soft minimax is at least the
minimax minus $\ln r/\beta$. For the feature rows,
$\ip{G_i}{B_t} = \sqrt{2/d}\,\ip{\Gamma_i}{v^t}$ where $v^t$ is the unit code
of $B_t$ and $\Gamma_i$ the row of Gaussian coordinates, so the array
$\sum_iy^a_i\ip{G_i}{B_t}$ is $\sqrt{2/d}$ times the bilinear process, and
the soft minimax is positively homogeneous after rescaling $\beta, \gamma$;
the formal statement is the scaled inequality.
\end{proof}

The paper obtains a high-probability margin from Gordon's tail inequality
and passes to expectations with a second-moment bound; the formal proof
uses the expectation form directly.

\begin{theorem}[Second-chaos interpolation for the soft minimax \formal{Soft.Comparison.paper\_main\_soft\_public\_comparison}, \formal{Soft.PublicFeatures.paper\_main\_public\_feature\_soft\_comparison\_coarse}, \formal{Multirow.IntegratedBound.paper\_multirow\_integrated\_bound\_results}; paper Lemmas~3.9, 3.10, 3.11 and 3.12 and Corollary~3.14, in the special case of the soft minimax]\label{thm:interp_soft}
Let $y^a$ ($a \in Y$) be unit vectors of $\R^N$ with $\abs{y^a_i} \le \theta/\sqd$,
and $B_t$ ($t \in S$) symmetric with $\norm{B_t}_\op \le k$. Let
$\Phi(A) = \Phi_{\beta,\gamma}(\sum_iy^a_i\ip{A_i}{B_t})$ for a family
$A = (A_i)_{i<N}$ of symmetric matrices. Then
\[
\abs{\E\Phi(W) - \E\Phi(G)} \le \frac{40}{3}\,k^3(\beta + \gamma)^2\theta\,\frac{d + 2\sqrt{d\ln N} + 2\ln N}{d^2},
\]
and if $\ln N \le d$ this is at most $\frac{200}{3}k^3(\beta + \gamma)^2\theta/d$.
\end{theorem}

\begin{proof}
\emph{Interpolation.} For $t \in [0,1]$ let $Z_i(t) = \sqrt t\,W_i + \sqrt{1-t}\,G_i$,
with $W_i$ built from the point $x_i \sim \gauss_d$ and $G_i$ from
independent Gaussian coordinates, and $f(t) = \E\Phi(Z(t))$; $f$ is
continuous on $[0,1]$ \formal{Multirow.ExpectationAssembly.mexp\_clause\_endpoints}
and $f(1) = \E\Phi(W)$, $f(0) = \E\Phi(G)$.

\emph{Derivative.} For $0 < t < 1$, by the chain rule and the two Stein
identities, with $\Psi_i = \nabla_i\Phi(Z(t))$ (the derivative of $\Phi$ in
its $i$-th matrix argument, a symmetric matrix),
\[
f'(t) = \sum_i\Bigl(\frac{1}{2\sqrt t}\E\ip{W_i}{\Psi_i} - \frac{1}{2\sqrt{1-t}}\E\ip{G_i}{\Psi_i}\Bigr).
\]
Gaussian integration by parts in the coordinates of $G_i$ (covariance
$(2/d)\Id$ in the orthonormal basis) gives $\frac{1}{2\sqrt{1-t}}\E\ip{G_i}{\Psi_i} = \frac1d\E\operatorname{div}_i\Psi_i$,
where $\operatorname{div}_i\Psi = \sum_\nu\ip{E_\nu}{D_i\Psi[E_\nu]}$. The
second-order Stein identity in the coordinates of $x_i$ applied to the
entries of $\Psi_i$ (with $\ip{W_i}{\Psi} = d^{-1/2}\sum_{a,b}(x_{ia}x_{ib} - \delta_{ab})(\Psi)_{ab}$
and $\partial_a\partial_b$ of $\Psi(Z_i(t))$ through
$D_xZ_i[e_a] = \sqrt{t/d}\,(e_ax_i^{\top} + x_ie_a^{\top})$ and
$D_x^2Z_i[e_a,e_b] = \sqrt{t/d}\,(e_ae_b^{\top} + e_be_a^{\top})$) gives
\[
\frac{1}{2\sqrt t}\E\ip{W_i}{\Psi_i} = \frac1d\E\operatorname{div}_i\Psi_i + \frac{\sqrt t}{2d^{3/2}}\E\mathcal{R}_i(t),
\quad
\mathcal{R}_i(t) = \frac12\sum_{a,b}D_i^3\Phi[U_a, U_b, V_{ab}],
\]
with $U_a = e_ax_i^{\top} + x_ie_a^{\top}$ and $V_{ab} = e_ae_b^{\top} + e_be_a^{\top}$
\formal{Multirow.FullStein.paper\_multirow\_full\_stein\_results}. The
divergence terms cancel and
$f'(t) = \frac{\sqrt t}{2d^{3/2}}\sum_i\E\mathcal{R}_i(t)$.

\emph{The remainder.} For a trilinear form given as $\sum_\ell c_\ell\prod_{j\le3}\tr(B_{\ell j}H_j)$
the contraction $\frac12\sum_{a,b}T[U_a,U_b,V_{ab}]$ equals
$4\sum_\ell c_\ell\,x_i^{\top}B_{\ell1}B_{\ell3}B_{\ell2}x_i$, so
$\abs{\mathcal{R}_i(t)} \le 4\norm{x_i}_2^2\,\norm{D_i^3\Phi}_{\pi,\op}$, the
projective operator norm being the infimum of $\sum_\ell\abs{c_\ell}\prod_j\norm{B_{\ell j}}_\op$
over such representations \formal{Soft.MatrixProjectiveBound}. For the soft
minimax, the third derivative is a sum of three terms (the inner third
cumulant, a covariance of the inner Hessian with the gradient, and the
outer third cumulant), each a combination of pure tensors
$\tr(y_iB_s\cdot)\tr(y_iB_{s'}\cdot)\tr(y_iB_{s''}\cdot)$ with Gibbs
weights, and the computation of its projective norm gives
\[
\sum_i\norm{D_i^3\Phi}_{\pi,\op} \le \frac{10\,k^3(\beta + \gamma)^2\theta}{\sqd}
\tag*{\formal{Soft.MatrixProjectiveBound.soft\_matrix\_projective\_sum\_bound},}
\]
using $\sum_i\abs{y_i}^3 \le \theta d^{-1/2}\sum_iy_i^2 = \theta/\sqd$ and
$\norm{B_s}_\op \le k$. Hence
$\abs{f'(t)} \le \frac{2\sqrt t}{d^{3/2}}\,\E\max_i\norm{x_i}_2^2\cdot\frac{10k^3(\beta+\gamma)^2\theta}{\sqd}$.

\emph{The radius.} $\E\max_{i<N}\norm{x_i}_2^2 \le d + 2\sqrt{d\ln N} + 2\ln N$
\formal{Radius.ExpectedMaximum.gaussian\_maximum\_radius\_expected\_bound}:
for $0 < \lambda < 1/2$, $\ln\E e^{\lambda(\norm{x}_2^2 - d)} = \frac d2(-2\lambda - \ln(1 - 2\lambda)) \le \frac{d\lambda^2}{1 - 2\lambda}$,
so by the log-sum-exp bound and Jensen, $\E(\max_i\norm{x_i}^2 - d) \le \frac{\ln N}{\lambda} + \frac{d\lambda}{1-2\lambda}$,
and $\lambda = \sqrt{\ln N}/(\sqd + 2\sqrt{\ln N})$ gives the claim (for
$N = 1$ let $\lambda \to 0$).

\emph{Integration.} $\int_0^12\sqrt t\,dt = 4/3$ gives
$\abs{f(1) - f(0)} \le \frac{4}{3}\cdot\frac{R_{N,d}}{d^{3/2}}\cdot\frac{10k^3(\beta+\gamma)^2\theta}{\sqd}$
\formal{Analysis.UnitIntervalComparison}, which is the first bound; the
second uses $R_{N,d} \le 5d$ when $\ln N \le d$.
\end{proof}

\begin{corollary}[Expected margin for the chaos rows \formal{Soft.ConeMargin.paper\_main\_data\_soft\_cone\_margin}, \formal{Head.SoftTail.paper\_main\_head\_soft\_expectation}; paper Corollary~3.14]\label{cor:chaos_margin}
In the setting of Lemma~\ref{lem:section_nets} with $\kappa \ge 1$, unit
$y^a$ capped by $\theta/\sqd$ and section elements $B_t$ (so
$\norm{B_t}_\op \le \kappa/\sqd$ by Lemma~\ref{lem:box}), for $\ln N \le d$,
\[
\E\,\Phi_{\beta,\gamma}\Bigl(\sum_iy^a_i\ip{W_i}{B_t}\Bigr)
\ge \sqrt{\tfrac2d}\bigl(w(\{B_t\}) - \sqrt N\bigr) - \frac{\ln r}{\beta} - \frac{200}{3}\,\frac{\kappa^3(\beta+\gamma)^2\theta}{d^2\sqd}.
\]
\end{corollary}

\begin{proof}
Proposition~\ref{prop:gordon_margin} and Theorem~\ref{thm:interp_soft} with
$k = \kappa/\sqd$.
\end{proof}

\begin{lemma}[Parameters \formal{Soft.ParameterChoice.paper\_main\_soft\_parameters}, \formal{Soft.ParameterChoice.paper\_main\_soft\_scaled\_margin}]\label{lem:params}
Let $g = 1/2 - \sqrt\alpha$, $\kappa \ge 1$ and $L \ge 0$. There are
$\varepsilon, A, \theta > 0$ with $\varepsilon \le g/16$, $L\varepsilon \le g/16$,
$C(\varepsilon)(\alpha + 1)/A \le g/16$ and $\frac{800}{3}\kappa^3A^2\theta \le g/16$.
With these, $\beta = \gamma = Ad\sqd$, $q \le \alpha d^2$, $\Dd_d \le d^2$,
$\ln r \le C(\varepsilon)q$, $\ln s \le C(\varepsilon)\Dd_d$ and a net width
$w \ge (1/2 - g/16)d$, the level of Corollary~\ref{cor:chaos_margin} minus
the entropy and net errors of Lemma~\ref{lem:section_nets}(b) is at least
$\frac g2\sqd$:
\[
\sqrt{\tfrac2d}\bigl(w - \varepsilon\sqrt{\Dd_d} - \sqrt q\bigr) - \frac{\ln r + \ln s}{Ad\sqd} - \frac{200}{3}\frac{\kappa^3(2Ad\sqd)^2\theta}{d^2\sqd} - L\sqd\,\varepsilon \ge \frac g2\sqd .
\]
\end{lemma}

\begin{proof}
$\sqrt{2/d}\,(w - \varepsilon\sqrt{\Dd_d} - \sqrt q) \ge \sqrt2\,(1/2 - g/16 - g/16 - \sqrt\alpha)\sqd = \sqrt2\,\frac{7g}{8}\sqd$;
the entropy terms are at most $C(\varepsilon)(\alpha + 1)d^2/(Ad\sqd) \le \frac{g}{16}\sqd$;
the interpolation term is $\frac{800}{3}\kappa^3A^2\theta\sqd \le \frac g{16}\sqd$;
and $L\varepsilon\sqd \le \frac g{16}\sqd$. Since $\sqrt2\cdot\frac78 - \frac3{16} > \frac12$,
the claim follows.
\end{proof}

\begin{lemma}[Clipping and concentration \formal{Data.ClippedScore.paper\_main\_data\_contraction\_score\_lipschitz}, \formal{Data.SoftConcentration.paper\_main\_data\_soft\_clipped\_lower\_margin}, \formal{Data.SoftScale.paper\_main\_data\_soft\_scaled\_clipping\_large}; paper Lemma~3.15, with coordinate clipping]\label{lem:clip}
Let $U > 0$ and let $\overline W_i$ be the feature of the point whose
coordinates are clipped to $[-U, U]$. For unit $y$ and symmetric $B$ with
$\norm{B}_\op \le k$, the score $x \mapsto \sum_iy_i\ip{\overline W_i}{B}$ is
$2kU$-Lipschitz for the Euclidean norm of all the $Nd$ coordinates, and so
is the soft minimax of such scores. If $b \le \E\Phi(W)$, then for $u > 0$
\[
\Prob\bigl\{\Phi(\overline W) \le b - \ell - u\bigr\} \le \exp\Bigl(-\frac{u^2}{2(2kU)^2}\Bigr),
\qquad \ell = \frac{V}{a} + a\cdot2Nd\,e^{-U^2/2},
\]
for every $a > 0$, where $V$ bounds the second moments of $\Phi(W)$ and
$\Phi(\overline W)$. At $U = d^{1/4}$, $k = \kappa/\sqd$, $\beta = \gamma = Ad\sqd$,
$a = d^6$ and $N \le Cd^2$, $\ell \le \eta'\sqd$ for every $\eta' > 0$ and
all large $d$, and the exponent at $u = c\sqd$ is $c^2d\sqd/(8\kappa^2)$.
\end{lemma}

\begin{proof}
For clipped rows $u, u'$ with $\norm{u}_2, \norm{u'}_2 \le U\sqd$,
$\ip{\overline W_i(u) - \overline W_i(u')}{B} = d^{-1/2}(u^{\top}Bu - u'^{\top}Bu')$
and $\abs{u^{\top}Bu - u'^{\top}Bu'} = \abs{(u - u')^{\top}B(u + u')} \le k\norm{u - u'}_2\cdot2U\sqd$;
Cauchy--Schwarz over the rows with $\norm{y}_2 \le 1$ gives the Lipschitz
constant $2kU$ with no factor $\sqrt N$ (the clipping is $1$-Lipschitz).
A soft minimax of $L$-Lipschitz functions is $L$-Lipschitz. The two
functionals agree outside the coordinate event at level $U$; on it,
$\abs{\Phi(W) - \Phi(\overline W)}$ is bounded through the second moments
(both scores grow at most like $k(\sum x_{ia}^2 + Nd)/\sqd$ plus the entropy
terms, whose squares have Gaussian means $O(d^4)$), and the elementary
inequality $\E\abs{F}\one_E \le \E F^2/a + a\Prob(E)$ gives
$\E\Phi(W) - \E\Phi(\overline W) \le \ell$. Theorem~\ref{thm:conc}
(through Proposition~\ref{prop:bridge}) gives the tail. The scale claims
are arithmetic: $V \le Md^4$, $2Nde^{-\sqd/2} \le 2C\cdot14^7/\sqd$, so
$\ell \le M/(2d^2) + 2C\,38^{19}/\sqd$ \formal{Clipping.Rates.paper\_main\_clipping\_mean\_error\_rate}.
\end{proof}

\begin{proposition}[The bad event of a fixed head \formal{Head.SubcriticalTail.paper\_main\_head\_subcritical\_clipped\_tail}, \formal{Head.CompactTail.paper\_main\_head\_compact\_clipped\_tail}; paper Lemma~3.15, one head]\label{prop:fixed_head}
Let $0 \le \alpha < 1/4$, $\kappa \ge 1$, $L \ge 0$ and $g = 1/2 - \sqrt\alpha$.
There are $\varepsilon, A, \theta, c > 0$ and $D$ such that for all $d \ge D$,
all $m, q$ with $q + 1 \le \alpha d^2$, every head sample $h$ of $m$ rows whose
section $\mathcal{D}$ has width at least $(1/2 - g/16)d$, and every
nonempty set $Y$ of unit vectors of $\R^{q+1}$ capped by $\theta/\sqd$: the
event of tail samples $x$ of $q + 1$ rows on which the clipped features
satisfy $\norm{\sum_iu_i\overline W_i}_\Fro \le L\sqd\norm{u}_2$ and the
section minimax of the clipped features over $Y$ is at most $\frac g4\sqd$,
has probability at most $e^{-cd\sqd}$.
\end{proposition}

\begin{proof}
Take the nets of Lemma~\ref{lem:section_nets} at precision $\varepsilon$ and
the parameters of Lemma~\ref{lem:params}. By Lemma~\ref{lem:section_nets}(a)
the net width is at least $w(\mathcal{D}) - \varepsilon\sqrt{\Dd_d} \ge (1/2 - g/16)d - \varepsilon d$,
and Corollary~\ref{cor:chaos_margin} with Lemma~\ref{lem:params} gives
$\E\Phi(W) \ge \frac g2\sqd + (\text{net and entropy errors})$, where the
errors are those that Lemma~\ref{lem:section_nets}(b) charges to pass from
$\Phi$ to the section minimax. By Lemma~\ref{lem:clip} with $u = \frac g8\sqd$
and $\ell \le \frac g{16}\sqd$ for large $d$, outside an event of probability
$\exp(-g^2d\sqd/(512\kappa^2))$ the clipped soft minimax is at least
$\frac g2\sqd - \frac g{16}\sqd - \frac g8\sqd + \text{errors}$, hence by
Lemma~\ref{lem:section_nets}(b) (applied to the clipped features, which have
the operator bound $L\sqd$ on the event) the section minimax of the clipped
features over $Y$ exceeds $\frac g4\sqd$.
\end{proof}

\begin{proposition}[All heads \formal{Head.SubcriticalFamily.paper\_main\_head\_subcritical\_family\_event}, \formal{Head.SubcriticalProbability.paper\_main\_head\_subcritical\_probability}, \formal{Head.SelectionLaw.paper\_main\_head\_selection\_law}, \formal{Head.FamilyCount.paper\_main\_head\_families\_exp}; paper Proposition~3.4 and Lemma~3.15, union over heads]\label{prop:all_heads}
Let $0 \le \alpha < 1/4$, $\kappa \ge 1$, $L \ge 0$. There is $\theta > 0$
such that for every $K \ge 0$ and every $N(d) \le \alpha d^2$ eventually there
are measurable events $E_d$ with $\Prob(E_d) \to 0$ with the following
property for all large $d$. Let $x$ be a sample outside $E_d$ and outside the
coordinate clipping event at $d^{1/4}$, with $\norm{\sum_iu_iW_i}_\Fro \le L\sqd\norm{u}_2$.
Then every set $J$ of at most $Kd\ln d$ distinct rows whose section has
width at least $(1/2 - g/16)d$ satisfies: some $P$ of condition $\kappa$ has
$\ip{W_i}{P} = 0$ on $J$, and for every unit $t$ supported off $J$ and
capped by $\theta/\sqd$, some such $P$ has $\sum_it_i\ip{W_i}{P} > 0$.
\end{proposition}

\begin{proof}
\emph{One head.} Fix $J$ with $k$ rows. There is a permutation $\pi$ of the
rows whose first $k$ images are $J$ in order; the permuted sample has the
same law, and the pair (head rows, tail rows) has the product law
$\Prob_{k,d}\otimes\Prob_{N-k,d}$ \formal{Head.SelectionLaw.paper\_main\_head\_selection\_law}
(permutation invariance of the product law of the coordinates,
Proposition~\ref{prop:gaussE}). Let $E_J$ be the event that the head $h$
has section width at least $(1/2 - g/16)d$ and the tail $x$ lies in the
bad event of Proposition~\ref{prop:fixed_head} for $Y$ the capped sphere of
$\R^{N-k}$ (nonempty when $N - k \ge d/\theta^2$). By Fubini on the product,
$\Prob(E_J) \le e^{-cd\sqd}$ \formal{Head.LawBounds.paper\_main\_head\_exception\_bound};
measurability holds since the section minimax of continuous features is a
measurable function of the pair \formal{Head.SectionMinimax.paper\_main\_head\_section\_minimax\_measurable}.
Outside $E_J$ and outside the coordinate event (so that $\overline W_i = W_i$),
the section minimax exceeds $\frac g4\sqd > 0$: for every capped unit $t$
the support value $\sup_{S\in\mathcal{D}}\sum_it_i\ip{W_i}{S}$ is positive,
so some $S \in \mathcal{D}$ has a positive score; and $\mathcal{D} \ne \varnothing$
since its width is finite and positive \formal{Head.SubcriticalGood.paper\_main\_head\_subcritical\_selected\_good}.

\emph{The union.} The number of sets $J$ of at most $Kd\ln d$ distinct rows
among $N \le \alpha d^2$ is at most $\exp(K(\ln(\alpha + 1) + 2)d(\ln d)^2)$
\formal{Head.FamilyCount.paper\_main\_head\_families\_entropy}, so
$E_d = \bigcup_JE_J$ has probability at most
$\exp(K'd(\ln d)^2 - cd\sqd) \to 0$ \formal{Clipping.Limits.paper\_main\_clipped\_head\_union\_tends\_to\_zero}.
\end{proof}

\begin{theorem}[Low-influence tail margin; paper Proposition~3.4]\label{thm:lowinf}
Proposition~\ref{prop:all_heads} is the formal form of the paper's
Proposition~3.4: the margin $c_*d$ on $\sqd\,\mathcal{D}$ corresponds to
the margin $\frac g4\sqd$ on $\mathcal{D}$ that the proof establishes, and
the conclusion retained is the positivity of the support value.
\end{theorem}

\subsection{Assembly of the satisfiable phase}

\begin{proof}[Proof of Theorem~\ref{thm:sat}]
Let $0 \le \alpha < 1/4$ and $g = 1/2 - \sqrt\alpha$. Apply
Theorem~\ref{thm:heads} with $\varepsilon = g/16$ to get $\kappa \ge 1$, and
Proposition~\ref{prop:all_heads} with this $\kappa$ and $L = \sqrt C$ to get
$\theta$; let $K = \frac{16}{\theta^2\ln2}$ be the head budget of
Lemma~\ref{lem:peel}. Let $E_d$ be the union of $B_d$, of the events $A_d$
of Theorem~\ref{thm:heads} (with the given $C, K, c$), of the events of
Proposition~\ref{prop:all_heads} (with this $K$), and of the coordinate
clipping event; $\Prob(E_d) \to 0$ \formal{Head.SubcriticalFit.paper\_main\_head\_subcritical\_fit\_exception}.
Outside $E_d$, for large $d$: the head Gram matrix is coercive with
constant $cd$ and $\norm{\sum_iu_iW_i}_\Fro \le \sqrt C\sqd\norm{u}_2$
(Lemma~\ref{lem:bridge_features}); every head of at most $Kd\ln d$ rows has
section width at least $(1/2 - g/16)d$; hence hypotheses (ii) and (iv) of
Proposition~\ref{prop:support_all} hold with $\mu = cd$ and $A = \sqrt C\sqd$,
and (v) holds for large $d$ since $\sqrt C\,d^{-2}(\kappa + \frac12)d \to 0$
while $\sqrt{cd}/4 \to \infty$. Proposition~\ref{prop:support_all} gives
$x \in \Cond_{N,d}((2\kappa+1)^{-1}, 2\kappa+1)$, and the probability of the
conditioned event, which is measurable (Proposition~\ref{prop:measurable}),
is at least $1 - \Prob(E_d) \to 1$
\formal{Conditioned.SpectralReduction.paper\_main\_event\_from\_vanishing\_exception}.
\end{proof}

\section{The unsatisfiable phase}\label{sec:unsat}

\subsection{Statement and strategy}

\begin{theorem}[Unsatisfiable phase at fixed density \formal{Unsat.Limit.paper\_main\_unsat\_fixed\_density\_limit}; paper Theorem~4.1]\label{thm:unsat}
For every $1/4 < \gamma < 1/2$, with $n_\gamma(d) = \lfloor\gamma d^2\rfloor$,
$\Prob_{n_\gamma(d),d}(\Fit_{n_\gamma(d),d}) \to 0$.
\end{theorem}

With Proposition~\ref{prop:fixed}(b) this is Theorem~\ref{thm:main}(b).

\emph{Strategy.} A fit yields a unit witness $S \succeq 0$ with
$x_i^{\top}Sx_i - \tr S = b$ for all $i$ (Lemma~\ref{lem:homog}), and $\abs{b}$
is bounded when the feature Gram matrix is (Lemma~\ref{lem:offset}). The
witness is split at a \emph{constant} spectral cutoff $\varepsilon_0$ into a
head of rank at most $1/\varepsilon_0^2$ and a bulk of operator norm at most
$\varepsilon_0$; the paper uses the cutoff $d^{-1/4}$ and heads of rank up to
$\sqrt d$. The orientation of the head is rounded to a finite net of
orthogonal matrices measured by the head projector, and the witness is
\emph{pinched} by the net projector, which keeps the quadratic risk small
(Section~\ref{sec:witness}); the paper rotates instead. For a fixed
orientation, the admissible candidates are approximated by a finite
\emph{prior} of block pairs, and the quantity studied is the entropy-shifted
soft minimum $\mathcal{Z}$ of the normalized empirical risks over the prior
(Section~\ref{sec:prior}). Under the hybrid law, where the bulk coordinates
are replaced by GOE rows, every candidate has risk at least $e$ outside an
event of probability $e^{-cd^2}$ (Section~\ref{sec:hybrid}); the expectation of
$\mathcal{Z}$ under the data law is close to its expectation under the
hybrid law, by a Lindeberg replacement with a replica expansion
(Section~\ref{sec:universality}); and $\mathcal{Z}$ concentrates around its
mean at scale $N^{-1/2}$ (Section~\ref{sec:concentration}). A union over the
net of orientations and over the ranks concludes (Section~\ref{sec:unsat_assembly}).

\begin{definition}[Risk \formal{Feature.LossControl.paper\_main\_quadratic\_risk}, \formal{Loss.Translated.paper\_main\_loss}]
$\varphi(t) = 1 - e^{-t^2}$, which satisfies $e^{-1}\min(t^2, 1) \le \varphi(t) \le \min(t^2, 1)$
and is $4$-Lipschitz \formal{Loss.Translated.paper\_main\_loss\_sandwich},
\formal{Loss.Translated.paper\_main\_loss\_lipschitz}. The \emph{quadratic
risk} of $(S, b)$ on the sample is
$\mathcal{L}(S, b) = \frac1N\sum_{i<N}\varphi(x_i^{\top}Sx_i - \tr S - b)$;
it vanishes exactly when $(S, b)$ is a witness.
\end{definition}

\begin{lemma}[Lipschitz control of the risk \formal{Feature.LossControl.paper\_main\_quadratic\_risk\_lipschitz}, \formal{Feature.LossControl.paper\_main\_feature\_witness\_offset}; paper (4.5), (4.22)]\label{lem:offset}
Suppose $QQ^{*} \preceq K\Id$ with $K \le NL^2$. Then
$\abs{\mathcal{L}(S, b) - \mathcal{L}(T, b)} \le 4L\norm{S - T}_\Fro$ for all
$S, T$, and $\abs{\mathcal{L}(S,b) - \mathcal{L}(S,b')} \le 4\abs{b - b'}$. If
$\norm{S}_\Fro = 1$ and $x_i^{\top}Sx_i - \tr S = b$ for all $i$, then $\abs{b} \le L$.
\end{lemma}

\begin{proof}
$\sum_i(\ip{Q_i}{S - T})^2 = (S - T)^{\top}QQ^{*}(S - T) \le K\norm{S - T}_\Fro^2 \le N(L\norm{S - T}_\Fro)^2$
(reading $S - T$ in Frobenius coordinates), and $\varphi$ is $4$-Lipschitz;
Cauchy--Schwarz gives the first bound \formal{Loss.Empirical.paper\_main\_loss\_empirical\_lipschitz}.
For the witness, $Nb^2 = \sum_i\ip{Q_i}{S}^2 \le K \le NL^2$.
\end{proof}

\subsection{Witnesses, orientations and pinching}\label{sec:witness}

\begin{proposition}[Ordered spectral split \formal{Orientation.Spectral.paper\_main\_orientation\_spectral\_decomposition}; paper (4.8), (4.9)]\label{prop:split}
Let $S \succeq 0$ with $\norm{S}_\Fro = 1$ and $\varepsilon_0 > 0$. There are
$r \le d$, an orthogonal $O$ and $a \in \R^d$ with $S = O^{\top}\diag(a)O$,
$a_i > \varepsilon_0$ exactly for $i < r$, $r \le 1/\varepsilon_0^2$, such that
with $P = O^{\top}\diag(\one_{i<r})O$ (the \emph{head projector}
\formal{Orientation.Net.paper\_main\_orientation\_projector}), $H = PSP$
and $B = S - H$: $H, B \succeq 0$, $\norm{B}_\op \le \varepsilon_0$,
$\norm{H}_\Fro \le 1$, $PHP = H$, $PS = SP$, $P^{\top} = P = P^2$.
\end{proposition}

\begin{proof}
Theorem~\ref{thm:spectral} gives $O$ and $a \ge 0$ with $\sum_ia_i^2 = 1$;
permute the rows of $O$ and the entries of $a$ so that the entries above
$\varepsilon_0$ come first \formal{Orientation.Spectral.paper\_main\_orientation\_spectral\_order}.
Their number $r$ satisfies $r\varepsilon_0^2 < \sum_ia_i^2 = 1$
\formal{Spectral.Cutoff.paper\_main\_spectral\_head\_cardinality}. The
claims follow from $H = O^{\top}\diag(a\one_{i<r})O$ and
$B = O^{\top}\diag(a\one_{i\ge r})O$.
\end{proof}

\begin{proposition}[Nets of orientations \formal{Orientation.Net.paper\_main\_orientation\_net}, \formal{Orientation.Family.paper\_main\_orientation\_family\_uniform}; replaces the Grassmannian nets of the paper, line 1755]\label{prop:onets}
For $r \le d$ and $\tau > 0$ there are orthogonal matrices $O_0, \dots, O_{s-1}$,
$\ln s \le (4 + 8r/\tau)\ln 2\,(r + 1)(d + 1)$, such that every orthogonal
$O$ has a $k$ with $\norm{P_r(O) - P_r(O_k)}_\Fro \le \tau$. For $r \le r_{\max}$
these families can be chosen with a common entropy bound
$\ln s_r \le (4 + 8r_{\max}/\tau)\ln2\,(r_{\max} + 1)(d + 1)$.
\end{proposition}

\begin{proof}
Let $W$ be the matrix of the first $r$ rows of $O$ (zero below), so
$P_r(O) = W^{\top}W$ and $\norm{W}_\Fro^2 = r$. For two orthogonal $O, O'$,
$\norm{P_r(O) - P_r(O')}_\Fro \le 2\sqrt r\norm{W - W'}_\Fro$
\formal{Projection.Frobenius.paper\_main\_frobenius\_product}. The codes
$W/\sqrt r$ are unit vectors of $\R^{rd}$; a net of the image of the
orthogonal matrices in the unit ball at mesh $\tau/(2r)$
\formal{Loss.BallNet.paper\_main\_loss\_ball\_net}, with a chosen orthogonal
representative for each net point, has $\ln(\text{size}) \le (4 + 8r/\tau)\ln2\,(rd + 1)$
by Theorem~\ref{thm:nets}, which is at most the stated bound. For $r = 0$
the single matrix $\Id$ works.
\end{proof}

\begin{definition}[Pinching and admissible candidates \formal{Projection.Pinching.paper\_main\_pinching}, \formal{Orientation.Witness.paper\_main\_orientation\_admissible}]
For a projector $U$ and a matrix $S$, the \emph{pinching} is
$\Pi_U(S) = USU + (\Id - U)S(\Id - U)$. Given $r$, an orthogonal $O$ with
$U = P_r(O)$, $\varepsilon \ge 0$ and $L \ge 0$, the pair $(A, b)$ is
\emph{admissible} when $A \succeq 0$, $\Pi_U(A) = A$,
$1/4 \le \norm{A}_\Fro^2 \le 4$, $\norm{(\Id - U)A(\Id - U)}_\op \le \varepsilon$
and $\abs{b} \le L$.
\end{definition}

\begin{proposition}[From a fit to an admissible candidate \formal{Orientation.Witness.paper\_main\_orientation\_fit\_witness}, \formal{Projection.PinchingLoss.paper\_main\_pinching\_witness\_risk}, \formal{Projection.PinchingAnnulus.paper\_main\_pinching\_annulus}, \formal{Projection.PinchingOperator.paper\_main\_pinching\_bulk\_opnorm}; replaces paper Lemma~4.6]\label{prop:witness}
Let $\varepsilon_0 > 0$, $0 \le \tau \le 1/8$, $\varepsilon \ge \varepsilon_0 + \tau^2$,
$L, K \ge 0$ and $1/\varepsilon_0^2 \le r_{\max} < d$. Let, for each
$r \le r_{\max}$, $(O_{r,k})_{k<s_r}$ be a $\tau$-net of orientations as in
Proposition~\ref{prop:onets}. Let $x \in \Fit_{N,d}$ with $N > 0$ and
$QQ^{*} \preceq K\Id$, $K \le NL^2$. Then there are $r \le r_{\max}$, $k < s_r$
and an admissible $(A, b)$ for $O_{r,k}$ with $\mathcal{L}(A, b) \le 16L\tau$.
\end{proposition}

\begin{proof}
Lemma~\ref{lem:homog} gives a unit witness $(S, b)$, $\abs{b} \le L$ by
Lemma~\ref{lem:offset}. Split $S$ by Proposition~\ref{prop:split} at
$\varepsilon_0$: $r \le 1/\varepsilon_0^2 \le r_{\max}$, $P = P_r(O)$. Choose
$k$ with $\norm{U - P}_\Fro \le \tau$, $U = P_r(O_{r,k})$, and let
$A = \Pi_U(S)$. Then: $A \succeq 0$ as a sum of two congruences of $S$;
$\Pi_U(A) = A$ since $U(\Id - U) = 0$; $\norm{A}_\Fro^2 = \norm{USU}_\Fro^2 + \norm{(\Id-U)S(\Id-U)}_\Fro^2 \in [1/4, 4]$
when $\norm{U - P}_\Fro \le 1/8$, because $PS = SP$ makes $\Pi_P(S) = S$ and
the pinching is $4\tau$-close to $S$: $\norm{A - S}_\Fro \le 4\tau \le 1/2$
and $\norm{S}_\Fro = 1$; the bulk block
$(\Id - U)A(\Id - U) = (\Id - U)(H + B)(\Id - U)$ has operator norm at most
$\norm{U - P}_\Fro^2\norm{H}_\Fro + \varepsilon_0 \le \tau^2 + \varepsilon_0 \le \varepsilon$,
since $(\Id - U)H(\Id - U) = (\Id - U)(P - U)H(P - U)(\Id - U)$ for $PHP = H$
\formal{Projection.PinchingOperator.paper\_main\_compression\_leakage}. Finally
$\mathcal{L}(S, b) = 0$ and Lemma~\ref{lem:offset} give
$\mathcal{L}(A, b) \le 4L\norm{A - S}_\Fro \le 16L\tau$.
\end{proof}

The bound $\norm{A - S}_\Fro \le 4\tau$ used twice above: with $PS = SP$,
$A - S = -(US(\Id - U) + (\Id - U)SU)$ and
$US(\Id - U) = (U - P)S(\Id - U) + PS(\Id - U) = (U - P)S(\Id - U) + SP(P - U)$,
whose Frobenius norm is at most $2\tau\norm{S}_\Fro$ (each factor $U - P$,
$P - U$ contributes $\tau$ and projectors have operator norm at most one);
the other term is its transpose \formal{Projection.Pinching.paper\_main\_pinching\_perturbation}.

\subsection{Finite priors and the soft minimum}\label{sec:prior}

Fix an orientation $O$ of $\R^d$, $d = r + m$, $m = m_i + 1$, with head
projector $U = P_r(O)$. For a sample $x$, the \emph{rotated sample}
$Ox = (Ox_i)_{i<N}$ has the same law \formal{Orientation.Rotation.paper\_main\_orientation\_rotate\_law}
(the rotation is a linear image with orthonormal rows, Proposition~\ref{prop:mulmx});
its rows split into head coordinates $g_i \in \R^r$ and bulk coordinates
$y_i \in \R^m$. An admissible $A$ is block diagonal in this orientation,
$O A O^{\top} = \diag(H, B)$ with $H \in \Spsd^r$, $B \in \Spsd^m$,
$1/4 \le \norm{H}_\Fro^2 + \norm{B}_\Fro^2 \le 4$, $\norm{B}_\op \le \varepsilon$
\formal{Orientation.Witness.paper\_main\_orientation\_admissible\_blocks},
and $x_i^{\top}Ax_i - \tr A = g_i^{\top}Hg_i - \tr H + y_i^{\top}By_i - \tr B$
\formal{Orientation.Rotation.paper\_main\_orientation\_head\_value}.

\begin{definition}[The finite-prior functionals \formal{Hybrid.FiniteArray.paper\_main\_hybrid\_finite\_data}, \formal{Hybrid.FiniteArray.paper\_main\_hybrid\_finite\_goe}, \formal{Loss.NormalizedArray.paper\_main\_loss\_normalized\_array}]
Let $I = \{0, \dots, s-1\}$ index a finite \emph{prior} of states
$(H_j, b_j, B_j)$, and let $\beta > 0$. The \emph{data row loss} of state
$j$ on the row $(g, y)$ is $\ell^{\mathrm{dat}}_j(g, y) = \varphi(g^{\top}H_jg - \tr H_j - b_j + y^{\top}B_jy - \tr B_j)$;
the \emph{hybrid row loss} replaces the bulk term by the linear form
$\sqrt m\ip{G}{B_j} = \sqrt2\ip{z}{\mathrm{code}(B_j)}$ of a GOE coordinate
vector $z \in \R^{\Dd_m}$, which has the same mean and covariance as the
quadratic term \formal{Projection.Public.paper\_main\_projection\_goe\_covariance}.
For a sample with rows $(g_i, y_i)$, resp. $(g_i, z_i)$, the \emph{array}
of state $j$ is $\mathcal{A}_j = \sum_{i<N}\ell_j(g_i, \cdot)$, and
\[
\mathcal{Z} = \frac1N\Bigl(\min\nolimits_\beta^{j\in I}\mathcal{A}_j + \frac{\ln s}{\beta}\Bigr),
\qquad \min\nolimits_\beta^{j}f_j = -\frac1\beta\ln\sum_{j\in I}e^{-\beta f_j},
\]
written $\mathcal{Z}_{\mathrm{dat}}$ and $\mathcal{Z}_{\mathrm{goe}}$.
\end{definition}

\begin{lemma}[Properties of $\mathcal{Z}$ \formal{Loss.NormalizedArray.paper\_main\_loss\_normalized\_array\_bounds}, \formal{Loss.NormalizedArray.paper\_main\_loss\_normalized\_array\_gap}, \formal{Loss.NormalizedArray.paper\_main\_loss\_normalized\_array\_witness}]\label{lem:Z}
$0 \le \mathcal{Z} \le 1$; if $\mathcal{A}_j \ge Ne$ for every $j \in I$ then
$\mathcal{Z} \ge e$; and $\mathcal{Z} \le \mathcal{A}_j/N + \ln s/(\beta N)$ for
every $j \in I$. $\mathcal{Z}$ is a continuous function of the coordinates
of the sample, and replacing one row changes it by at most $1/N$
\formal{Loss.Gibbs.paper\_main\_gibbs\_row\_change}.
\end{lemma}

\begin{proof}
The soft minimum satisfies $\min_jf_j - \ln s/\beta \le \min_\beta f \le \min_jf_j$;
the row losses take values in $[0,1]$, so a row change moves each
$\mathcal{A}_j$ by at most one and the soft minimum by at most one.
\end{proof}

\begin{proposition}[A finite prior for an orientation \formal{Loss.QuadraticPrior.paper\_main\_quadratic\_prior\_positive}, \formal{Orientation.Prior.paper\_main\_orientation\_prior}, \formal{Orientation.BlocksLoss.paper\_main\_orientation\_candidate\_bound}; replaces paper Lemma~4.9]\label{prop:prior}
Let $O$ be orthogonal, $L \ge 0$, $\delta > 0$, and suppose some admissible
candidate exists. There are $s > 0$ and states $(H_j, b_j, B_j)_{j<s}$ with
$H_j \in \Spsd^r$, $B_j \in \Spsd^m$, $1/4 \le \norm{H_j}_\Fro^2 + \norm{B_j}_\Fro^2 \le 4$,
$\norm{B_j}_\op \le \varepsilon$, $\abs{b_j} \le L$ and
$\ln s \le 3C(\delta/(L+3))\,d^2$, such that for every sample of $N > 0$ points
with $QQ^{*} \preceq K\Id$, $K \le NL^2$: if some admissible $(A, b)$ has
$\mathcal{L}(A, b) \le a$, then $\mathcal{Z}_{\mathrm{dat}}(Ox) \le a + 4(L+1)\delta + \ln s/(\beta N)$.
\end{proposition}

\begin{proof}
The admissible candidates form a subset of $\{(A, b) : A \text{ symmetric},\ \norm{A}_\Fro^2 \le 4,\ \abs{b} \le L\}$,
a ball of radius $\sqrt{8 + L^2} \le L + 3$ in the metric
$\sqrt{2\norm{A - A'}_\Fro^2 + (b - b')^2}$ on $\R^{\Dd_d + 1}$; take a
$\delta$-net of the admissible set inside itself, of size
$\ln s \le C(\delta/(L+3))(\Dd_d + 2) \le 3C(\delta/(L+3))d^2$
\formal{Loss.JointNet.paper\_main\_loss\_joint\_net}. By Lemma~\ref{lem:offset},
every admissible $(A, b)$ has a net point $(A_j, b_j)$ with
$\abs{\mathcal{L}(A, b) - \mathcal{L}(A_j, b_j)} \le 4L\norm{A - A_j}_\Fro + 4\abs{b - b_j} \le 4(L+1)\delta$.
The blocks of $A_j$ in the orientation are the states. Finally, for the
rotated sample the array of state $j$ is $N\mathcal{L}(A_j, b_j)$, and
Lemma~\ref{lem:Z} gives $\mathcal{Z}_{\mathrm{dat}}(Ox) \le \mathcal{L}(A_j, b_j) + \ln s/(\beta N)$.
\end{proof}

\subsection{The hybrid model}\label{sec:hybrid}

In this section the head sample $h = (g_i)_{i<N}$, $g_i \in \R^r$, and the
bulk array $z = (z_i)_{i<N}$, $z_i \in \R^{\Dd_m}$, are independent standard
Gaussian, $m = m_i + 1$, and the \emph{hybrid residual} of $(H, b, B)$ is
$u = \mathsf{K}(H, b) + \mathsf{R}(B) \in \R^N$ with
$\mathsf{K}(H,b)_i = g_i^{\top}Hg_i - \tr H - b$ and
$\mathsf{R}(B)_i = \sqrt2\ip{z_i}{\mathrm{code}(B)}$
\formal{Hybrid.Residual.paper\_main\_hybrid\_residual}; its risk is
$\frac1N\sum_i\varphi(u_i)$ \formal{Hybrid.Loss.paper\_main\_hybrid\_risk}.

\begin{theorem}[Hybrid risk gap \formal{Hybrid.Assembly.paper\_main\_hybrid\_gap\_event}; paper Proposition~4.4, strengthened]\label{thm:hybrid}
For every $\gamma > 1/4$ there are $\eta, e, c > 0$ and $d_0$ such that for
all $d \ge d_0$, $N$ with $\abs{N - \gamma d^2} \le \eta d^2$, $r \le \sqd + 1$
and $r + m_i + 1 \le d$, there is a measurable event $E$ of
$\Prob_{N,r}\otimes\gauss_{N\Dd_m}$-probability at most $e^{-cd^2}$ outside
which, for every symmetric $H$, every $B \succeq 0$ and every $b$ with
$\norm{H}_\Fro^2 + \norm{B}_\Fro^2 \ge 1/4$, the hybrid risk of $(H, b, B)$ is
at least $e$.
\end{theorem}

The paper requires $H \succeq 0$ and $\norm{H}_\Fro^2 + \norm{B}_\Fro^2 \le 4$;
neither is needed. The proof has three probabilistic inputs and one
deterministic step.

\begin{proposition}[Uniform head small ball \formal{Head.SmallBallUniform.paper\_main\_head\_small\_ball\_uniform\_poly}; paper Lemma~4.10]\label{prop:smallball}
Let $r = n + 1$ and $306432(r+1)^3 \le N$. There is a measurable event
$\mathcal{B}$ of head samples with $\Prob_{N,r}(\mathcal{B}) \le \exp(-\frac{2\ln2 - 1}{2432}N)$
such that outside $\mathcal{B}$, for every symmetric $H \in \R^{r\times r}$,
every $b$ and every $J \subseteq \{0,\dots,N-1\}$ with $\abs{J} \ge (1 - 1/1216)N$,
\[
\sum_{i\in J}\bigl(g_i^{\top}Hg_i - \tr H - b\bigr)^2 \ge \frac{N}{4864}\bigl(2\norm{H}_\Fro^2 + b^2\bigr).
\]
\end{proposition}

\begin{proof}
Write $Y_{H,b}(g) = g^{\top}Hg - \tr H - b$ and $\sigma^2(H,b) = 2\norm{H}_\Fro^2 + b^2$.

\emph{One row.} For $g \sim \gauss_r$, $\E Y_{H,b}^2 = \sigma^2$ and
$\E Y_{H,b}^4 \le 19\sigma^4$ \formal{Head.SmallBallFourth.paper\_main\_head\_small\_ball\_moments}:
diagonalize $H = O^{\top}\diag(a)O$, so that $Y = \sum_ia_i(w_i^2 - 1) - b$
with $w = Og \sim \gauss_r$, and compute with the moments $\E v = 0$,
$\E v^2 = 2$, $\E v^3 = 8$, $\E v^4 = 60$ of $v = w_i^2 - 1$:
$\E Y^4 = 48\sum a_i^4 + 12(\sum a_i^2)^2 - 32b\sum a_i^3 + 12b^2\sum a_i^2 + b^4$,
and $\sum a_i^4 \le (\sum a_i^2)^2$, $-32b\sum a_i^3 \le 16b^2\sum a_i^2 + 16\sum a_i^4$
give $\E Y^4 \le 19(2\sum a_i^2 + b^2)^2$. The paper's hypercontractivity
bound $\E Y^4 \le 81\sigma^4$ is not used. Paley--Zygmund
(Theorem~\ref{thm:pz} with $X = Y^2$, $\theta = \sigma^2/2$) gives
$\Prob\{Y^2 \ge \sigma^2/2\} \ge (\sigma^2/2)^2/(19\sigma^4) = 1/76$
\formal{Head.SmallBallPZ.paper\_main\_head\_small\_ball\_probability}.
Also, with $A(g) = gg^{\top} - \Id$, $\E\norm{A(g)}_\Fro^2 = r(r+1)$
\formal{Data.FeatureMoments.paper\_main\_data\_feature\_frobenius}, so
$\Prob\{\norm{A(g)}_\Fro^2 \ge \Psi\} \le 1/1216$ for $\Psi = 1216r(r+1)$.

\emph{Net.} The parameters $(H, b)$ with $\sigma^2(H,b) = 1$ form the unit
sphere of $\R^{\Dd_r + 1}$ in the coordinates $(\sqrt2\,\mathrm{code}(H), b)$,
and $\abs{Y_{H,b}(g) - Y_{H',b'}(g)} \le \sigma(H - H', b - b')\sqrt{\norm{A(g)}_\Fro^2/2 + 1}$
by Cauchy--Schwarz \formal{Head.SmallBallNet.paper\_main\_head\_small\_ball\_value\_diff}.
Take an $\epsilon$-net $(H_t, b_t)_{t<s}$ of the unit parameter sphere with
$\epsilon^2(\Psi/2 + 1) = 1/25$ and $\ln s \le C(\epsilon)(\Dd_r + 1) \le 504\ln2\,(r+1)^3$
\formal{Head.SmallBallNet.paper\_main\_head\_small\_ball\_param\_net}.

\emph{Counts.} For each $t$, the number of rows with $Y_{H_t,b_t}(g_i)^2 \ge 1/2$
is at least $N/304$ outside an event of probability $\exp(-\frac1{76}N(\frac12 - \frac{\ln2}{4}))$:
the count of rows in an event of probability $p$ is at most $pN/4$ with
probability at most $2^{pN/4}(1 - p/2)^N \le \exp(-pN(\frac12 - \frac{\ln 2}{4}))$
\formal{Head.SmallBallCount.paper\_main\_head\_row\_count\_lower\_exp}, by
the binomial recursion obtained by splitting off one row; and the number
of rows with $\norm{A(g_i)}_\Fro^2 \ge \Psi$ is below $N/608$ outside an event
of probability $\exp(N/1216 - (N/608)\ln 2) = \exp(-\frac{2\ln2 - 1}{1216}N)$
\formal{Head.SmallBallCount.paper\_main\_head\_row\_count\_upper\_exp}.
Let $\mathcal{B}$ be the union of these $s + 1$ events; its probability is at
most $\exp(504\ln2(r+1)^3 - \frac{N}{76}(\frac12 - \frac{\ln2}4)) + \exp(-\frac{2\ln2-1}{1216}N)$,
and the hypothesis $306432(r+1)^3 \le N$ makes this at most
$\exp(-\frac{2\ln2-1}{2432}N)$ \formal{Head.SmallBallUniform.paper\_main\_head\_small\_ball\_uniform}.

\emph{Deterministic step.} Outside $\mathcal{B}$, let $\sigma^2(H,b) = 1$ (the
general case follows by homogeneity) and let $t$ be a net point within
$\epsilon$. On the at least $N/304 - N/608 = N/608$ rows where
$Y_{H_t,b_t}^2 \ge 1/2$ and $\norm{A(g_i)}_\Fro^2 < \Psi$,
$\abs{Y_{H,b}(g_i)} \ge 1/\sqrt2 - \epsilon\sqrt{\Psi/2 + 1} \ge 1/\sqrt2 - 1/5 \ge 1/2$.
Removing at most $N/1216$ rows leaves at least $N/1216$ of them, whose
squares sum to at least $N/4864$ \formal{Head.SmallBallUniform.paper\_main\_head\_small\_ball\_good}.
\end{proof}

\begin{proposition}[Projected bulk escape \formal{Hybrid.Bulk.paper\_main\_hybrid\_bulk\_event}, \formal{Projected.BulkConditional.paper\_main\_projected\_psd\_conditional\_tail}, \formal{Projected.BulkEscape.paper\_main\_projected\_escape\_tail}; paper Lemma~4.11]\label{prop:bulk}
For $\gamma > 1/4$ there are $\eta, \delta, t, c, c_2 > 0$ and $d_0$, with
$\delta < 1/1216$, such that for $d \ge d_0$, $\abs{N - \gamma d^2} \le \eta d^2$,
$r \le \sqd + 1$, $r + m_i + 1 \le d$, $q \le r^2 + 1$, $k = \lfloor\delta N\rfloor$
and every head feature matrix $F \in \R^{N\times q}$ with entries measurable in
the head sample: the event $E_{\mathrm{bulk}}$ that some $J$ with $\abs{J} \ge N - k$,
some unit $v$ in the coded positive semidefinite sphere of $\Sym^m$ and some
$c \in \R^q$ have $\norm{S_J(Zv + Fc)}_2 \le \sqrt{\abs{J} - q - 1} - \sqrt{\Dd_m/2} - td$
is measurable, with probability at most $e^{-c_2d^2}$ under the joint law of
head and bulk, and outside it every such $J, v, c$ satisfy
$\norm{S_J(Zv + Fc)}_2 > cd$ and $\abs{J} \ge (1 - 1/1216)N$. Here $Z \in \R^{N\times\Dd_m}$
is the bulk array and $S_J$ the selection of the rows of $J$.
\end{proposition}

\begin{proof}
Condition on the head, so $F$ is fixed. For a fixed $J$ let $F_J = S_JF$;
there is a matrix $V$ with $s \ge \abs{J} - q$ orthonormal rows and $VF_J = 0$
(keep the rows of an orthogonal diagonalization of $F_JF_J^{\top}$ with
eigenvalue $0$; the others number at most $\rank F_J \le q$)
\formal{Projected.BulkComplement.paper\_main\_projected\_complement\_mx}.
Then $\norm{S_J(Zv + Fc)}_2 \ge \norm{VS_JZv}_2$ (Bessel), and $VS_JZ$ is an
$s\times\Dd_m$ array with independent standard Gaussian entries
(Proposition~\ref{prop:mulmx}, \formal{Projected.BulkLaw.paper\_main\_projected\_law}).
The coded positive semidefinite sphere $T$ has Gaussian width at most
$\sqrt{\Dd_m/2}$ (Proposition~\ref{prop:psdwidth}), so Theorem~\ref{thm:escape}
gives $\Prob\{\inf_{v\in T}\norm{VS_JZv}_2 \le \sqrt{s - 1} - \sqrt{\Dd_m/2} - u\} \le e^{-u^2/2}$
\formal{Gaussian.Lipschitz.EscapeTail.paper\_main\_gaussian\_escape\_sqrt\_tail}.
The number of sets $J$ with $\abs{J} \ge N - k$ is at most
$\exp(k(1 + \ln(N/k)))$ \formal{Projected.BulkUnion.paper\_main\_projected\_count\_entropy},
and the union over $J$ with $u = td$ has conditional probability at most
$\exp(k(1 + \ln(N/k)) - t^2d^2/2) \le e^{-c_2d^2}$ by the choice of $\delta$ and
$t$ \formal{Hybrid.Density.paper\_main\_hybrid\_density\_margin}; the same
choice gives $\sqrt{(N - k) - q - 1} - \sqrt{\Dd_m/2} - td \ge cd$, since
$N - k \ge (1 - \delta)(\gamma - \eta)d^2$, $q \le d + 2\sqd + 2$ and
$\Dd_m/2 \le (d+1)^2/4$, while $\sqrt{(1-\delta)(\gamma - \eta)} > 1/2$ for
$\delta, \eta$ small ($\gamma > 1/4$). Measurability of the joint event
holds because the infimum over $(v, c)$ is witnessed on countable nets
\formal{Projected.BulkConditional.paper\_main\_projected\_conditional\_tail}.
\end{proof}

\begin{lemma}[Bulk operator bound \formal{Hybrid.Assembly.paper\_main\_hybrid\_operator\_event}; paper Proposition~2.8, (2.16), and Lemma~4.11, upper bound]\label{lem:bulkop}
Let $\rho > 0$ with $\rho m^2 \le N$. Outside a measurable event of
probability at most $2e^{-N/2}$, $\norm{\mathsf{R}(B)}_2 \le C(\rho)\sqrt N\norm{B}_\Fro$
for every symmetric $B$, with $C(\rho) = 4\sqrt{20\ln2(1 + 2/\rho) + 1}$.
\end{lemma}

\begin{proof}
The $N\times\Dd_m$ array $Z$ of the bulk coordinates has operator norm at most
$2B$ outside probability $2\exp(20\ln2(N + \Dd_m) - B^2/2)$ (nets of the two
spheres and Proposition~\ref{prop:linear}, Lemma~\ref{lem:opnet}); with
$B^2 = 2(20\ln2(N + \Dd_m) + N/2)$ and $\Dd_m \le m^2 \le N/\rho$ this is
$2e^{-N/2}$ and $2B \le C(\rho)\sqrt N$ \formal{Gaussian.Symmetric.ForwardConstant.paper\_main\_gaussian\_forward\_threshold}.
\end{proof}

\begin{lemma}[Operator norms by nets \formal{Gaussian.Tail.Operator.paper\_main\_gaussian\_operator\_tail}; paper Lemma~2.5, with worse constants]\label{lem:opnet}
For an $N\times M$ array $\Gamma$ of independent standard Gaussian entries,
$N, M \ge 1$ and $B > 0$, $\Prob\{\norm{\Gamma}_\op > 2B\} \le 2\exp(20\ln2\,(N + M) - B^2/2)$.
\end{lemma}

\begin{proof}
For $1/4$-nets $(z_y)$, $(v_t)$ of the two unit spheres made of unit vectors,
with $\ln(\abs{z}\abs{v}) \le 20\ln2(N+M)$ (Theorem~\ref{thm:nets} with
$(2e(1 + 8))^N \le 2^{20N}$, through the grid nets of \formal{Gordon.Entropy}),
$\max_{y,t}\abs{z_y^{\top}\Gamma v_t} \le B$ implies $\norm{\Gamma}_\op \le 2B$
\formal{Gordon.OperatorNet.paper\_main\_gordon\_operator\_discretization},
and each $z_y^{\top}\Gamma v_t$ is a unit linear form of the entries
(Proposition~\ref{prop:linear}).
\end{proof}

\begin{proposition}[The deterministic gap \formal{Hybrid.Residual.paper\_main\_hybrid\_residual\_gap}, \formal{Hybrid.Loss.paper\_main\_hybrid\_risk\_gap}]\label{prop:detgap}
Let $\beta_0, C_0 > 0$, $k$ with $\delta N < k + 1$, and suppose:
(1) for every symmetric $H$, real $b$ and $J$ with $\abs{J} \ge N - k$,
$\sum_{i\in J}\mathsf{K}(H,b)_i^2 \ge \frac N{4864}(2\norm{H}_\Fro^2 + b^2)$;
(2) for every $J$ with $\abs{J} \ge N - k$, every unit coded positive
semidefinite $v$ and every head coefficient vector $c$,
$\norm{S_J(Zv + Fc)}_2 > \beta_0\sqrt N$, where $F$ is the $N\times(r^2+1)$
matrix with rows $(\mathrm{vec}(g_ig_i^{\top} - \Id), -1)$, so that
$\mathsf{K}(H,b) = F\,(\mathrm{vec}H, b)$;
(3) $\norm{\mathsf{R}(B)}_2 \le C_0\sqrt N\norm{B}_\Fro$ for every symmetric $B$.
Then for every symmetric $H$, $B \succeq 0$ and $b$ with
$\norm{H}_\Fro^2 + \norm{B}_\Fro^2 \ge 1/4$, the hybrid risk of $(H, b, B)$ is at
least $e^{-1}\min(\delta, c_4^2)$, where $c_4 = c_4(\beta_0, C_0) > 0$.
\end{proposition}

\begin{proof}
Fix $(H, b, B)$ and $u = \mathsf{K}(H,b) + \mathsf{R}(B)$. For every $J$ with
$\abs{J} \ge N - k$ we claim $\norm{S_Ju}_2 \ge c_4\sqrt N$. If $\norm{B}_\Fro \ge \theta_0$,
with $\theta_0 = \min(1/4, 1/(8\cdot4864^{1/2}C_0))$ say: by (2) and
homogeneity, $\norm{S_J(\mathsf{R}(B) + \mathsf{K}(H,b))}_2 = \sqrt2\norm{B}_\Fro\norm{S_J(Zv + Fc)}_2 \ge \sqrt2\beta_0\sqrt N\norm{B}_\Fro$
with $v = \mathrm{code}(B)/\norm{B}_\Fro$ and $c = (\mathrm{vec}H, b)/(\sqrt2\norm{B}_\Fro)$.
Otherwise $\norm{H}_\Fro^2 \ge 1/4 - \theta_0^2 \ge 3/16$, and (1), (3) give
$\norm{S_Ju}_2 \ge \norm{S_J\mathsf{K}}_2 - \norm{\mathsf{R}(B)}_2 \ge \sqrt{N/4864}\sqrt{2\cdot3/16} - C_0\theta_0\sqrt N$,
which is positive for $\theta_0$ small. Now let $I = \{i : \abs{u_i} > 1\}$. If
$\abs{I} > k$ then the risk is at least $e^{-1}(k+1)/N > e^{-1}\delta$; otherwise
$J = I^c$ has $\abs{J} \ge N - k$, every $i \in J$ has $\varphi(u_i) \ge e^{-1}u_i^2$,
and the risk is at least $e^{-1}\norm{S_Ju}_2^2/N \ge e^{-1}c_4^2$
\formal{Hybrid.Loss.paper\_main\_hybrid\_loss\_gap}.
\end{proof}

\begin{proof}[Proof of Theorem~\ref{thm:hybrid}]
Take the constants of Proposition~\ref{prop:bulk} and $\rho = \gamma/2$ in
Lemma~\ref{lem:bulkop} (so $\rho m^2 \le N$ for large $d$). Let $E$ be the
union of the head event of Proposition~\ref{prop:smallball} (for $r \ge 1$;
for $r = 0$, $\mathsf{K}(0, b) = -b\one$ satisfies (1) trivially), of the bulk
event of Proposition~\ref{prop:bulk} with $F$ the head feature matrix, and of
the operator event. Its probability is at most
$\exp(-\frac{2\ln2-1}{2432}N) + e^{-c_2d^2} + 2e^{-N/2} \le e^{-cd^2}$ for large
$d$. Outside $E$, Proposition~\ref{prop:detgap} applies with
$\beta_0 = c/\sqrt{\gamma + \eta}$ (as $cd \ge \beta_0\sqrt N$) and
$C_0 = C(\rho)$.
\end{proof}

\begin{corollary}[Expectation under the hybrid law \formal{Hybrid.FiniteMean.paper\_main\_hybrid\_finite\_goe\_mean}; paper Lemma~4.12, mean bound]\label{cor:goe_mean}
With $\gamma, \eta, e, c, d_0$ as in Theorem~\ref{thm:hybrid}, for all
$d \ge d_0$, $N$, $r$, $m_i$ as there, every $\beta > 0$ and every prior
$(H_j, b_j, B_j)_{j\in I}$ with $H_j$ symmetric, $B_j \succeq 0$ and
$\norm{H_j}_\Fro^2 + \norm{B_j}_\Fro^2 \ge 1/4$:
$\E\mathcal{Z}_{\mathrm{goe}} \ge e(1 - e^{-cd^2})$.
\end{corollary}

\begin{proof}
Outside $E$ (transported to the row layout of $\mathcal{Z}_{\mathrm{goe}}$
by the coordinate permutation, which preserves the law), every state $j$
has $\mathcal{A}_j = N\cdot(\text{hybrid risk of }(H_j, b_j, B_j)) \ge Ne$, so
$\mathcal{Z}_{\mathrm{goe}} \ge e$ (Lemma~\ref{lem:Z}); and $\mathcal{Z} \ge 0$
everywhere. Hence $\E\mathcal{Z}_{\mathrm{goe}} \ge e\Prob(E^c)$
\formal{Loss.ExpectedGap.paper\_main\_loss\_expected\_gap}.
\end{proof}

\subsection{Universality for the finite prior}\label{sec:universality}

\begin{theorem}[Comparison of the two means \formal{Hybrid.FiniteMean.paper\_main\_hybrid\_finite\_mean\_comparison}; replaces paper Propositions~4.3 and 4.8 and Lemma~4.9]\label{thm:universality}
For every $\beta > 0$ and $\eta > 0$ there is $C = C(\beta, \eta) \ge 0$ such
that for all $N \ge 1$, all dimensions $r, m_i$, all finite priors with $B_j$
symmetric, $\norm{B_j}_\op \le \varepsilon$ and $\norm{B_j}_\Fro \le S_f$, and all
offsets,
\[
\abs{\E\mathcal{Z}_{\mathrm{dat}} - \E\mathcal{Z}_{\mathrm{goe}}} \le 2\eta + C\varepsilon S_f .
\]
\end{theorem}

The free-entropy universality of Bandeira and Maillard, which the paper
cites, is replaced by the following three steps.

\begin{proposition}[Row replacement \formal{Loss.Replacement.paper\_main\_loss\_array\_comparison}, \formal{Loss.BlockComparison.paper\_main\_loss\_conditional\_prefix\_comparison}; replaces paper Proposition~4.8, step 1]\label{prop:replacement}
Let $\ell^{\mathrm{Q}}_{ij}$, $\ell^{\mathrm{G}}_{ij}$ ($i < N$, $j \in I$) be
continuous row losses with values in $[0,1]$ on rows of lengths $p$ and $q$.
If for every $i$ and every $f : I \to \R$,
$\abs{\E^{(p)}\min_\beta^j(f_j + \ell^{\mathrm{Q}}_{ij}) - \E^{(q)}\min_\beta^j(f_j + \ell^{\mathrm{G}}_{ij})} \le B$,
then for every $f$,
$\abs{\E^{(Np)}\min_\beta^j(f_j + \sum_i\ell^{\mathrm{Q}}_{ij}) - \E^{(Nq)}\min_\beta^j(f_j + \sum_i\ell^{\mathrm{G}}_{ij})} \le NB$.
\end{proposition}

\begin{proof}
Replace the rows one at a time: with $i$ rows of the second kind and $N - i$
of the first, the difference between consecutive hybrids is, after
integrating out all the other rows (Fubini, with the other rows fixed, the
functional is a soft minimum of $f'_j = f_j + \sum_{i'\ne i}\ell_{i'j}$ plus
the $i$-th row loss), at most $B$ by hypothesis. Sum over $i$.
\end{proof}

\begin{proposition}[One row \formal{Loss.SoftCost.paper\_main\_loss\_row\_value\_comparison}, \formal{Loss.RowIncrement.paper\_main\_loss\_row\_increment\_comparison}; replaces paper Proposition~4.8, step 2]\label{prop:onerow}
For $\beta, \eta > 0$ there is $C \ge 0$ such that for every $n$, every
$f : I \to \R$, offsets $(a_j)$ and symmetric $T_j \in \R^{(n+1)\times(n+1)}$
with $\norm{T_j}_\op \le \varepsilon$ and $\norm{T_j}_\Fro \le S_f$, with
$Q_j(x) = x^{\top}T_jx - \tr T_j$ for $x \sim \gauss_{n+1}$ and
$G_j(z) = \ip{z}{\mathrm{loading}(T_j)}$ the Gaussian linear forms with the
same covariances, $z \sim \gauss_{\Dd_n}$,
\[
\Bigl|\E\min\nolimits_\beta^j\bigl(f_j + \varphi(a_j + Q_j(x))\bigr) - \E\min\nolimits_\beta^j\bigl(f_j + \varphi(a_j + G_j(z))\bigr)\Bigr| \le 2\eta + C\varepsilon S_f .
\]
\end{proposition}

\begin{proof}
\emph{Replica expansion.} For losses $\ell_j \in [0,1]$,
$\min_\beta^j(f_j + \ell_j) - \min_\beta^jf_j = -\beta^{-1}\ln\mathcal{A}$, where
$\mathcal{A} = \sum_j\pi_je^{-\beta\ell_j}$ is the average of $e^{-\beta\ell_j}$
under the Gibbs weights $\pi_j \propto e^{-\beta f_j}$, and
$e^{-\beta} \le \mathcal{A} \le 1$ \formal{Loss.Gibbs.paper\_main\_gibbs\_row\_increment}.
With $u_j = 1 - e^{-\beta\ell_j} \in [0, 1 - e^{-\beta}]$,
$-\ln\mathcal{A} = -\ln(1 - \sum_j\pi_ju_j) = \sum_{k\ge1}\frac1k\bigl(\sum_j\pi_ju_j\bigr)^k$,
and the truncation at degree $k_0$ has error at most
$q^{k_0}/(1 - q)$ with $q = 1 - e^{-\beta}$, uniformly in everything
\formal{Loss.LogPolynomial.paper\_main\_loss\_log\_remainder}; choose $k_0$
with $q^{k_0}/((1-q)\beta) \le \eta$ \formal{Loss.Replica.paper\_main\_loss\_row\_poly\_uniform}.
Each term $(\sum_j\pi_ju_j)^k$ is the average over $k$ independent replicas
$j_1, \dots, j_k$ drawn from $\pi$ of the product $\prod_{l\le k}u_{j_l}$
\formal{Loss.Replica}.

\emph{Comparison of a product.} With $\ell_j = \varphi(a_j + Q_j(x))$, the
factor $u_j = 1 - \exp(-\beta\varphi(a_j + Q_j(x)))$ is a smooth bounded
function of $Q_j(x)$ with derivatives bounded by $C_\beta = 1 + 10\beta + 16\beta^2$
\formal{Loss.SmoothFactor.paper\_main\_loss\_factor\_smooth}, and the
product of $k$ factors is a $C^2$ function $F$ of the $k$ features
$(Q_{j_1}, \dots, Q_{j_k})$ with $F$ and its first two derivatives bounded by
$(4C_\beta)^k$ \formal{Loss.SmoothProduct.paper\_main\_loss\_test\_smooth}.
By Proposition~\ref{prop:smooth_comparison} below,
$\abs{\E F(Q_{j_1}(x), \dots) - \E F(G_{j_1}(z), \dots)} \le 2k^2(4C_\beta)^k\varepsilon S_f$
\formal{Loss.ReplicaComparison.paper\_main\_loss\_product\_comparison};
averaging over the replicas keeps the bound
\formal{Loss.ReplicaIntegral.paper\_main\_loss\_replica\_comparison}, and
summing the $k_0$ terms with their factors $1/k$ gives a constant
$C' = C'(\beta, k_0)$ with
$\abs{\E(\text{truncated}^{\mathrm{Q}}) - \E(\text{truncated}^{\mathrm{G}})} \le C'\varepsilon S_f$
\formal{Loss.RowPolynomial.paper\_main\_loss\_row\_polynomial\_comparison}.
Adding the two truncation errors $\eta$ gives the claim with $C = C'/\beta$.
\end{proof}

\begin{proposition}[Fixed-dimensional smooth comparison \formal{Projection.PublicOperator.paper\_main\_projection\_public\_comparison}, \formal{Projection.Comparison.paper\_main\_projection\_smooth\_comparison}; replaces paper Proposition~4.8, step 3]\label{prop:smooth_comparison}
Let $T_0, \dots, T_{K-1}$ be symmetric $(n+1)\times(n+1)$ matrices with
$\norm{T_a}_\op \le \varepsilon$ and $\norm{T_a}_\Fro \le S_f$, and let $F : \R^K \to \R$
be of class $C^2$ with $F$ and its first and second partial derivatives
bounded by $C_F$. Then
\[
\Bigl|\E F\bigl((x^{\top}T_ax - \tr T_a)_{a<K}\bigr) - \E F\bigl((\ip{z}{\mathrm{loading}(T_a)})_{a<K}\bigr)\Bigr| \le 2K^2C_F\varepsilon S_f ,
\]
for $x \sim \gauss_{n+1}$ and $z \sim \gauss_{\Dd_n}$, where the loadings are
chosen so that the covariances agree: $\E G_aG_b = \E Q_aQ_b = 2\ip{T_a}{T_b}$.
\end{proposition}

\begin{proof}
Put the two Gaussian vectors side by side and interpolate the features:
$\Xi_a(t) = \sqrt t\,Q_a(x) + \sqrt{1-t}\,G_a(z)$. The map
$t \mapsto \E F(\Xi(t))$ is differentiable on $(0,1)$ with derivative
$\E\sum_a\partial_aF(\Xi(t))\bigl(\frac{Q_a}{2\sqrt t} - \frac{G_a}{2\sqrt{1-t}}\bigr)$
\formal{Projection.Interpolation.paper\_main\_projection\_interpolation\_derivative}.
Integration by parts in $z$ (Theorem~\ref{thm:ibp}) turns the $G$ term into
$\frac12\sum_{a,b}\E\,\partial_a\partial_bF(\Xi(t))\,\E G_aG_b$, and integration
by parts in $x$ turns the $Q$ term, since $Q_a$ is a centered quadratic
form, into $\frac12\sum_{a,b}\E\,\partial_a\partial_bF(\Xi(t))\cdot2\ip{T_a}{T_b}$
plus an error term that involves the second derivatives of $F$ and the
\emph{cross products} $x^{\top}T_aT_bx$ centered at their means
\formal{Projection.Stein.paper\_main\_projection\_quadratic\_stein},
\formal{Projection.Error.paper\_main\_projection\_mixed\_error}. The
covariance terms cancel, and the error term is bounded, by Cauchy--Schwarz
and the moments of quadratic forms, by $2C_F\sum_{a,b}\norm{T_a}_\op\norm{T_b}_\Fro \le 2K^2C_F\varepsilon S_f$
\formal{Projection.ErrorBound.paper\_main\_projection\_error\_bound};
integrate over $[0,1]$. The loadings and the orthogonality of the two blocks
are handled by \formal{Projection.Independent.paper\_main\_projection\_independent\_comparison}
and \formal{Projection.Public.paper\_main\_projection\_goe\_covariance}.
\end{proof}

\begin{proof}[Proof of Theorem~\ref{thm:universality}]
For a row $(g, y)$ or $(g, z)$ with the head $g$ fixed, the offsets
$a_j = g^{\top}H_jg - \tr H_j - b_j$ are numbers, and
Proposition~\ref{prop:onerow} with $T_j = B_j$ bounds the difference of the
expected row values by $2\eta + C\varepsilon S_f$; integrating over $g$
(Fubini, \formal{Loss.BlockComparison}) gives the one-row hypothesis of
Proposition~\ref{prop:replacement} with $B = 2\eta + C\varepsilon S_f$, and
the normalization by $N$ and the entropy shift (which is the same on both
sides) give the claim \formal{Hybrid.FiniteMean.paper\_main\_hybrid\_finite\_row\_comparison}.
\end{proof}

\subsection{Concentration}\label{sec:concentration}

\begin{proposition}[Bounded differences for Gaussian rows \formal{Loss.ConcentrationRow.paper\_main\_concentration\_rows}, \formal{Loss.ConcentrationTail.paper\_main\_concentration\_lower\_tail}, \formal{Loss.ConcentrationSoftMin.paper\_main\_softmin\_lower\_tail}; replaces McDiarmid's inequality, paper line 2330]\label{prop:bdd}
Let $F$ be a continuous bounded function of $N$ rows of $p$ Gaussian
coordinates such that changing one row changes $F$ by at most $c$. Then
$\E e^{\lambda(F - \E F)} \le e^{\lambda^2Nc^2/2}$ for every real $\lambda$, and
$\Prob\{F \le \E F - t\} \le e^{-t^2/(2Nc^2)}$ for $t > 0$. For
$F = \mathcal{Z}_{\mathrm{dat}}$ (Lemma~\ref{lem:Z}, $c = 1/N$):
$\Prob\{\mathcal{Z}_{\mathrm{dat}} \le \E\mathcal{Z}_{\mathrm{dat}} - t\} \le e^{-Nt^2/2}$.
\end{proposition}

\begin{proof}
Induction on the rows: integrating out the last row, $F - \E_{\text{last}}F$
has absolute value at most $c$ for fixed other rows, and
$\E_{\text{last}}e^{\lambda(F - \E_{\text{last}}F)} \le e^{\lambda^2c^2/2}$,
by the convexity of the exponential on $[-c, c]$: for a centered variable
$V$ with $\abs{V} \le c$, $e^{\lambda V} \le \cosh(\lambda c) + V\sinh(\lambda c)/c$,
whose mean is $\cosh(\lambda c) \le e^{\lambda^2c^2/2}$. The function
$\E_{\text{last}}F$ of the remaining rows has the same oscillation bound,
so the induction closes. Chernoff's bound gives the tail.
\end{proof}

\begin{corollary}[Lower tail of the data functional \formal{Unsat.FiniteTail.paper\_main\_unsat\_ordered\_tail}, \formal{Hybrid.FiniteTail.paper\_main\_hybrid\_finite\_data\_lower\_tail}; replaces paper Lemma~4.12]\label{cor:tail}
For $\gamma > 1/4$ there are $\kappa \le \gamma/2$, $e, c > 0$ and $d_0$ such
that for all $\beta, \eta > 0$ there is $C \ge 0$ with: for $d \ge d_0$,
$\abs{N - \gamma d^2} \le \kappa d^2$, $r \le \sqd + 1$, $r + m_i + 1 \le d$, every
prior with $H_j$ symmetric, $B_j \succeq 0$, $\norm{H_j}_\Fro^2 + \norm{B_j}_\Fro^2 \ge 1/4$,
$\norm{B_j}_\op \le \varepsilon$, $\norm{B_j}_\Fro \le S_f$, and every $t > 0$,
\[
\Prob\Bigl\{\mathcal{Z}_{\mathrm{dat}} \le e(1 - e^{-cd^2}) - (2\eta + C\varepsilon S_f) - t\Bigr\} \le e^{-Nt^2/2}.
\]
\end{corollary}

\begin{proof}
Corollary~\ref{cor:goe_mean}, Theorem~\ref{thm:universality} and
Proposition~\ref{prop:bdd}.
\end{proof}

\subsection{Assembly of the unsatisfiable phase}\label{sec:unsat_assembly}

\begin{proposition}[The exceptional event of an orientation \formal{Orientation.Exception.paper\_main\_orientation\_exception}, \formal{Unsat.Orientation.paper\_main\_unsat\_orientation\_constants}; replaces paper Proposition~4.5 for one net point]\label{prop:exception}
For $\gamma > 1/4$ there are $\kappa \le \gamma/2$, $e, c, L > 0$ with
$\frac\gamma2L^2 = 109$, $\varepsilon > 0$, $0 < \tau \le 1/8$ with
$\varepsilon/2 + \tau^2 \le \varepsilon$, and $d_0$, such that for all $d \ge d_0$,
every orthogonal $O$, every $r < d$ with $r \le \sqd + 1$, every $N$ with
$\abs{N - \gamma d^2} \le \kappa d^2$ and $e^{-cd^2} \le 1/4$: there is a
measurable event $E_O$ with $\Prob_{N,d}(E_O) \le \exp(-N(e/4)^2/2)$ that
contains every sample with $QQ^{*} \preceq K\Id$, $K \le NL^2$, which admits an
admissible $(A, b)$ for $O$ (at $r$, $\varepsilon$, $L$) with
$\mathcal{L}(A, b) \le 16L\tau$.
\end{proposition}

\begin{proof}
Choose, in this order: $\kappa, e, c, d_0$ from Corollary~\ref{cor:tail};
$L$ with $\frac\gamma2L^2 = 109$ \formal{Loss.EntropyBudget.paper\_main\_loss\_feature\_constant};
$\delta, \beta, \eta$ with $4(L+1)\delta \le e/16$, $2\eta = e/32$ and
$3C(\delta/(L+3))/(\frac\gamma2\beta) \le e/16$ \formal{Unsat.SupercriticalParameters.paper\_main\_unsat\_net\_parameters};
then $C = C(\beta,\eta)$ from Corollary~\ref{cor:tail}, and $\varepsilon, \tau$
with $\tau \le 1/8$, $\varepsilon/2 + \tau^2 \le \varepsilon$, $2C\varepsilon \le e/32$
and $16L\tau \le e/16$ \formal{Unsat.SupercriticalParameters.paper\_main\_unsat\_cutoff\_parameters}.
Given $O$ and $r$, take the prior of Proposition~\ref{prop:prior} for the
pinched admissible set (if it is empty, $E_O = \varnothing$), with
$S_f = 2$. Let $E_O$ be the pullback under the rotation by $O$ of the
event $\{\mathcal{Z}_{\mathrm{dat}} \le z\}$ with
$z = 16L\tau + 4(L+1)\delta + 3C(\delta/(L+3))/(\frac\gamma2\beta) \le 3e/16$;
it is measurable and has the same probability, since the rotation
preserves the law. By Proposition~\ref{prop:prior} (the density
$\frac\gamma2d^2 \le N$ bounds $\ln s/(\beta N)$ by $3C(\delta/(L+3))/(\frac\gamma2\beta)$),
every sample with a low-risk admissible candidate lies in $E_O$. By
Corollary~\ref{cor:tail} with $t = e/4$: the threshold there is at least
$e\cdot\frac34 - \frac e{16} - \frac e4 = \frac{7e}{16} > \frac{3e}{16} \ge z$,
so $\Prob(E_O) \le e^{-N(e/4)^2/2}$.
\end{proof}

\begin{proof}[Proof of Theorem~\ref{thm:unsat}]
Fix $\gamma \in (1/4, 1/2)$ and the constants of Proposition~\ref{prop:exception};
let $\varepsilon_0 = \varepsilon/2$ and $r_{\max} = \lceil 1/\varepsilon_0^2\rceil$,
a constant. Let $N = n_\gamma(d)$, so $\abs{N - \gamma d^2} \le 1 \le \kappa d^2$
and $N \le d^2$ for large $d$. For each $r \le r_{\max}$ take the $\tau$-net of
orientations of Proposition~\ref{prop:onets} and the events
$E_{r,k} = E_{O_{r,k}}$; their union $F_d$ is measurable with
\[
\Prob(F_d) \le (r_{\max} + 1)\exp\bigl((4 + 8r_{\max}/\tau)\ln2(r_{\max}+1)(d+1)\bigr)\,e^{-N(e/4)^2/2} \to 0,
\]
since $N(e/4)^2/2 \ge ad^2$ with $a = \gamma(e/4)^2/4$
\formal{Orientation.Union.paper\_main\_orientation\_union\_fit},
\formal{Unsat.Rates.paper\_main\_unsat\_orientation\_union\_limit}. Outside
the upper-edge event of Proposition~\ref{prop:upperedge}, $QQ^{*} \preceq 109d^2\Id \preceq NL^2\Id$
(as $\frac\gamma2d^2 \le N$), so by Proposition~\ref{prop:witness} every fit
lies in some $E_{r,k}$. Hence
$\Prob(\Fit_{N,d}) \le \Prob(\text{upper-edge event}) + \Prob(F_d) \to 0$
\formal{Unsat.Limit.paper\_main\_unsat\_fixed\_density\_limit\_from\_events}.
\end{proof}

\section{Assembly of the three theorems}\label{sec:assembly}

\begin{proof}[Proof of Theorem~\ref{thm:cond}]
Let $0 < \beta < 1/4$ and $N(d) = \lfloor\beta d^2\rfloor$. By
Theorem~\ref{thm:loweredge} there are events $E_d$ of vanishing probability
outside which $QQ^{*}$ is coercive with constant $d^2/128$; by
Proposition~\ref{prop:upperedge} there are events of vanishing probability
outside which $QQ^{*} \preceq 109d^2\Id$ (the counts satisfy $0 < N(d) \le d^2$
and the density is at least $\beta/2$ eventually
\formal{Fixed.CountBounds.paper\_main\_fixed\_count\_density\_bounds}).
Theorem~\ref{thm:sat} with $\alpha = \beta$, $c = 1/128$ and $C = 109$ gives
$\kappa = \kappa(\beta) \ge 1$ with
$\Prob_{N(d),d}(\Cond_{N(d),d}((2\kappa+1)^{-1}, 2\kappa+1)) \to 1$
\formal{Conditioned.FromLower.paper\_main\_conditioned\_from\_lower}. The
events are measurable (Proposition~\ref{prop:measurable}), and
Proposition~\ref{prop:fixed}(a) concludes
\formal{PaperMainConditionedLimit.paper\_main\_conditioned\_limit}.
\end{proof}

\begin{proof}[Proof of Theorem~\ref{thm:main}]
$\Fit_{n,d}$ is measurable by Proposition~\ref{prop:measurable}.
(a) If $\limsup n(d)/d^2 < 1/4$, pick $\alpha^* < 1/4$ with
$n(d)/d^2 \le \alpha^*$ eventually; Theorem~\ref{thm:cond} gives
$\lambda_\pm$ with $\Prob(\Cond_{n(d),d}(\lambda_-,\lambda_+)) \to 1$, and
$\Cond \subseteq \Fit$ \formal{Asymptotics.paper\_main\_conditioned\_subset},
so $\Prob(\Fit_{n(d),d}) \to 1$ (probabilities are at most one)
\formal{Asymptotics.paper\_main\_sat\_from\_conditioned}.
(b) Theorem~\ref{thm:unsat} and Proposition~\ref{prop:fixed}(b)
\formal{Unsat.Limit.paper\_main\_unsat\_supercritical\_limit}.
\end{proof}

\begin{proof}[Proof of Theorem~\ref{thm:bal}]
$\Bal_{n,d}$ is measurable by Proposition~\ref{prop:measurable}. For
$d > 0$, no sample is both in $\Fit_{n,d}$ and in $\Bal_{n,d}$
(Theorem~\ref{thm:alternative}, exclusion), so
$\Prob(\Bal_{n,d}) \le 1 - \Prob(\Fit_{n,d})$
\formal{Balanced.Probability.paper\_main\_balanced\_probability\_le\_complement};
below the threshold, Theorem~\ref{thm:main}(a) gives $\Prob(\Bal) \to 0$
\formal{Balanced.Probability.paper\_main\_balanced\_from\_sat}. Above the
threshold, $0 < d \le n(d)$ for all large $d$
\formal{Balanced.Supercritical.paper\_main\_supercritical\_spanning\_count},
the sample spans almost surely (Lemma~\ref{lem:spanning}), and a spanning
sample that is not a fit is in $\Bal_{n,d}$ (Theorem~\ref{thm:alternative}),
so $\Bal_{n,d}$ and the complement of $\Fit_{n,d}$ agree almost surely and
$\Prob(\Bal_{n(d),d}) = 1 - \Prob(\Fit_{n(d),d}) \to 1$ by
Theorem~\ref{thm:main}(b) \formal{Balanced.Supercritical.paper\_main\_balanced\_from\_unsat},
\formal{Balanced.Limit.paper\_main\_balanced\_supercritical\_limit}.
\end{proof}

\begin{remark}[Where the threshold enters]
In the satisfiable phase, the density $\alpha < 1/4$ enters through
$g = 1/2 - \sqrt\alpha > 0$: the width of the bounded-condition cone is
close to $d/2$ (Proposition~\ref{prop:conewidth}), the section widths keep
all but $\frac g{16}d$ of it (Theorem~\ref{thm:heads}), and Gordon's
inequality compares the width with $\sqrt N \le \sqrt\alpha\,d$
(Proposition~\ref{prop:gordon_margin}). The lower feature edge is also used
there, and it is proved only for $4N \le d^2$ (Theorem~\ref{thm:loweredge});
the paper claims it up to density $1/2$. In the unsatisfiable phase, the
density $\gamma > 1/4$ enters through the escape estimate for the projected
bulk: $\sqrt{(1-\delta)\gamma} - 1/2 > 0$ (Proposition~\ref{prop:bulk}). The
other ingredients (small ball, nets, universality, concentration) work at
any density; the formal proof restricts to $\gamma < 1/2$ only through the
upper-edge bound $N \le d^2$.
\end{remark}

\begin{remark}[What is not computed]
The constants $\kappa(\alpha^*)$, hence $\lambda_\pm$ of
Theorem~\ref{thm:cond}, and all thresholds in $d$ are produced by the proof
but not computed. The walk-count estimates of Section~\ref{sec:edge} need
$d \ge 2^{9600}$; the other thresholds are of the usual kind (polynomial
versus logarithmic comparisons).
\end{remark}

\begin{remark}[Axioms]
The formal development uses, besides the logic of Rocq, the classical
axioms of MathComp-Analysis (excluded middle, functional and propositional
extensionality, and a choice principle), which also underlie the informal
proof.
\end{remark}

\appendix
\section{Correspondence with the paper}\label{sec:correspondence}

This appendix is generated from the reconciled comparison
\texttt{doc/paper\_correspondence.md} of the repository, which holds the
full notes: for every row, the paper's lines, the formal statements, and,
for the rows on which the two independent comparisons disagreed, the
evidence and the verdict. The \emph{statement relation} compares the
formal statement with the paper's: \emph{same} (equivalent up to
notation, the values of unspecified constants and the thresholds on $d$),
\emph{weaker} (implied by the paper's statement but not conversely: worse
constants or rates, a missing direction or case, a special case, or an
added hypothesis on the sample), \emph{stronger} (implies the paper's
statement and says more), \emph{absent} (not needed). The \emph{argument
relation} compares the proofs: \emph{paper's route} (the same
construction, decomposition and intermediate claims, possibly with other
constants or other proofs of individual steps) or \emph{replaced} (another
argument). For a result the paper cites without proof, the argument
relation compares the formal proof with the proof of the cited source.

\subsection{The numbered results of the paper}\label{sec:corr_table}

The table lists the 39 theorems, propositions, lemmas and corollaries of
sections 1.2 to 4 of the paper in the paper's order (remarks and the
conjecture share the counter), where each appears in this document, the
formal counterparts, and the two relations.

{\small\setlength{\tabcolsep}{3pt}
\begin{longtable}{@{}p{23mm}p{22mm}p{47mm}p{13mm}p{15mm}p{22mm}@{}}
\caption{The numbered results of the paper and their formal counterparts.}\label{tab:corr}\\
\hline
Paper result & Here & Formal counterparts & Statement & Argument & Note\\
\hline
\endfirsthead
\multicolumn{6}{@{}l}{\small\emph{Table~\ref{tab:corr}, continued.}}\\
\hline
Paper result & Here & Formal counterparts & Statement & Argument & Note\\
\hline
\endhead
\hline
\endfoot
\hline
\endlastfoot
Thm.~1.2 \texttt{\small thm:main} & Thm.~\ref{thm:main} & \formal{PaperMainFinalAssembly.paper\_main\_results\_from\_conditioned}, \formal{PaperMainResults.paper\_main\_results}, \formal{MainStatement.main\_statement}, \formal{Asymptotics.paper\_main\_sat\_from\_conditioned}, \formal{Unsat.Limit.paper\_main\_unsat\_supercritical\_limit} & same & paper's route & measurability added\\
Thm.~1.3 \texttt{\small thm:conditioned} & Thm.~\ref{thm:cond} & \formal{PaperMainConditionedLimit.paper\_main\_conditioned\_limit}, \formal{PaperMainConditionedLimit.paper\_main\_conditioned\_fixed\_limit}, \formal{Fixed.Density.paper\_main\_conditioned\_from\_fixed}, \formal{Conditioned.FromLower.paper\_main\_conditioned\_from\_lower}, \formal{Conditioned.FromLower.paper\_main\_conditioned\_from\_lower\_exists} & same & paper's route & constants $(2\kappa+1)^{\pm1}$, not computed\\
Cor.~1.4 \texttt{\small cor:alternative} & Thm.~\ref{thm:bal} & \formal{Balanced.Probability.paper\_main\_balanced\_from\_sat}, \formal{Balanced.Limit.paper\_main\_balanced\_supercritical\_limit}, \formal{Balanced.Supercritical.paper\_main\_balanced\_from\_unsat}, \formal{Gaussian.Spanning.AlmostSure.paper\_main\_gaussian\_spanning\_ae}, \formal{Measurability.Balanced.pm\_balanced\_pd\_event\_measurable} & same & paper's route & spanning proved almost surely\\
Lem.~2.1 \texttt{\small lem:psd-width} & Prop.~\ref{prop:psdwidth} & \formal{Gaussian.Width.Psd.paper\_main\_gaussian\_psd\_projection\_second\_moment}, \formal{Gaussian.Width.Psd.paper\_main\_gaussian\_psd\_width\_bounds}, \formal{Cone.Projection.paper\_main\_psd\_moreau} & same & paper's route & Poincar\'e inequality for the width bounds\\
Lem.~2.2 \texttt{\small lem:widthcone} & Prop.~\ref{prop:conewidth} & \formal{Gaussian.Width.ConditionBound.paper\_main\_gaussian\_condition\_width\_lower}, \formal{Gaussian.Width.ConditionBound.paper\_main\_gaussian\_condition\_width\_large}, \formal{Cone.Condition.paper\_main\_cone\_shift\_scaled} & stronger & paper's route & explicit, for every $d$\\
Lem.~2.3 \texttt{\small lem:gordon} & Thm.~\ref{thm:gordon}, Cor.~\ref{cor:bil}, Prop.~\ref{prop:gordon_margin} & \formal{Gordon.Width.paper\_main\_gordon\_finite\_width\_margin}, \formal{Gordon.SoftMargin.paper\_main\_gordon\_soft\_width\_margin}, \formal{Gordon.Processes.paper\_main\_gordon\_process\_comparison}, \formal{Gaussian.Lipschitz.LowerMargin.paper\_main\_gaussian\_lipschitz\_lower\_margin}, \formal{GaussianAnalysis.comparison.escape.auxT\_bilT\_minmax}, \formal{GaussianAnalysis.comparison.gordon.gordon} & weaker & replaced & expectation form on finite nets\\
Lem.~2.4 \texttt{\small lem:escape} & Thm.~\ref{thm:escape} & \formal{Gaussian.Lipschitz.EscapeTail.paper\_main\_gaussian\_escape\_sqrt\_tail}, \formal{Gordon.EscapeExpectation.paper\_main\_gordon\_escape\_expectation}, \formal{GaussianAnalysis.comparison.escape.gordon\_escape} & weaker & paper's route & threshold $\sqrt{N-1}$\\
Lem.~2.5 \texttt{\small lem:gauss-rect} & Lem.~\ref{lem:opnet}; Thm.~\ref{thm:opnorm} (library, unused) & \formal{Gaussian.Tail.Operator.paper\_main\_gaussian\_operator\_tail}, \formal{Gaussian.Symmetric.FeatureTail.paper\_main\_gaussian\_operator\_exponential\_tail}, \formal{GaussianAnalysis.comparison.escape.gauss\_opnorm\_tail} & weaker & replaced & net bound, constant $20\ln 2$\\
Lem.~2.6 \texttt{\small lem:squarefree-edge} & Prop.~\ref{prop:upperedge}, Thm.~\ref{thm:loweredge} & \formal{SquareFree.EdgeUpperLimit.paper\_main\_squarefree\_upper\_bound}, \formal{SquareFree.EdgeUpperLimit.paper\_main\_squarefree\_upper\_event\_tends\_to\_zero}, \formal{SquareFree.LowerProbability.paper\_main\_squarefree\_lower\_probability}, \formal{SquareFree.Coercivity.paper\_main\_squarefree\_coercivity}, \formal{SquareFree.Coercivity.paper\_main\_squarefree\_coercivity\_generic} & weaker & replaced & bounds, lower edge for $4N \le d^2$\\
Lem.~2.7 \texttt{\small lem:diagonal-feature-block} & Prop.~\ref{prop:diag} & \formal{Feature.DiagonalLimit.paper\_main\_diag\_block\_tends\_to\_zero}, \formal{Feature.DiagonalLimit.paper\_main\_diag\_trace\_event\_bound}, \formal{Feature.DiagonalLimit.paper\_main\_diag\_trace\_event\_markov}, \formal{Feature.DiagonalMoment.paper\_main\_diag\_trace\_moment} & stronger & replaced & fourth trace moment, constant 3\\
Prop.~2.8 \texttt{\small prop:feature-edge} & Thm.~\ref{thm:loweredge}, Prop.~\ref{prop:upperedge}, Lem.~\ref{lem:bulkop} & \formal{Fixed.LowerProbability.paper\_main\_fixed\_lower\_probability}, \formal{Feature.UpperProbability.paper\_main\_feature\_upper\_edge}, \formal{Feature.UpperProbability.paper\_main\_feature\_upper\_event\_bound}, \formal{Feature.GramBlocks.paper\_main\_feature\_gram\_blocks}, \formal{Gaussian.Symmetric.PublicFeature.paper\_main\_public\_gaussian\_feature\_exceptional\_event}, \formal{Gaussian.Symmetric.ForwardConstant.paper\_main\_gaussian\_forward\_threshold} & weaker & paper's route & one sequence $N(d)$\\
Cor.~2.9 \texttt{\small cor:feature-sat} & Lem.~\ref{lem:bridge_features} & \formal{Feature.SpectralBridge.paper\_main\_raw\_feature\_coercive}, \formal{Feature.SpectralBridge.paper\_main\_raw\_feature\_upper}, \formal{Feature.SpectralBridge.paper\_main\_raw\_feature\_operator}, \formal{Conditioned.SpectralReduction.paper\_main\_subcritical\_raw\_spectral\_reduction} & weaker & paper's route & densities below $1/4$\\
Lem.~2.10 \texttt{\small lem:data} & Lems.~\ref{lem:radius}--\ref{lem:sparse} & \formal{Gaussian.Sample.Radius.paper\_main\_data\_radius\_good}, \formal{Gaussian.Sample.Radius.paper\_main\_data\_radius\_probability}, \formal{Data.SampleTail.paper\_main\_data\_coordinate\_tail}, \formal{Data.CoherenceLimit.paper\_main\_data\_coherence\_log\_good}, \formal{Data.CoherenceLimit.paper\_main\_data\_coherence\_log\_tail}, \formal{Head.DataLog.paper\_main\_head\_data\_log\_event}, \formal{Head.DataLog.paper\_main\_head\_data\_log\_probability} & stronger & paper's route & explicit tails; coordinates used at $d^{1/4}$\\
Lem.~2.11 \texttt{\small lem:hw} & Lem.~\ref{lem:hw} & \formal{Gaussian.Orthogonal.HansonWright.paper\_main\_gaussian\_quadratic\_tail}, \formal{Gaussian.Sample.ResidualTail.paper\_main\_gaussian\_residual\_tail}, \formal{Gaussian.Orthogonal.Tail.paper\_main\_gaussian\_orthogonal\_square\_two\_sided\_tail}, \formal{GaussianAnalysis.comparison.chi\_square.gauss\_centered\_chi2\_abs\_tail} & same & replaced & Gaussian proof, $c = 1/8$\\
Lem.~2.12 \texttt{\small lem:homogeneous} & Lem.~\ref{lem:homog}, Thm.~\ref{thm:alternative} & \formal{Geometry.Homogeneous.paper\_main\_homogeneous\_kernel}, \formal{Geometry.Homogeneous.paper\_main\_homogeneous\_witness}, \formal{Geometry.Separation.paper\_main\_homogeneous\_alternative}, \formal{Geometry.Separation.paper\_main\_homogeneous\_existence}, \formal{Geometry.Separation.paper\_main\_fit\_alternative}, \formal{Geometry.Exclusion.paper\_main\_fit\_balanced\_exclusive} & same & paper's route & exclusion without spanning\\
Thm.~3.1 \texttt{\small thm:sat} & Thm.~\ref{thm:sat} & \formal{PaperMainConditionedLimit.paper\_main\_conditioned\_limit}, \formal{PaperMainConditionedLimit.paper\_main\_conditioned\_fixed\_limit}, \formal{Conditioned.FromLower.paper\_main\_conditioned\_from\_lower}, \formal{Conditioned.SpectralReduction.paper\_main\_subcritical\_spectral\_reduction}, \formal{Head.SubcriticalFit.paper\_main\_head\_subcritical\_fit\_exception}, \formal{Cone.HeadTailFit.paper\_main\_cone\_head\_tail\_conditioned\_fit} & same & paper's route & constants $(2\kappa+1)^{\pm1}$\\
Lem.~3.2 \texttt{\small lem:convex-separation} & Prop.~\ref{prop:kernel}, Lem.~\ref{lem:kernel_fit} & \formal{Conditioned.SupportFit.paper\_main\_box\_support\_conditioned\_fit}, \formal{Conditioned.SupportFit.paper\_main\_cone\_kernel\_conditioned\_fit}, \formal{Cone.Kernel.paper\_main\_cone\_kernel\_exists\_support}, \formal{Cone.Kernel.paper\_main\_cone\_kernel\_exists}, \formal{Cone.KernelSlice.paper\_main\_cone\_kernel\_slice\_compact} & weaker & paper's route & sufficient direction; nearest point on the slice\\
Prop.~3.3 \texttt{\small prop:heads} & Thm.~\ref{thm:heads}, Prop.~\ref{prop:width_transfer} & \formal{Head.RegularProbability.paper\_main\_head\_width\_coercive\_probability\_uniform}, \formal{Head.RegularProbability.paper\_main\_head\_width\_coercive\_probability\_fixed}, \formal{Head.WidthProbability.paper\_main\_head\_width\_regular\_probability\_fixed}, \formal{Head.WidthRegular.paper\_main\_head\_width\_regular\_event\_fixed}, \formal{Head.WidthRegular.paper\_main\_head\_width\_regular\_transfer}, \formal{Head.WidthActual.paper\_main\_head\_width\_actual\_selected} & weaker & paper's route & coercivity as hypothesis; cone $\cone_{2k_0^2}$\\
Prop.~3.4 \texttt{\small prop:lowinf} & Thm.~\ref{thm:lowinf}, Prop.~\ref{prop:all_heads} & \formal{Head.SubcriticalProbability.paper\_main\_head\_subcritical\_probability}, \formal{Head.SubcriticalFixed.paper\_main\_head\_subcritical\_fixed\_event}, \formal{Head.SubcriticalFixed.paper\_main\_head\_list\_positive}, \formal{Head.SubcriticalFamily.paper\_main\_head\_subcritical\_family\_event}, \formal{Head.SubcriticalGood.paper\_main\_head\_subcritical\_selected\_good}, \formal{Head.SubcriticalEvent.paper\_main\_head\_subcritical\_bad} & weaker & paper's route & positivity only\\
Lem.~3.5 \texttt{\small lem:regularity} & Lem.~\ref{lem:peel} & \formal{Head.TailDecompositionLog.paper\_main\_head\_tail\_regularity}, \formal{Head.TailDecomposition.head\_tail\_decomposition\_explicit}, \formal{Head.Peeling.peeling\_iteration} & same & paper's route & $C_\theta = 16/(\theta^2\ln 2)$\\
Lem.~3.8 \texttt{\small lem:head-correction-bounds} & Prop.~\ref{prop:correction}, Lem.~\ref{lem:rates} & \formal{Geometry.HeadSpectral.paper\_main\_head\_correction\_bounds\_cone}, \formal{Geometry.HeadEntries.paper\_main\_head\_correction\_bounds}, \formal{Geometry.HeadCorrection.paper\_main\_head\_condition\_transfer}, \formal{Geometry.HeadCorrection.paper\_main\_head\_support\_transfer}, \formal{Head.CorrectionScale.paper\_main\_head\_correction\_eventual}, \formal{Head.WidthActual.paper\_main\_head\_width\_actual\_selected} & weaker & paper's route & rates in $t = d^{1/32}$\\
Lem.~3.9 \texttt{\small lem:interpolation} & Thm.~\ref{thm:interp_soft} & \formal{Soft.Comparison.paper\_main\_soft\_public\_comparison}, \formal{Soft.PublicFeatures.paper\_main\_public\_feature\_soft\_comparison}, \formal{Multirow.IntegratedBound.paper\_multirow\_integrated\_bound\_results} & weaker & paper's route & soft minimax only\\
Lem.~3.10 \texttt{\small lem:chaos-ibp} & Thm.~\ref{thm:ibp}, Thm.~\ref{thm:interp_soft} & \formal{Gaussian.Vector.Stein.gaussian\_vector\_second\_stein}, \formal{Gaussian.Law.LibraryBridge.gaussian\_iterated\_coordinate\_second\_stein\_gauss}, \formal{GaussianAnalysis.gaussian.gauss\_ibp.gauss\_ibp2}, \formal{Quadratic.TimeStein.quadratic\_gaussian\_time\_stein}, \formal{Joint.SteinIdentities.joint\_gaussian\_stein\_cancellation}, \formal{Row.ConditionalStein.gaussian\_row\_full\_stein\_of\_regular}, \formal{Multirow.FullStein.paper\_multirow\_full\_stein\_results}, \formal{Chaos.SecondContractionBounds.chaos\_projective\_bound}, \formal{Multirow.IntegratedBoundSteps.multirow\_integrated\_remainder\_sum\_bound} & weaker & paper's route & soft minimax only; union of both counterpart sets\\
Lem.~3.11 \texttt{\small lem:max-radius} & proof of Thm.~\ref{thm:interp_soft} & \formal{Radius.ExpectationBound.gaussian\_expected\_maximum\_radius}, \formal{Radius.ExpectedMaximum.gaussian\_maximum\_radius\_expected\_bound} & same & paper's route & identical bound\\
Lem.~3.12 \texttt{\small lem:soft} & Def.~\ref{def:softmm}, Thm.~\ref{thm:interp_soft} & \formal{Soft.Extrema.soft\_minimax\_bounds}, \formal{Soft.MatrixProjectiveBound.soft\_matrix\_projective\_sum\_bound}, \formal{Tensor.ProjectiveSum.matrix\_nested\_projective\_sum\_low\_influence}, \formal{Soft.BlockThirdBound.lower\_block\_third\_bound} & same & paper's route & $C = 10$\\
Lem.~3.13 \texttt{\small lem:gaussian-soft-head} & Prop.~\ref{prop:gordon_margin} & \formal{Soft.FeatureMargin.paper\_main\_public\_goe\_soft\_margin}, \formal{Gordon.SoftMargin.paper\_main\_gordon\_soft\_width\_margin}, \formal{Head.SectionApproxWidth.paper\_main\_head\_section\_net\_width}, \formal{Head.SectionApproxWidth.paper\_main\_head\_section\_approx\_width\_exists} & weaker & replaced & mean bound only\\
Cor.~3.14 \texttt{\small cor:chaos-soft-head} & Cor.~\ref{cor:chaos_margin} & \formal{Head.SoftTail.paper\_main\_head\_soft\_expectation}, \formal{Soft.ConeMargin.paper\_main\_data\_soft\_cone\_margin}, \formal{Soft.ParameterChoice.paper\_main\_soft\_scaled\_margin}, \formal{Soft.ParameterChoice.paper\_main\_soft\_parameters} & same & paper's route & scores divided by $\sqd$\\
Lem.~3.15 \texttt{\small lem:clipped-soft-heads} & Lem.~\ref{lem:clip}, Props.~\ref{prop:fixed_head}, \ref{prop:all_heads} & \formal{Data.SoftConcentration.paper\_main\_data\_soft\_clipped\_lower\_margin}, \formal{Data.SoftConcentration.paper\_main\_data\_soft\_clipping\_mean}, \formal{Head.SubcriticalTail.paper\_main\_head\_subcritical\_clipped\_tail}, \formal{Head.CompactTail.paper\_main\_head\_compact\_clipped\_tail}, \formal{Head.SoftTail.paper\_main\_head\_soft\_clipped\_lower\_tail}, \formal{Head.SubcriticalFamily.paper\_main\_head\_subcritical\_family\_event}, \formal{Clipping.Limits.paper\_main\_clipped\_head\_union\_tends\_to\_zero} & weaker & paper's route & coordinate clipping, rate $e^{-cd^{3/2}}$\\
Thm.~4.1 \texttt{\small thm:unsat} & Thm.~\ref{thm:unsat} & \formal{Unsat.Limit.paper\_main\_unsat\_supercritical\_limit}, \formal{Unsat.Limit.paper\_main\_unsat\_fixed\_density\_limit}, \formal{Unsat.Limit.paper\_main\_unsat\_fixed\_density\_limit\_from\_events}, \formal{Fixed.Density.paper\_main\_unsat\_from\_fixed} & same & paper's route & measurability added\\
Lem.~4.2 \texttt{\small lem:homogeneous-witness} & Lem.~\ref{lem:homog} & \formal{Projection.PinchingLoss.paper\_main\_fit\_unit\_witness}, \formal{Geometry.Homogeneous.paper\_main\_homogeneous\_kernel}, \formal{Geometry.Homogeneous.paper\_main\_homogeneous\_witness} & stronger & paper's route & forward direction without spanning\\
Prop.~4.3 \texttt{\small prop:translated-univ} & Thm.~\ref{thm:universality} & \formal{Hybrid.FiniteMean.paper\_main\_hybrid\_finite\_mean\_comparison}, \formal{Hybrid.FiniteMean.paper\_main\_hybrid\_finite\_row\_comparison}, \formal{Loss.Replacement.paper\_main\_loss\_array\_comparison}, \formal{Loss.RowIncrement.paper\_main\_loss\_row\_increment\_comparison}, \formal{Unsat.FiniteTail.paper\_main\_unsat\_ordered\_tail} & same & replaced & finite priors, soft minima\\
Prop.~4.4 \texttt{\small prop:hybrid-gap} & Thm.~\ref{thm:hybrid} & \formal{Hybrid.Assembly.paper\_main\_hybrid\_gap\_event}, \formal{Hybrid.Assembly.paper\_main\_hybrid\_gap\_event\_joint}, \formal{Hybrid.Loss.paper\_main\_hybrid\_risk\_gap}, \formal{Hybrid.Residual.paper\_main\_hybrid\_residual\_gap} & stronger & paper's route & fewer hypotheses\\
Prop.~4.5 \texttt{\small prop:true-net-gap} & Prop.~\ref{prop:exception}, Sec.~\ref{sec:unsat_assembly} & \formal{Unsat.Orientation.paper\_main\_unsat\_orientation\_constants}, \formal{Orientation.Exception.paper\_main\_orientation\_exception}, \formal{Orientation.Union.paper\_main\_orientation\_union\_fit}, \formal{Orientation.Family.paper\_main\_orientation\_family\_uniform}, \formal{Unsat.FiniteTail.paper\_main\_unsat\_ordered\_tail} & same & replaced & net of orientations only\\
Lem.~4.6 \texttt{\small lem:transport} & Prop.~\ref{prop:witness} & \formal{Projection.Pinching.paper\_main\_pinching\_perturbation}, \formal{Projection.PinchingLoss.paper\_main\_pinching\_witness\_risk}, \formal{Projection.PinchingLoss.paper\_main\_pinching\_cutoff\_bulk}, \formal{Projection.PinchingAnnulus.paper\_main\_pinching\_annulus}, \formal{Orientation.Witness.paper\_main\_orientation\_fit\_witness} & same & replaced & pinching instead of rotation\\
Prop.~4.8 \texttt{\small prop:finite-prior-free-entropy} & Thm.~\ref{thm:universality}, Props.~\ref{prop:replacement}--\ref{prop:smooth_comparison} & \formal{Hybrid.FiniteMean.paper\_main\_hybrid\_finite\_mean\_comparison}, \formal{Loss.Replacement.paper\_main\_loss\_array\_comparison}, \formal{Loss.RowIncrement.paper\_main\_loss\_row\_increment\_comparison}, \formal{Loss.ReplicaComparison.paper\_main\_loss\_product\_comparison}, \formal{Loss.ReplicaComparison.paper\_main\_loss\_replicated\_comparison}, \formal{Loss.LogPolynomial.paper\_main\_loss\_log\_remainder}, \formal{Loss.BlockComparison.paper\_main\_loss\_conditional\_prefix\_comparison}, \formal{Projection.PublicOperator.paper\_main\_projection\_public\_comparison}, \formal{Projection.Comparison.paper\_main\_projection\_smooth\_comparison} & same & replaced & complete proof, own route\\
Lem.~4.9 \texttt{\small lem:ground-state-transfer} & Prop.~\ref{prop:prior} & \formal{Orientation.Prior.paper\_main\_orientation\_prior}, \formal{Loss.QuadraticPrior.paper\_main\_quadratic\_prior\_positive}, \formal{Orientation.BlocksLoss.paper\_main\_orientation\_candidate\_bound}, \formal{Orientation.BlocksLoss.paper\_main\_orientation\_candidate\_bound\_rotate}, \formal{Loss.JointNet.paper\_main\_loss\_joint\_net} & same & replaced & finite prior, no limit\\
Lem.~4.10 \texttt{\small lem:head-smallball} & Prop.~\ref{prop:smallball} & \formal{Head.SmallBallUniform.paper\_main\_head\_small\_ball\_uniform\_poly}, \formal{Head.SmallBallUniform.paper\_main\_head\_small\_ball\_uniform}, \formal{Head.SmallBallPZ.paper\_main\_head\_small\_ball\_probability}, \formal{Head.SmallBallFourth.paper\_main\_head\_small\_ball\_moments}, \formal{Head.SmallBallCount.paper\_main\_head\_row\_count\_lower\_exp}, \formal{Head.SmallBallNet.paper\_main\_head\_small\_ball\_param\_net}, \formal{GaussianAnalysis.probability.moments.paley\_zygmund} & same & replaced & $a_0^2 = 1/4864$; nets and binomial tails\\
Lem.~4.11 \texttt{\small lem:projected-bulk-escape} & Prop.~\ref{prop:bulk}, Lem.~\ref{lem:bulkop} & \formal{Hybrid.Bulk.paper\_main\_hybrid\_bulk\_event}, \formal{Projected.BulkConditional.paper\_main\_projected\_psd\_conditional\_tail}, \formal{Projected.BulkEscape.paper\_main\_projected\_escape\_tail}, \formal{Projected.BulkUnion.paper\_main\_projected\_count\_entropy}, \formal{Projected.BulkLaw.paper\_main\_projected\_law}, \formal{Projected.BulkComplement.paper\_main\_projected\_complement\_mx}, \formal{Hybrid.Assembly.paper\_main\_hybrid\_operator\_event} & stronger & paper's route & no small-ball condition\\
Lem.~4.12 \texttt{\small lem:conditional-true-mean-gap} & Cor.~\ref{cor:tail}, Cor.~\ref{cor:goe_mean} & \formal{Hybrid.FiniteMean.paper\_main\_hybrid\_finite\_goe\_mean}, \formal{Hybrid.FiniteMean.paper\_main\_hybrid\_finite\_mean\_comparison}, \formal{Loss.ExpectedGap.paper\_main\_loss\_expected\_gap}, \formal{Hybrid.FiniteTail.paper\_main\_hybrid\_finite\_data\_lower\_tail}, \formal{Unsat.FiniteTail.paper\_main\_unsat\_ordered\_tail} & weaker & replaced & no conditioning on the head\\
\end{longtable}}

Count by statement relation: same 18, weaker 15, stronger 6, absent 0. Count by argument relation: paper's route 26, replaced 13.

\subsection{Results the paper cites without proof}\label{sec:corr_cited}

The formal proof assumes nothing from the paper or the literature: every
cited result is proved, replaced by a weaker statement that it proves, or
avoided. One line per cited result; the details are in section (a) of the
markdown document.

{\small\setlength{\tabcolsep}{3pt}
\begin{longtable}{@{}p{46mm}p{34mm}p{70mm}@{}}
\caption{Results the paper cites without proof.}\label{tab:cited}\\
\hline
Cited result & Used by the paper in & The formal proof\\
\hline
\endfirsthead
\multicolumn{3}{@{}l}{\small\emph{Table~\ref{tab:cited}, continued.}}\\
\hline
Cited result & Used by the paper in & The formal proof\\
\hline
\endhead
\hline
\endfoot
\hline
\endlastfoot
Kogan, Nandy and Huang, extreme eigenvalues of kernel matrices, Thm~1.2(2a) & Lemma 2.6 (l.543) & Own moment method with non-sharp constants: Prop.~\ref{prop:upperedge}, Thm.~\ref{thm:loweredge} (lower edge for $4n \le d^2$ only).\\
Tropp, matrix Bernstein, Thm~6.6.1 & Lemma 2.7 (l.585) & Not used; a fourth trace moment and Markov's inequality, Prop.~\ref{prop:diag}.\\
Vershynin, norms of rectangular Gaussian matrices, Thm~7.3.1 & Lemma 2.5 (l.495), used in Prop. 2.8, Lemma 2.10, eq:Gstar-bound, Lemma 4.11 & A net bound with worse constants, Lem.~\ref{lem:opnet}; the sharp form, Thm.~\ref{thm:opnorm}, is in the library but unused.\\
Rudelson and Vershynin, Hanson--Wright inequality & Lemma 2.11 (l.660), in prop:heads & Own proof for Gaussian vectors by diagonalization and chi-square tails, Lem.~\ref{lem:hw}.\\
Gordon (1985), min--max comparison of finite Gaussian arrays & Lemma 2.3 (l.448-454) & The expectation form for finite families, Thm.~\ref{thm:gordon} and Cor.~\ref{cor:bil}; the probability form is never proved (R2).\\
Gordon (1988), escape through a mesh & Lemma 2.4 (l.478), in Lemma 4.11 & Proved from Gordon's comparison and concentration, Thm.~\ref{thm:escape}, with threshold $\sqrt{N-1}$.\\
eq:chi-mean, $\sqrt{s-1} \le \E\norm{g_s} \le \sqrt s$ & Lemmas 2.3, 2.4, 4.11 & Proved (\formal{Gaussian.Lipschitz.EuclideanNorm.paper\_main\_gaussian\_norm\_expectation}, \formal{Gaussian.Lipschitz.NormConcentration.paper\_main\_gaussian\_norm\_mean\_lower}); the Gamma-function formula is not.\\
Boucheron, Lugosi and Massart, Gaussian concentration, Thm~5.6 & Lemmas 2.3, 3.15, eq:radius (l.469, l.1551) & Herbst's argument in the library, Thm.~\ref{thm:conc}.\\
Boucheron, Lugosi and Massart, McDiarmid's inequality, Thm~6.2 & Prop. 4.5 (l.2331) & Own bounded-difference inequality on full rows, Prop.~\ref{prop:bdd}, with exponent $Nt^2/2$.\\
Amelunxen, Lotz, McCoy and Tropp, Prop.~10.2 & Lemma 2.1 (l.365) & Written out for the PSD cone with the Gaussian Poincar\'e inequality, Prop.~\ref{prop:psdwidth}.\\
Anderson, Guionnet and Zeitouni, Thm~2.1.22, and the semicircle law & Lemma 2.2 (l.390-401) & Not used; a net bound $\norm{Z}_\op \le 20\sqd$ and the second moment of the trace, Prop.~\ref{prop:conewidth}.\\
Bandeira and Maillard, Lemma~3.6 & Lemma 2.2 (the paper reproves it) & Not used.\\
Bandeira and Maillard, free-entropy universality, Thm~4.6 & Prop. 4.8 (only sketched, l.1955) & Not used; a complete proof by another route, Thm.~\ref{thm:universality} (R10).\\
Lindeberg's replacement principle & Prop. 4.8 (l.1955) & The library's Thm.~\ref{thm:lindeberg} is unused; the replacement needed is Prop.~\ref{prop:replacement}.\\
Second-chaos interpolation and Gaussian comparison with matching covariances & Lemma 3.9 (l.1254-1275), Prop. 4.8 (l.1955) & The library's Thms.~\ref{thm:interp} and \ref{thm:kahane} are unused by the tree; the theories use their own interpolation calculus, Thm.~\ref{thm:interp_soft} and Prop.~\ref{prop:smooth_comparison}.\\
Aliprantis and Border, measurable maximum theorem, Thm~18.19 & Prop. 3.3 (l.964), sec. 3.5 (l.1553) & Not used; a measurable approximate maximizer over a finite net, Lem.~\ref{lem:selector}, with the conditional step Lem.~\ref{lem:residual}.\\
Second-order Gaussian integration by parts (stated without reference) & Lemma 3.10 (l.1208) & Proved for $f$ of polynomial growth, Thm.~\ref{thm:ibp}.\\
The chi-square moment generating function & Lemma 3.11 (l.1242) & Proved from the library's normal integral (\formal{GaussianAnalysis.comparison.chi\_square.integral\_normal\_expR\_sqr}), used in Thm.~\ref{thm:chi2} and in the proof of Thm.~\ref{thm:interp_soft}.\\
Weyl's inequality & Lemma 3.8 (l.1042) & Replaced by addition in the Loewner order, Prop.~\ref{prop:correction}.\\
Packing bounds for nets, $(1 + 2/\rho)^N$ and the like & sec. 3.5, Prop. 4.5, Lemma 4.9 & Grid nets with log-cardinality $(4 + 4/\varepsilon)\ln 2$ times the dimension, Lem.~\ref{lem:section_nets} and Prop.~\ref{prop:onets}.\\
Szarek, nets of the Grassmannian & Prop. 4.5 (l.1755) & Own net of orthogonal frames in the projector distance, Prop.~\ref{prop:onets}.\\
The binomial entropy bounds & Lemma 4.11 (l.2171), sec. 3.5 (l.1562) & Replaced by $\sum_{j \le k}\binom nj \le \exp(k(1 + \ln(n/k)))$, Prop.~\ref{prop:bulk}, and by the count of head families, Prop.~\ref{prop:all_heads}.\\
The principal-angle identity & Lemma 4.6 & Not needed; pinching, Prop.~\ref{prop:witness}.\\
Gaussian hypercontractivity for degree two & Lemma 4.10 (l.2088) & Not used; $\E Y^4 \le 19\sigma^4$ is computed, Prop.~\ref{prop:smallball}.\\
Paley--Zygmund inequality & Lemma 4.10 (l.2088) & The library's Thm.~\ref{thm:pz}.\\
Sauer's lemma and the VC inequality & Lemma 4.10 (l.2093-2098) & Not used; a net of parameters, a Markov truncation and binomial tails, Prop.~\ref{prop:smallball}.\\
Gordan--Stiemke duality & Cor. 1.4, Lemma 2.12 (l.173) & Proved directly, Thm.~\ref{thm:alternative}.\\
Rotation invariance of the standard Gaussian (implicit) & eq:decompose, Lemma 4.11 & Proved in the library, Prop.~\ref{prop:mulmx}; used for the bulk array (Prop.~\ref{prop:bulk}) and the orientations (Sec.~\ref{sec:prior}).\\
Spectral theorem and spectral projectors (implicit) & eq:split-def, sgn(M\_u) in thm:sat, and throughout & The library's Thm.~\ref{thm:spectral}; $\sign(M_u)$ is avoided (Prop.~\ref{prop:support_all}).\\
Operator norm from traces of powers and from quadratic forms & l.2158-2167 and the edge arguments & Props.~\ref{prop:opnorm_trace} and \ref{prop:abab}.\\
Almost sure spanning of the sample for $n \ge d$ & Lemma 2.12, cor:alternative, thm:unsat (l.678, l.1649) & Proved where needed, Lem.~\ref{lem:spanning}; the unsatisfiable phase does not need it.\\
\end{longtable}}

\subsection{Changes of route}\label{sec:corr_routes}

Section (b) of the markdown document describes fourteen places where the
formal proof takes another route than the paper. Each is summarized here
with the place of this document that gives the formal argument.

\begin{description}
\item[R1. Widths of cones (sec. 2.1, Lemmas 2.1 and 2.2).] The projection onto the PSD cone is built as a nearest point with Moreau's decomposition (Prop.~\ref{prop:moreau}), and Lemma 2.1 uses the Poincar\'e inequality (Prop.~\ref{prop:psdwidth}). Lemma 2.2 keeps the paper's candidate but bounds its value on and off the event $\norm{Z}_\op \le 20\sqd$, so no semicircle law and no uniform integrability are used (Prop.~\ref{prop:conewidth}).
\item[R2. Gordon's inequality in expectation (secs. 2.2, 3.3, 3.5; l.434-476, l.1451-1471).] Gordon's inequality is applied in expectation to the finite nets of dual vectors and of section elements, then the log-sum-exp bound (Thm.~\ref{thm:gordon}, Cor.~\ref{cor:bil}, Prop.~\ref{prop:gordon_margin}); the net loss is also in expectation (Lem.~\ref{lem:section_nets}). Probability enters only through the concentration of the clipped data functional (Lem.~\ref{lem:clip}).
\item[R3. Feature edge (sec. 2.3, l.538-616).] Both feature edges are proved by moment methods with explicit non-sharp constants, the lower edge only up to density $1/4$ (Sec.~\ref{sec:edge}: Prop.~\ref{prop:diag}, Prop.~\ref{prop:upperedge}, Thm.~\ref{thm:loweredge}). The head correction needs only the lower edge (Prop.~\ref{prop:correction}).
\item[R4. Clipping (secs. 2.4, 3.5; l.1517-1566).] Every coordinate is clipped at $d^{1/4}$ instead of the radius at $2\sqd$. The clipped functional is $2\kappa d^{-1/4}$-Lipschitz, so one head fails with probability $e^{-cd^{3/2}}$, which still beats the union over heads (Lem.~\ref{lem:clip}, Props.~\ref{prop:fixed_head} and \ref{prop:all_heads}).
\item[R5. Head sections (sec. 3.2, Prop. 3.3, Lemma 3.8; l.960-1096).] The Borel maximizer is replaced by the first maximizer over a finite net (Lem.~\ref{lem:selector}), the residuals are bounded at the coarse threshold $d^{-1/4}$ with a Markov step at $\tau = d^{-4}$ (Lem.~\ref{lem:residual}), and the Weyl step uses only the spectral box, so the corrected matrices lie in $\cone_{2k_0^2}$ (Prop.~\ref{prop:correction}, Lem.~\ref{lem:rates}, Prop.~\ref{prop:width_transfer}, Thm.~\ref{thm:heads}).
\item[R6. Assembly of thm:sat (sec. 3.1; l.805-821, l.889-924).] The separation is done on the trace-one slice of the cone, by the paper's nearest-point argument, and the fit comes from a kernel element (Prop.~\ref{prop:kernel}, Lem.~\ref{lem:kernel_fit}). In the tiny-tail case the perturbation is $\delta M_u/\norm{M_u}_\Fro$, so no spectral sign is needed (Prop.~\ref{prop:support_all}); the constants are $(2\kappa+1)^{\pm1}$ (Thm.~\ref{thm:sat}).
\item[R7. Low-influence tail (secs. 3.4 and 3.5; l.1254-1275).] The interpolation is proved for the soft minimax only (Thm.~\ref{thm:interp_soft}). The head is fixed, the tail probability is bounded under the law of the tail rows and integrated over the head by a section bound, and the bad event is the exact min-max, so the passages from nets to sets and from soft to hard extrema happen inside it (Props.~\ref{prop:fixed_head} and \ref{prop:all_heads}, Thm.~\ref{thm:lowinf}).
\item[R8. Constant spectral cutoff (sec. 4.1, eq:epsd, eq:split-def; l.1665-1683).] A witness is split at a constant cutoff $\varepsilon_0$ instead of $d^{-1/4}$, so the head rank is bounded by a constant and the bulk satisfies only $\norm{B}_\op \le \varepsilon$ (Props.~\ref{prop:split} and \ref{prop:witness}). The universality error is made small by choosing $\varepsilon$ after $C(\beta,\eta)$, not by letting $d$ grow (Sec.~\ref{sec:unsat_assembly}).
\item[R9. Nets, pinching and finite priors (secs. 4.1 and 4.4; l.1755-1851).] A net of orientations in the projector distance replaces the net of tuples $(r, U, H, b)$ (Prop.~\ref{prop:onets}), a witness is pinched instead of rotated (Prop.~\ref{prop:witness}), and for each orientation the sample is rotated and one finite prior of candidates enters a soft minimum (Prop.~\ref{prop:prior}). Nothing is conditioned on the head: the hybrid mean is over heads and bulks together (Cor.~\ref{cor:goe_mean}, Prop.~\ref{prop:exception}).
\item[R10. Bulk universality (sec. 4.2; l.1917-2025).] Rows are replaced one at a time (Prop.~\ref{prop:replacement}); the soft-min increment of a row is expanded into replica averages of products of loss factors (Prop.~\ref{prop:onerow}); each product is compared with its Gaussian counterpart by a fixed-dimensional smooth comparison (Prop.~\ref{prop:smooth_comparison}). The error $2\eta + C(\beta,\eta)\varepsilon\norm{B}_\Fro$ is uniform in $N$, the offsets and the prior (Thm.~\ref{thm:universality}).
\item[R11. Hybrid gap (sec. 4.3; l.2081-2108).] The head small ball is proved with a net of the parameters, a Markov truncation and binomial tails, and a computed fourth moment (Prop.~\ref{prop:smallball}). The projected bulk bound uses a complement basis and Bessel's inequality instead of the projection $P_J$, so no small-ball condition is needed (Prop.~\ref{prop:bulk}).
\item[R12. Concentration in the unsatisfiable phase (sec. 4.4; l.2329-2349).] McDiarmid's inequality is replaced by a bounded-difference inequality for Gaussian rows applied to full rows, with oscillation $1/N$ per row (Prop.~\ref{prop:bdd}, Cor.~\ref{cor:tail}).
\item[R13. Measurability (everywhere).] Every event is proved measurable: the three events of the statement by countable rational approximation (Prop.~\ref{prop:measurable}), the exceptional events by continuity of the functionals of the sample.
\item[R14. Fixed densities (everywhere).] The reductions to fixed densities are statements on probabilities: deleting points cannot increase the fit probabilities (Lem.~\ref{lem:drop}, Prop.~\ref{prop:fixed}).
\end{description}

\end{document}
